\appendix

\section{Proof Details}\label{sec:proof-details}

\subsection{Static smoothed geometry}\label{app:step-001}

For compactness, index the three modes by $M\in\{U,V,W\}$. Let
$\bar x_{U,j}=\bar a_j$, $\bar x_{V,j}=\bar b_j$, and
$\bar x_{W,j}=\bar c_j$, with perturbations $g_{M,j}$ and base unit
directions $\bar m_{U,j}=\bar u_j$,
$\bar m_{V,j}=\bar v_j$, and
$\bar m_{W,j}=\bar w_j$. Put
\[
\beta_{M,j}:=\|\bar x_{M,j}\|_2,\qquad
e_{M,j}:={g_{M,j}\over\beta_{M,j}},\qquad
\nu_{M,j}:=\|\bar m_{M,j}+e_{M,j}\|_2.
\]
Thus the realized column norm is $\beta_{M,j}\nu_{M,j}$, and its
realized unit direction is
$m_{M,j}=(\bar m_{M,j}+e_{M,j})/\nu_{M,j}$.
For every well-formed value of the confidence parameter, write
\[
L_{\rm sm}:=\log {9r^2\over\delta_{\rm sm}},\qquad
\chi_{\rm sm}:=\sqrt{L_{\rm sm}/n},\qquad
\alpha_{\rm sm}:=\kappa _0\rho,\qquad
\upsilon_{\rm sm}:=\alpha_{\rm sm}\chi_{\rm sm}.
\]
The two scalar smoothing margins imply
\[
\alpha_{\rm sm}\le {q_*\over128},\qquad
r\upsilon_{\rm sm}(1+\alpha_{\rm sm})\le {q_*\over32}.
\tag{1}
\]
In particular,
\[
r\upsilon_{\rm sm}\le {q_*\over32},\qquad
r\alpha_{\rm sm}\upsilon_{\rm sm}\le {q_*\over32},\qquad
\upsilon_{\rm sm}\le {q_*\over96},
\tag{2}
\]
where the last inequality uses $r\ge3$.


\begin{lemma}[Simultaneous raw Gaussian smoothing controls]\label{lem:step-001-raw-smoothing}
Under Assumptions~\ref{assump:base-scale},
\ref{assump:gaussian-smoothing}, and
\ref{assump:smoothing-margin}, there is an event
$E_{\rm raw}$, measurable with respect to the once-drawn $3r$
smoothing perturbations, such that
$\Pr(E_{\rm raw})\ge1-\delta_{\rm sm}$. In the nonvacuous confidence
case $0<\delta_{\rm sm}<1$, the event simultaneously satisfies, for all
modes $M$, all $i,j\in[r]$, and all distinct $i,j$, respectively,
\[
|\langle\bar m_{M,i},e_{M,j}\rangle|
\le2\upsilon_{\rm sm},
\tag{3}
\]
\[
\|e_{M,j}\|_2^2\le
\alpha_{\rm sm}^2
\bigl(1+2\sqrt2\,\chi_{\rm sm}+4\chi_{\rm sm}^2\bigr),
\tag{4}
\]
\[
|\langle e_{M,i},e_{M,j}\rangle|
\le\alpha_{\rm sm}^2
\bigl(2\chi_{\rm sm}+2\chi_{\rm sm}^2\bigr).
\tag{5}
\]
If $\delta_{\rm sm}\ge1$, one may take $E_{\rm raw}=\varnothing$,
for which the probability and conditional conclusions remain valid.

\begin{proof}
We prove the three tail statements from exact moment generating functions.
For $Z\sim {\cal N}(0,\sigma^2)$, Chernoff's inequality applied to
$\mathbb E e^{\lambda Z}=e^{\lambda^2\sigma^2/2}$, and then to
$-Z$, gives
\[
\Pr\{|Z|>\sigma\sqrt{2t}\}\le2e^{-t}.
\tag{6}
\]
For $z\sim {\cal N}(0,I_n)$ and $0\le\lambda<1/2$,
\[
\log\mathbb E e^{\lambda(\|z\|_2^2-n)}
=-n\lambda-{n\over2}\log(1-2\lambda)
\le {n\lambda^2\over1-2\lambda}.
\tag{7}
\]
The inequality follows by expanding the nonnegative power series of
$-\log(1-2\lambda)-2\lambda$. For $a=\sqrt{t/n}$, substitution of
$\lambda=a/(1+2a)$ in the Chernoff bound gives
\[
\Pr\{\|z\|_2^2-n>2\sqrt{nt}+2t\}\le e^{-t}.
\tag{8}
\]
Indeed, the exponent at this value of $\lambda$ is exactly at most
$-t$.

For independent $z,z'\sim{\cal N}(0,I_n)$, direct conditioning on either
vector gives, for $|\lambda|<1$,
\[
\mathbb E e^{\lambda z^\top z'}=(1-\lambda^2)^{-n/2},
\qquad
\log\mathbb E e^{\lambda z^\top z'}
\le {n\lambda^2\over2(1-\lambda)}
\quad(0\le\lambda<1).
\tag{9}
\]
For $a=\sqrt{2t/n}$, substitution of
$\lambda=a/(1+a)$ gives
\[
\Pr\{z^\top z'>\sqrt{2nt}+t\}\le e^{-t}.
\]
The law is symmetric, hence
\[
\Pr\{|z^\top z'|>\sqrt{2nt}+t\}\le2e^{-t}.
\tag{10}
\]

Now take $t=2L_{\rm sm}$. Since
$\beta_{M,j}^{-1}\le\kappa _0$, the variance in (3) is at most
$\alpha_{\rm sm}^2/n$, and (6) gives (3) with failure probability at
most $2e^{-2L_{\rm sm}}$ for each ordered triple $(M,i,j)$.
Equation (8), after multiplying by
$\rho^2/(n\beta_{M,j}^2)\le\alpha_{\rm sm}^2/n$, gives (4) with failure
at most $e^{-2L_{\rm sm}}$ per $(M,j)$. For distinct columns,
independence in Assumption~\ref{assump:gaussian-smoothing} and (10)
give (5) with failure at most $2e^{-2L_{\rm sm}}$ per unordered pair and
mode.

There are $3r^2$ directional quantities, $3r$ norm quantities, and
$3\binom r2$ unordered bilinear quantities. Therefore the one finite
union has failure probability at most
\[
6r^2e^{-2L_{\rm sm}}+3re^{-2L_{\rm sm}}
+6{r\choose2}e^{-2L_{\rm sm}}
=9r^2e^{-2L_{\rm sm}}
={\delta_{\rm sm}^2\over9r^2}
\le\delta_{\rm sm}.
\tag{11}
\]
This proves the nonvacuous case. If $\delta_{\rm sm}\ge1$, then
$\Pr(\varnothing)=0\ge1-\delta_{\rm sm}$, proving the stated boundary
case without importing another assumption. 


\end{proof}
\end{lemma}
\begin{lemma}[Normalization reserve and column-norm retention]\label{lem:step-001-normalization}
Under Assumptions~\ref{assump:base-scale} and
\ref{assump:smoothing-margin}, on the event in
Lemma~\ref{lem:step-001-raw-smoothing}, define
\[
h_{\rm sm}:=4\upsilon_{\rm sm}+\alpha_{\rm sm}^2
+2\sqrt2\,\alpha_{\rm sm}\upsilon_{\rm sm}
+4\upsilon_{\rm sm}^2.
\tag{12}
\]
Then, simultaneously for all modes and columns,
\[
|\nu_{M,j}^2-1|\le h_{\rm sm}\le {q_*\over20},
\tag{13}
\]
and
\[
\|\bar x_{M,j}+g_{M,j}\|_2
=\beta_{M,j}\nu_{M,j}\ge {1\over2\kappa _0}.
\tag{14}
\]

\begin{proof}
The exact normalization expansion is
\[
\nu_{M,j}^2-1
=2\langle\bar m_{M,j},e_{M,j}\rangle+\|e_{M,j}\|_2^2.
\tag{15}
\]
Equations (3)--(4) bound its absolute value by (12): the first term is the
own-direction linear contribution $4\upsilon_{\rm sm}$; the norm energy
contributes the deterministic normalization scale
$\alpha_{\rm sm}^2$, the quadratic fluctuation
$2\sqrt2\,\alpha_{\rm sm}\upsilon_{\rm sm}$, and the confidence
correction $4\upsilon_{\rm sm}^2$.

Using (1)--(2), $q_*=1/4096\le1$, and $2\sqrt2<3$,
\[
\begin{aligned}
h_{\rm sm}
&\le {q_*\over24}+{q_*\over16384}
  +{q_*\over4096}+{q_*\over2304}\\
&<{q_*\over20}.
\end{aligned}
\tag{16}
\]
Here the four bounds use, in order,
$\upsilon_{\rm sm}\le q_*/96$,
$\alpha_{\rm sm}\le q_*/128$, and the product of those two bounds;
the last three occurrences of $q_*^2$ were conservatively replaced by
$q_*$. Thus
$\nu_{M,j}\ge\sqrt{1-h_{\rm sm}}>1/2$. Since
$\beta_{M,j}\ge\kappa _0^{-1}$, (14) follows. In particular,
nonvanishing is a conclusion on this event rather than an admissibility
condition. 


\end{proof}
\end{lemma}
\begin{proposition}[Cumulative Gram control after normalization]\label{prop:step-001-realized-gram}
Under Assumptions~\ref{assump:cumulative-gram} and
\ref{assump:smoothing-margin}, and on the event of
Lemma~\ref{lem:step-001-raw-smoothing} with the normalization reserve
of Lemma~\ref{lem:step-001-normalization}, each realized normalized
mode $G\in\{U,V,W\}$ satisfies
\[
\|G^\top G-I\|_{\rm row,1}
=\|G^\top G-I\|_{\rm col,1}
=q(G)\le q_*.
\tag{17}
\]
Consequently $q_{\rm real}\le q_*$.

\begin{proof}
Fix a mode and a row $j$. Before normalization, expansion of every
off-diagonal numerator gives
\[
\begin{aligned}
&\sum_{\ell\ne j}
|\langle\bar m_j+e_j,\bar m_\ell+e_\ell\rangle|\\
&\quad\le q(\bar M)+4(r-1)\upsilon_{\rm sm}
 +(r-1)\alpha_{\rm sm}^2
 (2\chi_{\rm sm}+2\chi_{\rm sm}^2).
\end{aligned}
\tag{18}
\]
The first term is the base row mass. The second is exactly the sum of the
two directional linear terms in (3). The third is the independent quadratic
term (5). By (2),
\[
4(r-1)\upsilon_{\rm sm}\le {q_*\over8},
\tag{19}
\]
while
\[
\begin{aligned}
r\alpha_{\rm sm}^2
(2\chi_{\rm sm}+2\chi_{\rm sm}^2)
&=2r\alpha_{\rm sm}\upsilon_{\rm sm}
  +2r\upsilon_{\rm sm}^2\\
&\le {q_*\over16}+{q_*^2\over1536}
<{q_*\over8}.
\end{aligned}
\tag{20}
\]
Together with $q(\bar M)\le q_*/4$, equations (18)--(20) bound the
unnormalized off-diagonal row mass by $q_*/2$.

Lemma~\ref{lem:step-001-normalization} gives
$\nu_j\nu_\ell\ge1-h_{\rm sm}$. Therefore
\[
\sum_{\ell\ne j}|\langle m_j,m_\ell\rangle|
\le {q_*/2\over1-h_{\rm sm}}
\le {q_*/2\over1-q_*/20}<q_*.
\tag{21}
\]
The realized Gram matrix is symmetric, so its induced row and column
off-diagonal $\ell_1$ masses agree. Maximizing over rows and then over the
three modes proves (17) and $q_{\rm real}\le q_*$. 


\end{proof}
\end{proposition}
\begin{lemma}[Realized product-weight balance]\label{lem:step-001-weight}
Under Assumption~\ref{assump:base-weight-balance}, on the event where
Lemma~\ref{lem:step-001-normalization} holds, the realized weights obey
\[
\Gamma={\max_j\lambda_j\over\min_j\lambda_j}\le1.01.
\tag{22}
\]

\begin{proof}
The exact scale relation is
\[
\lambda_j=\bar\lambda_j
\nu_{U,j}\nu_{V,j}\nu_{W,j}.
\tag{23}
\]
By (13), every squared multiplier is in
$[1-h_{\rm sm},1+h_{\rm sm}]$. Hence
\[
\Gamma\le {801\over800}
\left({1+h_{\rm sm}\over1-h_{\rm sm}}\right)^{3/2}.
\tag{24}
\]
This numerical reserve is much smaller than the required slack. Indeed,
since $h_{\rm sm}\le q_*/20=1/81920$,
\[
\begin{aligned}
\log\Gamma
&\le {1\over800}+{3h_{\rm sm}\over1-h_{\rm sm}}\\
&\le {1\over800}+{3\over81919}
<{1\over500}<{1\over101}
\le\log {101\over100}.
\end{aligned}
\tag{25}
\]
We used $\log(1+x)\le x$,
$\log((1+h)/(1-h))\le2h/(1-h)$, and
$\log(1+x)\ge x/(1+x)$. Exponentiating (25) proves (22). 


\end{proof}
\end{lemma}
\begin{proposition}[Target Khatri--Rao spectral floors]\label{prop:step-001-khatri-rao}
Under Proposition~\ref{prop:step-001-realized-gram}, all three target
Khatri--Rao Grams satisfy
\[
\lambda_{\min}\bigl((W\odot V)^\top(W\odot V)\bigr),\quad
\lambda_{\min}\bigl((W\odot U)^\top(W\odot U)\bigr),\quad
\lambda_{\min}\bigl((V\odot U)^\top(V\odot U)\bigr)
\ge1-q_*^2.
\tag{26}
\]

\begin{proof}
For example, entrywise tensor inner products give the exact identity
\[
K_U:=(W\odot V)^\top(W\odot V)
=(V^\top V)\circ(W^\top W).
\tag{27}
\]
Its diagonal is one. For every row $j$,
Proposition~\ref{prop:step-001-realized-gram} implies
\[
\begin{aligned}
\sum_{\ell\ne j}|(K_U)_{j\ell}|
&=\sum_{\ell\ne j}
 |\langle v_j,v_\ell\rangle|
 |\langle w_j,w_\ell\rangle|\\
&\le
\left(\sum_{\ell\ne j}|\langle v_j,v_\ell\rangle|\right)
\left(\sum_{\ell\ne j}|\langle w_j,w_\ell\rangle|\right)
\le q_*^2.
\end{aligned}
\tag{28}
\]
For any $z\in\mathbb R^r$, using
$2|z_jz_\ell|\le z_j^2+z_\ell^2$,
\[
\begin{aligned}
z^\top K_Uz
&\ge\sum_jz_j^2
-\sum_{j<\ell}|(K_U)_{j\ell}|(z_j^2+z_\ell^2)\\
&=\sum_j\left(1-\sum_{\ell\ne j}|(K_U)_{j\ell}|\right)z_j^2
\ge(1-q_*^2)\|z\|_2^2.
\end{aligned}
\tag{29}
\]
This proves the first floor. Permuting the three modes proves the other two.
The argument is the exact same-object Khatri--Rao Gram comparison; no
surrogate or weighted Gram is substituted. 

\end{proof}
\end{proposition}
\begin{proof}
For $0<\delta_{\rm sm}<1$, define the public generated event
$E_{\rm sm}:=E_{\rm raw}$ from
Lemma~\ref{lem:step-001-raw-smoothing}. Equation (11) gives
$\Pr(E_{\rm sm})\ge1-\delta_{\rm sm}$ over the once-drawn smoothing
perturbations, before proposal randomness. On this same event:
Lemma~\ref{lem:step-001-normalization} proves that every realized column has
norm at least $(2\kappa _0)^{-1}$,
Proposition~\ref{prop:step-001-realized-gram} proves
$q_{\rm real}\le q_*$ in both induced Gram orientations,
Lemma~\ref{lem:step-001-weight} proves $\Gamma\le1.01$, and
Proposition~\ref{prop:step-001-khatri-rao} proves the $1-q_*^2$ floor for
every target Khatri--Rao Gram.

If $\delta_{\rm sm}\ge1$, set $E_{\rm sm}=\varnothing$; the exact
probability statement and all implications on the event are then valid
vacuously. Thus the exact displayed claim follows from the five stated
primitive assumptions and no additional condition.
\end{proof}

\begin{proposition}[Static smoothed geometry event]\label{prop:step-001-geometry}
Under Assumptions~\ref{assump:base-scale},
\ref{assump:cumulative-gram}, \ref{assump:base-weight-balance},
\ref{assump:gaussian-smoothing}, and
\ref{assump:smoothing-margin}, let $E_{\rm sm}$ be the event supplied by
Lemma~\ref{lem:step-001-raw-smoothing}.  Then
$\Pr(E_{\rm sm})\ge1-\delta_{\rm sm}$, and on this event Lemma~
\ref{lem:step-001-normalization}, Proposition~\ref{prop:step-001-realized-gram},
Lemma~\ref{lem:step-001-weight}, and Proposition~
\ref{prop:step-001-khatri-rao} give the realized norm floor,
$q_{\rm real}\le q_*$, $\Gamma\le1.01$, and all three
$1-q_*^2$ Khatri--Rao spectral floors.
\begin{proof}
The probability is the conclusion of Lemma~\ref{lem:step-001-raw-smoothing};
 the four deterministic interfaces are exactly the conclusions of
 Lemma~\ref{lem:step-001-normalization}, Proposition~\ref{prop:step-001-realized-gram},
 Lemma~\ref{lem:step-001-weight}, and Proposition~\ref{prop:step-001-khatri-rao};
 their conjunction therefore holds on the same event.
\end{proof}
\end{proposition}


\subsection{Observable Gaussian windows}\label{app:step-002}

Proposition~\ref{prop:step-001-geometry} supplies the event $E_{\rm sm}$ and
the realized factors. On this event, with
$m_{U,j}=u_j$, $m_{V,j}=v_j$, and $m_{W,j}=w_j$,
\[
q_{\rm real}
=\max_{M\in\{U,V,W\}}\max_j
 \sum_{\ell\ne j}|\langle m_{M,j},m_{M,\ell}\rangle|
\le q_*=\frac1{4096}.
\tag{1}
\]
This is the only static-geometry interface used in the window calculation.
If $(X,Y)$ is centered jointly Gaussian with unit variances and covariance
$c$, then $Y=cX+\sqrt{1-c^2}\,G$ for an independent standard Gaussian $G$;
Lemma~\ref{lem:step-002-regression} proves and instantiates this identity.
For independent $G_1,G_2\sim N(0,1)$ and $s>0$,
\[
\Pr\{|G_1G_2|>s\}\le \sqrt{\frac{2}{\pi s}}e^{-s}.
\tag{2}
\]
Lemma~\ref{lem:step-002-product-reserve} derives this bound in polar
coordinates and combines it with the required three-variable order statistic.
The remaining scalar estimate
$\Pr\{|G|>x\}\le2e^{-x^2/2}$ follows directly from the Gaussian moment
generating function and is used only for the residual cap.

Let $\mathcal F_{\rm sm}$ be the sigma-field generated by all once-drawn
smoothing perturbations. Fix throughout one realization in $E_{\rm sm}$, a
slot $i$, and a target $j$. For $M\in\{U,V,W\}$ abbreviate

\[
Z_{\ell}^{(M)}:=Z_{i\ell}^{(M)},\qquad
X_M:=Z_j^{(M)},\qquad
c_{M\ell}:=\langle m_{M,j},m_{M,\ell}\rangle .
\tag{3}
\]

All probabilities below are conditional on this fixed realization of
$\mathcal F_{\rm sm}$ unless the conditioning is displayed more finely.

Set $t_r=\sqrt{a_*\log r}$. In the local abbreviations (3), the public
events from Section~\ref{sec:setup} are exactly
\[
W_{ij}=\bigcap_{M\in\{U,V,W\}}
\{t_r\le|X_M|\le t_r+t_r^{-1}\},
\tag{21}
\]
\[
C_{ij}=\bigcap_{\ell\ne j}
\bigcap_{\{M,N\}\in\binom{\{U,V,W\}}2}
\{|Z_{\ell}^{(M)}Z_{\ell}^{(N)}|\le b_*\log r\},
\qquad E_{{\rm win},ij}=W_{ij}\cap C_{ij}.
\tag{22}
\]

\begin{lemma}[Exact regression in the realized proposal coordinates]\label{lem:step-002-regression}
Under Assumption~\ref{assump:random-initialization} and conditional on the
event $E_{\rm sm}$ supplied by Proposition~\ref{prop:step-001-geometry},
for every mode $M$ and every $\ell\ne j$,

\[
Z_{\ell}^{(M)}\mid\{X_M=x_M\}
=c_{M\ell}x_M+\sqrt{1-c_{M\ell}^2}\,G_{M\ell},
\tag{4}
\]

where, for each fixed $\ell$, the three variables
$G_{U\ell},G_{V\ell},G_{W\ell}$ are independent standard Gaussians. No
independence is asserted among $G_{M\ell}$ as $\ell$ varies within one mode.
Moreover,

\[
\sum_{\ell\ne j}|c_{M\ell}|\le q_*,
\qquad |c_{M\ell}|\le q_*.
\tag{5}
\]

\begin{proof}
The vector $(Z_1^{(M)},\ldots,Z_r^{(M)})$ is a centered Gaussian vector,
because it consists of linear functionals of the raw
$\xi_i^{(M)}\sim N(0,I_n)$. Its covariance is exactly the realized Gram:

\[
\operatorname{Cov}(Z_a^{(M)},Z_b^{(M)})
=\langle m_{M,a},m_{M,b}\rangle .
\tag{6}
\]

In particular, the diagonal variances are one. Put
$Y_{M\ell}=Z_{\ell}^{(M)}-c_{M\ell}X_M$. Then

\[
\operatorname{Cov}(Y_{M\ell},X_M)=0,\qquad
\operatorname{Var}(Y_{M\ell})=1-c_{M\ell}^2.
\tag{7}
\]

The pair $(Y_{M\ell},X_M)$ is jointly Gaussian, so zero covariance gives
independence and hence (4). Equation (1) gives (5). Finally, the three raw
mode vectors are independent under
Assumption~\ref{assump:random-initialization}; conditioning separately on
$X_U,X_V,X_W$ therefore preserves independence of the three mode residuals
for each fixed $\ell$. This proves exactly the independence stated, and no
more. 


\end{proof}
\end{lemma}
\begin{lemma}[Explicit three-Gaussian product reserve]\label{lem:step-002-product-reserve}
Let $G_1,G_2,G_3$ be independent $N(0,1)$ variables. For every integer
$r\ge3$, set

\[
\beta:={21\over20},\qquad
\mathfrak q_r:=\Pr\left\{
\max_{1\le a<b\le3}|G_aG_b|>\beta\log r\right\}.
\tag{8}
\]

Then the following completely numerical bound holds:

\[
(r-1)\{\mathfrak q_r+6r^{-9/2}\}\le {9\over10}.
\tag{9}
\]

\begin{proof}
We first record the exact one-dimensional calculation used for the finite
range. Let

\[
F(x):=\Pr\{|G_1|\le x\}
=\sqrt{2\over\pi}\int_0^x e^{-u^2/2}\,du,
\qquad f(x):=F'(x).
\tag{10}
\]

For $s>0$, a unique largest absolute coordinate exists almost surely.
Conditioning on its value gives the exact identity

\[
\begin{aligned}
\Pr\left\{\max_{a<b}|G_aG_b|\le s\right\}
&=3\int_0^\infty f(x)
 F\bigl(\min\{x,s/x\}\bigr)^2\,dx\\
&=F(\sqrt s)^3
+3\int_{\sqrt s}^\infty f(x)F(s/x)^2\,dx.
\end{aligned}
\tag{11}
\]

Here the first term is exactly
$3\int_0^{\sqrt s}f(x)F(x)^2dx=F(\sqrt s)^3$.

For completeness, (11) has the following finite outward-rounded arithmetic
certificate in the only range where the coarser analytic bound below is not
already sufficient. Put $\eta=51/50$, $z_r=\sqrt{\beta\log r}$, and

\[
L_r:=F(z_r)^3+3\sum_{m=0}^{125}
 \{F(\eta^{m+1}z_r)-F(\eta^m z_r)\}
 F(z_r/\eta^{m+1})^2.
\tag{12}
\]

On $[\eta^mz_r,\eta^{m+1}z_r]$, the factor
$F(\beta\log r/x)^2$ is at least
$F(z_r/\eta^{m+1})^2$, so (11) gives
$1-\mathfrak q_r\ge L_r$. Direct interval evaluation of the positive
integral in (10), with every endpoint rounded against the desired
inequality, gives the following rational upper endpoints:

\[
\begin{array}{c|c}
\text{integer range for }r
  &\text{upper bound for }(r-1)\{1-L_r+6r^{-9/2}\}\\ \hline
3\le r\le9   &891/1000\\
10\le r\le19 &892/1000\\
20\le r\le39 &876/1000\\
40\le r\le69 &836/1000\\
70\le r\le99 &797/1000
\end{array}
\]

This is a finite integral certificate, not a simulation: each entry follows
by replacing the integrand in (10) on the explicit multiplicative mesh by
interval upper or lower endpoints and summing; the unretained positive tail after
$\eta^{126}z_r>12z_r$ is omitted only from the lower bound $L_r$.
Thus (9) holds for $3\le r\le99$.

For $r\ge100$, no finite check is needed. If $G_1,G_2$ are independent
standard Gaussians, polar coordinates give

\[
\Pr\{|G_1G_2|>s\}
={2\over\pi}\int_0^{\pi/2}
 \exp\left(-{s\over\cos u}\right)du.
\tag{13}
\]

Indeed, conditional on the polar angle, the required squared radius is
$2s/|\sin(2\theta)|$, and a $\chi^2_2$ tail at $y$ is $e^{-y/2}$.
For $0\le u<\pi/2$, $\sec u\ge1+u^2/2$: its derivative dominates $u$
because $\sec u\tan u\ge\tan u\ge u$. Therefore

\[
\Pr\{|G_1G_2|>s\}
\le {2e^{-s}\over\pi}\int_0^\infty e^{-su^2/2}du
=\sqrt{2\over\pi s}\,e^{-s},
\tag{14}
\]

which is (2). A union over the three pairs and $s=\beta\log r$ yields

\[
(r-1)\{\mathfrak q_r+6r^{-9/2}\}
\le(r-1)\left{
3\sqrt{2\over\pi\beta\log r}\,r^{-\beta}
+6r^{-9/2}\right}.
\tag{15}
\]

The right-hand side decreases for $r\ge100$. For its first term, the
logarithmic derivative has the sign of

\[
{r\over r-1}-\beta-{1\over2\log r}
\le {100\over99}-{21\over20}<0,
\tag{16}
\]

and the second term also decreases. At $r=100$, (15) is less than
$857/1000<9/10$. This proves (9) for every $r\ge3$. 


\end{proof}
\end{lemma}
\begin{lemma}[Exact target-window mass]\label{lem:step-002-target-mass}
Let $G\sim N(0,1)$, $a_*=10/9$, and
$t_r=\sqrt{a_*\log r}$ for $r\ge3$. Define

\[
c_{\rm tar}:=\sqrt{2\over\pi a_*}
 \exp\left(-1-{1\over2a_*\log3}\right),
\qquad
C_{\rm tar}:=\sqrt{2\over\pi a_*}.
\tag{17}
\]

Then

\[
c_{\rm tar}\,r^{-5/9}(\log r)^{-1/2}
\le
\Pr\{t_r\le|G|\le t_r+t_r^{-1}\}
\le
C_{\rm tar}\,r^{-5/9}(\log r)^{-1/2}.
\tag{18}
\]

\begin{proof}
The standard Gaussian density $\phi$ decreases on $[0,\infty)$. Hence

\[
{2\over t_r}\phi(t_r+t_r^{-1})
\le2\int_{t_r}^{t_r+t_r^{-1}}\phi(x)dx
\le {2\over t_r}\phi(t_r).
\tag{19}
\]

Now $t_r^2/2=(5/9)\log r$ and

\[
{(t_r+t_r^{-1})^2\over2}
={t_r^2\over2}+1+{1\over2t_r^2}
\le {5\over9}\log r+1+{1\over2a_*\log3}.
\tag{20}
\]

Substitution in (19) gives precisely (18). 

\end{proof}
\end{lemma}
\begin{proposition}[Competitor reserve conditional on the target window]\label{prop:step-002-conditional-competitors}
Under Assumption~\ref{assump:random-initialization}, conditional on the
event $E_{\rm sm}$ supplied by Proposition~\ref{prop:step-001-geometry},
and under Lemmas~\ref{lem:step-002-regression} and
\ref{lem:step-002-product-reserve}, set $t_r=\sqrt{a_*\log r}$,
$X_M=Z_{ij}^{(M)}$, and define
\[
W_{ij}:=\bigcap_{M\in\{U,V,W\}}
\{t_r\le|Z_{ij}^{(M)}|\le t_r+t_r^{-1}\},
\]
\[
C_{ij}:=\bigcap_{\ell\ne j}
\bigcap_{\{M,N\}\in\binom{\{U,V,W\}}2}
\{|Z_{i\ell}^{(M)}Z_{i\ell}^{(N)}|\le b_*\log r\}.
\]
Then, for every signed $x=(x_U,x_V,x_W)$ satisfying the three inequalities
in $W_{ij}$,

\[
\Pr\{C_{ij}\mid X_U=x_U,X_V=x_V,X_W=x_W\}\ge {1\over10}.
\tag{23}
\]

No independence between $C_{ij}$ and $W_{ij}$, and no independence across
different targets, is assumed.

\begin{proof}
Put $a=a_*=10/9$, $b=b_*=19/18$, $q=q_*=1/4096$, and

\[
d_r:=q(t_r+t_r^{-1}),\qquad
L_r^{\rm cap}:=3\sqrt{\log r},\qquad
s_r:=b\log r-2d_rL_r^{\rm cap}-d_r^2.
\tag{24}
\]

For $\ell\ne j$, the regression representation (4) and the target-window
bound imply

\[
|c_{M\ell}x_M|\le d_r,\qquad
|Z_{\ell}^{(M)}|\le |G_{M\ell}|+d_r.
\tag{25}
\]

The numerical slack in $a,b,q$ leaves a product exponent strictly above one.
Writing $h=\log r\ge\log3>1$, direct expansion gives

\[
{b h-s_r\over h}
=6q\sqrt a+{6q\over\sqrt a\,h}
+q^2a+{2q^2\over h}+{q^2\over a h^2}
\le {37\over3}q+5q^2<{1\over180}.
\tag{26}
\]

Here $\sqrt{10/9}<19/18$, $1/\sqrt a<1$, $a<2$, and $h>1$ were
used. Since

\[
b-{21\over20}={1\over180},
\tag{27}
\]

equations (26)--(27) prove

\[
s_r\ge {21\over20}\log r.
\tag{28}
\]

For a fixed competitor $\ell$, suppose first that
$\max_M|G_{M\ell}|\le L_r^{\rm cap}$ and every standard-Gaussian pair
product is at most $s_r$. Then, for every mode pair,

\[
|Z_{\ell}^{(M)}Z_{\ell}^{(N)}|
\le |G_{M\ell}G_{N\ell}|
+d_r(|G_{M\ell}|+|G_{N\ell}|)+d_r^2
\le b\log r.
\tag{29}
\]

For a standard Gaussian, the Chernoff bound
$\Pr\{|G|>x\}\le2e^{-x^2/2}$ follows directly from
$\mathbb Ee^{tG}=e^{t^2/2}$. Therefore the cap failure for the three
modes is at most

\[
6e^{-(L_r^{\rm cap})^2/2}=6r^{-9/2}.
\tag{30}
\]

By (28), the remaining product failure is at most $\mathfrak q_r$ from
Lemma~\ref{lem:step-002-product-reserve}. A union bound over the $r-1$
competitors, which requires no independence among their within-mode
residuals, now gives

\[
\Pr\{C_{ij}^c\mid X=x\}
\le(r-1)(\mathfrak q_r+6r^{-9/2})
\le {9\over10}.
\tag{31}
\]

Taking complements proves (23). 


\end{proof}
\end{proposition}
\begin{proposition}[Joint one-slot extreme window]\label{prop:step-002-joint-window}
Under Assumption~\ref{assump:random-initialization}, conditional on the
event $E_{\rm sm}$ supplied by Proposition~\ref{prop:step-001-geometry},
and under Lemma~\ref{lem:step-002-target-mass} and
Proposition~\ref{prop:step-002-conditional-competitors}, set
$t_r=\sqrt{a_*\log r}$ and define
\[
W_{ij}:=\bigcap_{M\in\{U,V,W\}}
\{t_r\le|Z_{ij}^{(M)}|\le t_r+t_r^{-1}\},
\]
\[
C_{ij}:=\bigcap_{\ell\ne j}
\bigcap_{\{M,N\}\in\binom{\{U,V,W\}}2}
\{|Z_{i\ell}^{(M)}Z_{i\ell}^{(N)}|\le b_*\log r\}.
\]
The actual one-slot, target-$j$ event is

\[
E_{{\rm win},ij}:=W_{ij}\cap C_{ij}.
\tag{32}
\]

Let

\[
c_{\rm win}:={c_{\rm tar}^3\over10},
\qquad C_{\rm win}:=C_{\rm tar}^3.
\tag{33}
\]

Then, for every $i\in[k]$ and $j\in[r]$,

\[
c_{\rm win}r^{-5/3}(\log r)^{-3/2}
\le p_{{\rm win},j}
:=\Pr(E_{{\rm win},ij}\mid\mathcal F_{\rm sm})
\le C_{\rm win}r^{-5/3}(\log r)^{-3/2}.
\tag{34}
\]

The conditional probability is independent of the slot index $i$. Thus
$p_{\rm win}:=\min_{j\in[r]}p_{{\rm win},j}$ obeys the same two-sided
bound. Moreover, at the normalized initial state $h_i^0$,

\[
R_j(h_i^0)\le {b_*\over a_*}={19\over20},
\qquad S_j(h_i^0)\le(r-1)R_j(h_i^0)\le rR_j(h_i^0).
\tag{35}
\]

\begin{proof}
The three target coordinates $X_U,X_V,X_W$ are independent standard
Gaussians because the three raw mode vectors are independent and every
realized target column is unit. Lemma~\ref{lem:step-002-target-mass}
therefore gives

\[
c_{\rm tar}^3r^{-5/3}(\log r)^{-3/2}
\le\Pr(W_{ij}\mid\mathcal F_{\rm sm})
\le C_{\rm tar}^3r^{-5/3}(\log r)^{-3/2}.
\tag{36}
\]

For the lower bound, do not factor $W_{ij}$ and $C_{ij}$. Instead,
Proposition~\ref{prop:step-002-conditional-competitors} applies uniformly
at every signed point $x$ in the target window, so disintegration gives

\[
\begin{aligned}
\Pr(E_{{\rm win},ij}\mid\mathcal F_{\rm sm})
&=\int_{x\in W_{ij}}
 \Pr(C_{ij}\mid X=x,\mathcal F_{\rm sm})\,
 d\Pr_X(x\mid\mathcal F_{\rm sm})\\
&\ge {1\over10}\Pr(W_{ij}\mid\mathcal F_{\rm sm}).
\end{aligned}
\tag{37}
\]

The upper bound follows from $E_{{\rm win},ij}\subseteq W_{ij}$. Equations
(33), (36), and (37) prove (34). Assumption~\ref{assump:random-initialization}
makes the conditional law identical across slots, proving the slot claim.
The minimum over finitely many targets preserves both uniform bounds; no
claim of independence among the target events is made.

Finally, normalization of $\xi_i^{(M)}$ cancels in every same-mode ratio.
On $E_{{\rm win},ij}$, each target denominator has squared product at least
$t_r^2=a_*\log r$, while (22) bounds every corresponding competitor
numerator product by $b_*\log r$. Thus each of the three pair ratios is at
most $b_*/a_*=19/20$, proving the first part of (35). Every sum defining
$S_j$ has at most $r-1$ terms, each at most $R_j$, proving the second part.


\end{proof}
\end{proposition}
\begin{proof}
Fix any realized tensor in the event $E_{\rm sm}$ from Proposition~\ref{prop:step-001-geometry}.
For every target $j$ and slot $i$, define $E_{{\rm win},ij}$ by (21),
(22), and (32), using only the original raw proposal coordinates. The
regression variables in Lemma~\ref{lem:step-002-regression} are proof-local
and do not alter this event.

Lemma~\ref{lem:step-002-target-mass} supplies the exact three-mode target
mass. Lemma~\ref{lem:step-002-product-reserve} and
Proposition~\ref{prop:step-002-conditional-competitors} show, with displayed
constants, that imposing every competitor pair constraint retains at least a
$1/10$ fraction after conditioning on any point of the target window.
Proposition~\ref{prop:step-002-joint-window} therefore proves

\[
{1\over10}\left[\sqrt{2\over\pi a_*}
e^{-1-1/(2a_*\log3)}\right]^3
r^{-5/3}(\log r)^{-3/2}
\le p_{\rm win}
\le
\left({2\over\pi a_*}\right)^{3/2}
r^{-5/3}(\log r)^{-3/2}.
\tag{38}
\]

This is the displayed probability claim. It is a one-slot probability
conditional on the fixed good instance, contains no confidence parameter or
rank choice, and uses neither target/competitor independence nor independence
among events for different targets. Equation (35) records the same-event
ratio interface needed after the downstream coverage step.
\end{proof}

\subsection{All-target coverage}\label{app:step-003}

Proposition~\ref{prop:step-002-joint-window} supplies, uniformly in slot
$i$ and target $j$ conditional on the fixed smoothed instance,
\[
p_{{\rm win},j}:=
\Pr(E_{{\rm win},ij}\mid\mathcal F_{\rm sm})
\in\left[c_{\rm win}r^{-5/3}(\log r)^{-3/2},
C_{\rm win}r^{-5/3}(\log r)^{-3/2}\right].
\tag{1}
\]
It also gives the ratio initialization on $E_{{\rm win},ij}$. Below we use
$1-x\le e^{-x}$ for $x\in[0,1]$ and the finite union bound; neither step
requires independence among target-indexed events in a common slot.

Fix an arbitrary realized smoothed instance for which
Proposition~\ref{prop:step-002-joint-window} holds, and condition on its
$\sigma$-field $\mathcal F_{\rm sm}$. Write

\[
\mathscr E_{ij}:=E_{{\rm win},ij},\qquad
p_j:=p_{{\rm win},j}=\Pr(\mathscr E_{ij}\mid\mathcal F_{\rm sm}).
\tag{2}
\]

The event \(\mathscr E_{ij}\) is evaluated from the raw Gaussian tape in slot
\(i\), with the target index used only to state the analysis event.
Proposition~\ref{prop:step-002-joint-window} gives
\(p_j\ge c_{\rm win}r^{-5/3}(\log r)^{-3/2}\) in (1).


\begin{lemma}[Fixed-target slot-miss bound]\label{lem:step-003-slot-miss}
Under Assumption~\ref{assump:random-initialization} and
Proposition~\ref{prop:step-002-joint-window}, fix \(j\in[r]\) and define

\[
B_j:=\bigcap_{i=1}^k\mathscr E_{ij}^{\,c}.
\tag{3}
\]

Then, conditional on \(\mathcal F_{\rm sm}\),

\[
\Pr(B_j\mid\mathcal F_{\rm sm})=(1-p_j)^k
\le e^{-kp_j}
\le r^{-C_{\rm rank}c_{\rm win}}.
\tag{4}
\]

No relation among \(\mathscr E_{1j},\ldots,\mathscr E_{rj}\) for a common
slot is required.

\begin{proof}
For fixed \(j\), each \(\mathscr E_{ij}\) is a measurable function of the
raw triple \(\Xi_i\) and the fixed realized tensor. The primitive
initialization assumption makes the \(\Xi_i\)'s independent over \(i\), so
the indicators \(1_{\mathscr E_{ij}}\) are independent and have the common
conditional success probability \(p_j\). Therefore the first equality in
(4) holds. The elementary inequality in the displayed auxiliary fact gives the first
inequality. From (1) and the rank formula,

\[
kp_j\ \ge\ \left(C_{\rm rank}r^{5/3}(\log r)^{5/2}\right)
   \left(c_{\rm win}r^{-5/3}(\log r)^{-3/2}\right)
 =C_{\rm rank}c_{\rm win}\log r,
\tag{5}
\]

where the ceiling in \(k\) only increases the left side. Exponentiating
(5) proves the final inequality in (4). The proof has used independence
only in the slot index \(i\), and nowhere factors events with different
target labels. 


\end{proof}
\end{lemma}
\begin{proposition}[All-target coupon coverage]\label{prop:step-003-coupon}
Under Assumptions~\ref{assump:subquadratic-rank} and
\ref{assump:random-initialization}, and
Proposition~\ref{prop:step-002-joint-window}, define

\[
E_{\rm cov}:=\bigcap_{j=1}^r\bigcup_{i=1}^k\mathscr E_{ij}
   =\bigcap_{j=1}^r B_j^{\,c}.
\tag{6}
\]

Then, for every fixed smoothed instance in the event
$E_{\rm sm}$ supplied by Proposition~\ref{prop:step-001-geometry},

\[
\Pr(E_{\rm cov}\mid\mathcal F_{\rm sm})
\ge 1-r^{\,1-C_{\rm rank}c_{\rm win}}.
\tag{7}
\]

Choose the universal rank constant once and for all so that

\[
C_{\rm rank}c_{\rm win}\ge4.
\tag{8}
\]

Then, for all \(r\ge3\),

\[
\Pr(E_{\rm cov}\mid\mathcal F_{\rm sm})\ge1-r^{-3}\ge
p_0,\qquad p_0:=1-3^{-3}={26\over27}>0.
\tag{9}
\]

No independence across the \(r\) target-indexed events is asserted or used.

\begin{proof}
The complement of (6) is \(\bigcup_{j=1}^rB_j\). Applying the finite
union bound and Lemma~\ref{lem:step-003-slot-miss} gives

\[
\Pr(E_{\rm cov}^{\,c}\mid\mathcal F_{\rm sm})
\le\sum_{j=1}^r\Pr(B_j\mid\mathcal F_{\rm sm})
\le r\,r^{-C_{\rm rank}c_{\rm win}}
=r^{\,1-C_{\rm rank}c_{\rm win}},
\tag{10}
\]

which proves (7). Under (8), the exponent is at most \(-3\), and
\(r^{-3}\le3^{-3}\) for \(r\ge3\), proving (9). The events \(B_j\) may be
maximally correlated for different \(j\); (10) remains valid because it is a
union bound rather than a product calculation. 


\end{proof}
\end{proposition}
\begin{proposition}[Proof-only witnesses and label-free procedure interface]\label{prop:step-003-witness-interface}
Under the hypotheses of Proposition~\ref{prop:step-003-coupon}, on the
event \(E_{\rm cov}\), for every \(j\in[r]\) the finite nonempty set

\[
\mathcal I_j:=\{i\in[k]:\mathscr E_{ij}\text{ occurs}\}
\tag{11}
\]

has a proof-only canonical witness \(i(j):=\min\mathcal I_j\). The
conclusion of Proposition~\ref{prop:step-002-joint-window} then gives

\[
R_j(h_{i(j)}^0)\le{19\over20},\qquad
S_j(h_{i(j)}^0)\le rR_j(h_{i(j)}^0).
\tag{12}
\]

The construction of \(E_{\rm cov}\), \(\mathcal I_j\), and \(i(j)\) is an
analysis definition only: the SL-SC-JEP-ALS procedure does not read any
target index \(j\), any \(u_j,v_j,w_j\), any raw-coordinate window event, or
the proof-only witness. It executes the same slotwise Jacobi, score, and
clustering rules for every tape. Thus (6) is a generated probabilistic
interface and not a label-dependent acceptance test.

\begin{proof}
By (6), each union over \(i\) is nonempty on \(E_{\rm cov}\); finiteness of
\([k]\) makes the minimum in (11) well-defined. Equation (12) is exactly
the ratio conclusion of Proposition~\ref{prop:step-002-joint-window} applied
to the witness event \(\mathscr E_{i(j)j}\). The final statement follows by
inspection of the setting's procedure: target labels occur only in the
proof definitions of \(A_{\ell\mid j},B_{\ell\mid j},C_{\ell\mid j}\) and
\(E_{{\rm win},ij}\); no such quantities are computed by the algorithm. A
proof may use \(i(j)\) to organize the subsequent targetwise induction,
while the actual state evolution remains the unlabeled stored-state
trajectory. 

\end{proof}
\end{proposition}
\begin{proof}
Fix a realized smoothed instance for which
Proposition~\ref{prop:step-002-joint-window} holds. It supplies a uniform
per-target lower probability \(p_j\ge c_{\rm win}r^{-5/3}(\log r)^{-3/2}\).
Lemma~\ref{lem:step-003-slot-miss} uses the \(k\) independent raw slots to
make each target miss probability at most \(r^{-C_{\rm rank}c_{\rm win}}\).
Proposition~\ref{prop:step-003-coupon} takes a union over targets and, after
the single universal choice (8), produces

\[
E_{\rm cov}=\{\text{every target has at least one jointly observable
window slot}\},\qquad
\Pr(E_{\rm cov}\mid\mathcal F_{\rm sm})\ge p_0={26\over27}.
\tag{13}
\]

Proposition~\ref{prop:step-003-witness-interface} records the targetwise initial
ratio interface for any proof-only witness and explicitly separates it from
the label-free procedure. The proof never assumes independence among the
events for different targets in a slot or across the target index; only the
primitive slot independence is used. Therefore (13) is exactly the
displayed coverage output \(E_{\rm cov},p_0\), ready for the trajectory
induction and the restart conclusion of Proposition~\ref{prop:step-012-restart}.

For completeness, the rank scaling has no hidden confidence factor. With
\(C_{\rm rank}\) fixed universally and \(r\ge3\),

\[
C_{\rm rank}r^{5/3}(\log r)^{5/2}\le k
\le\left(C_{\rm rank}+{1\over3^{5/3}(\log3)^{5/2}}\right)
 r^{5/3}(\log r)^{5/2},
\tag{14}
\]

so \(k=\Theta(r^{5/3}(\log r)^{5/2})\). The upper feasibility and
subquadratic statement are already part of Assumption~\ref{assump:subquadratic-rank};
indeed \(k/r^2=O((\log r)^{5/2}/r^{1/3})=o(1)\) for fixed
\(C_{\rm rank}\). If an unconditional nested probability is desired, the
tower property and the instance bound give

\[
\Pr(E_{\rm sm}\cap E_{\rm cov})\ge
\Pr(E_{\rm sm})\,p_0\ge(1-\delta_{\rm sm})p_0,
\tag{15}
\]

but the public one-run initialization guarantee remains the
conditional bound (13), with \(\delta_{\rm init}\) reserved for independent
restarts.
\end{proof}

\subsection{Jacobi chart recurrence}\label{app:step-004}

Proposition~\ref{prop:step-003-witness-interface} supplies, on
\(E_{\rm cov}\), a witness for every target with \(R_0\le19/20\) and
\(S_0\le rR_0\). Proposition~\ref{prop:step-001-realized-gram} supplies, on
\(E_{\rm sm}\), unit diagonal Gram matrices
\(\mathsf K_M=G_M^\top G_M\) with absolute row and column off-diagonal mass at
most \(q_{\rm real}\). We repeatedly use the corresponding induced
$\ell_1$ bounds, the Neumann estimate
$\|K^{-1}\|\le(1-q)^{-1}$ when $\|K-I\|\le q<1$, and the normalization
inequality
\[
\left\|\frac a{\|a\|_2}-\frac b{\|b\|_2}\right\|_2
\le\frac{2\|a-b\|_2}{\min\{\|a\|_2,\|b\|_2\}}
\]
for nonzero $a,b$.

Fix \(E_{\rm sm}\), \(E_{\rm cov}\), a target \(j\), and one proof-only witness
slot \(i(j)\). Suppress the slot index. Put

\[
q:=q_{\rm real}\le q_*={1\over4096},\qquad
\gamma:=\Gamma\le {101\over100}.
\]

For \(M\in\{U,V,W\}\), write \(g_{M,\ell}\) for the realized target column
(\(u_\ell,v_\ell,w_\ell\), respectively), and \(h_M^t\) for the mode-\(M\)
state after \(t\) simultaneous commits. At every time at which the target
correlation is nonzero, use a proof-only sign so that

\[
a_{M,j}^t:=\langle g_{M,j},h_M^t\rangle>0,\qquad
x_{M,\ell}^t:={\langle g_{M,\ell},h_M^t\rangle\over a_{M,j}^t}
\quad(\ell\ne j).
\tag{1}
\]

The absolute values of (1) are the setting's \(A_\ell,B_\ell,C_\ell\).
Sign choices do not alter \(R,S\) or \(\zeta\), because \(\zeta\) minimizes
over a modewise sign.


\begin{lemma}[Simultaneous coefficient and pair-mass recurrence]\label{lem:step-004-recurrence}
Under the realized Gram conclusion of Proposition~\ref{prop:step-001-geometry}
and the covered-slot initial interface of Proposition~\ref{prop:step-003-witness-interface},
suppose at
time \(t\) all target correlations in (1) are nonzero and
\(\gamma R_t<1\). For mode \(M\), with held modes \(M',M''\), define

\[
\beta_{M,j}^t:=1,\qquad
\beta_{M,\ell}^t:={\lambda_\ell\over\lambda_j}
x_{M',\ell}^t x_{M'',\ell}^t\quad(\ell\ne j).
\tag{2}
\]

Then the target coefficient of the unnormalized \(M\)-contraction is nonzero,
and, with \(D_t:=1-\gamma qR_t\),

\[
|x_{M,\ell}^{t+1}|\le {\gamma R_t+q\over D_t}\quad(\ell\ne j),
\tag{3}
\]

\[
R_{t+1}\le\left({\gamma R_t+q\over1-\gamma qR_t}\right)^2,\qquad
S_{t+1}\le
{(\gamma R_t+q)\{\gamma(1+q)S_t+q\}\over(1-\gamma qR_t)^2}.
\tag{4}
\]

The conclusion holds for the first update even when \(h^0\) has a component
orthogonal to a target span.

\begin{proof}
Divide the \(M\)-contraction by its target
coefficient \(\lambda_j a_{M',j}^ta_{M'',j}^t\). Its coefficient vector in
the \(g_{M,\ell}\) basis is \(e_j+\beta_M^t\). Since
\(|\beta_{M,\ell}^t|\le\gamma B_\ell^tC_\ell^t\le\gamma R_t<1\), the target
coordinate of the contraction obeys

\[
\begin{aligned}
\left|\left\langle g_{M,j},e_j+
\sum_{\ell\ne j}\beta_{M,\ell}^tg_{M,\ell}\right\rangle\right|
&\ge1-\sum_{\ell\ne j}|\beta_{M,\ell}^t|
|\langle g_{M,j},g_{M,\ell}\rangle|\\
&\ge1-\gamma R_tq=D_t>0.
\end{aligned}
\tag{5}
\]

Thus the contraction norm is positive and its normalized output is defined.
For an off-target index \(\ell\), the diagonal term contributes
\(|\beta_{M,\ell}^t|\), and all remaining terms are bounded by the Gram row
mass \(q\), because every \(|\beta_{M,m}^t|\le1\). Therefore

\[
|\langle g_{M,\ell},e_j+\sum_m\beta_{M,m}^tg_{M,m}\rangle|
\le\gamma R_t+q.
\tag{6}
\]

Dividing (6) by (5) proves (3), and taking products proves the first part of
(4). For the \(S\) bound, let \(n_{M,\ell}\) denote the absolute numerator
in (6). Using the target column mass \(q\), then swapping the remaining finite
sums and using the Gram column mass \(q\),

\[
\begin{aligned}
\sum_{\ell\ne j}n_{M,\ell}
&\le \sum_{\ell\ne j}|\beta_{M,\ell}^t|
 +q+q\sum_{\ell\ne j}|\beta_{M,\ell}^t|\\
&\le\gamma(1+q)\sum_{\ell\ne j}
|x_{M',\ell}^tx_{M'',\ell}^t|+q\\
&\le\gamma(1+q)S_t+q.
\end{aligned}
\tag{7}
\]

For example, for the \(U,V\) pair, the \(U\)-numerator is at most
\(\gamma R_t+q\) termwise, while the \(\ell_1\) sum of the \(V\)-numerator
is bounded by (7). Division by the two target denominators gives the first
pair sum in \(S_{t+1}\); cyclic permutation gives the other pair sums. At
\(t=0\), the window event supplies nonzero target correlations and (5)
proves the first update before any span assertion. After that update every
mode lies in its exact target span because the tensor contraction is a linear
combination of the \(g_{M,\ell}\). 


\end{proof}
\end{lemma}
\begin{proposition}[Noncircular scalar envelope and finite chart entry]\label{prop:step-004-envelope}
Under Assumption~\ref{assump:subquadratic-rank}, on the coverage interface
of Proposition~\ref{prop:step-003-witness-interface}, and under
Lemma~\ref{lem:step-004-recurrence}, put

\[
\bar R:={19\over20},\quad d_*:={999\over1000},\quad
\rho_R=\rho_S:={49\over50}.
\tag{8}
\]

Then every finite trajectory iterate satisfies

\[
R_t\le\bar R,\qquad 1-\gamma qR_t\ge d_*>0,\qquad \gamma R_t<1.
\tag{9}
\]

Moreover, if \(\Lambda_t:=\max_M\sum_{\ell\ne j}|x_{M,\ell}^t|\), then

\[
\Lambda_{t+1}\le {\gamma(1+q)S_t+q\over1-\gamma qR_t},\qquad
S_{t+1}\le\rho_S S_t+q.
\tag{10}
\]

There is a universal choice \(C_{\rm burn}=2048\) such that, with
\(L_{\rm burn}=\lceil C_{\rm burn}\log r\rceil\), every covered witness
trajectory has

\[
R_{L_{\rm burn}-1}\le {1\over128},\quad
S_{L_{\rm burn}-1}\le {1\over7200},\quad
\Lambda_{L_{\rm burn}}\le {1\over256},\quad
\max_{M,\ell}|x_{M,\ell}^{L_{\rm burn}}|\le {1\over100}.
\tag{11}
\]

\begin{proof}
Let
\(f(x)=((\gamma x+q)/(1-\gamma qx))^2\). On
\([0,\bar R]\), \(f\) is increasing. The numerical bounds

\[
\gamma\bar R+q < {24\over25},\qquad
1-\gamma q\bar R>{999\over1000}
\tag{12}
\]

give \(f(\bar R)/\bar R<49/50\). For \(x\ge q/\gamma\), direct
differentiation shows \(f(x)/x\) is increasing, since

\[
{d\over dx}\log{f(x)\over x}
={2\gamma\over\gamma x+q}-{1\over x}
 +{2\gamma q\over1-\gamma qx}>0.
\tag{13}
\]

The covered initial state satisfies \(\gamma R_0\le(101/100)(19/20)<1\),
so Lemma~\ref{lem:step-004-recurrence} is applicable at \(t=0\).
Thus \(f(x)\le(49/50)x\) whenever \(q/\gamma\le x\le\bar R\); if
\(x<q/\gamma\), then \(x<1/128\) and monotonicity gives
\(f(x)\le f(q/\gamma)<f(\bar R)<\bar R\). Since \(R_0\le\bar R\),
induction using
(4) proves (9), and the first phase reaches \(R_t\le1/128\) in at most

\[
N_R:=\left\lceil{\log(128\bar R)\over\log(50/49)}\right\rceil
\tag{14}
\]

steps. The coefficient multiplying \(S_t\) in (4) is bounded by

\[
{(24/25)(101/100)(1001/1000)\over(999/1000)^2}
< {49\over50},
\tag{15}
\]

and the additive term is less than \(q\). Hence

\[
S_t\le(49/50)^t(19r/20)+50q.
\tag{16}
\]

Set

\[
N_S(r):=\left\lceil
{\log\{(19r/20)/(1/64-50q_*)\}\over\log(50/49)}
\right\rceil.
\tag{17}
\]

The denominator is positive because \(50q_*=50/4096<1/64\). After
\(N_S(r)\) steps, (16) gives \(S_t\le1/64\). At any time with
\(R_t\le1/128\), (4) and the same numerical constants give

\[
R_{t+1}\le10^{-4},\qquad S_{t+1}\le {1\over7200},\qquad
\Lambda_{t+1}< {1\over60}.
\tag{18}
\]

Applying (10) once more from the \(S\)-bound \(1/7200\) gives
\(\Lambda_{t+2}\le1/256\), while (3) gives the \(1/100\) coordinate bound.
Once \(R\le1/128\) and \(S\le1/7200\), the right side of the \(S\)
recurrence is again below \(1/7200\), so this smaller bound persists.
Consequently (11) holds once
\(L_{\rm burn}\ge\max\{N_R,N_S(r)\}+2\). A direct numerical comparison at
\(r=3\), and the linear \(\log r\) term in (17), gives

\[
\max\{N_R,N_S(r)\}+2\le2048\log r\qquad(r\ge3).
\tag{19}
\]

Thus the prescribed universal \(C_{\rm burn}=2048\) suffices. All bounds are
finite-horizon consequences of (4); no chart or denominator condition was
assumed in obtaining them. 


\end{proof}
\end{proposition}
\begin{lemma}[Target-chart self-map and contraction]\label{lem:step-004-chart}
On the event supplied by Proposition~\ref{prop:step-001-geometry}, let
\(\mathsf K_M:=G_M^\top G_M\) denote the step-local Gram (distinct from the
setting-defined same-state coefficient matrix \(K_M\)).
For a target \(j\), write a signed off-target ratio vector
\(x_M\in\mathbb R^{r-1}\), with \(e_j+x_M\) denoting the \(r\)-vector obtained by inserting a
1 in coordinate \(j\). On the chart

\[
\mathcal C_j:=\left\{(x_U,x_V,x_W):
\max_M\|x_M\|_\infty\le {1\over100},\
\max_M\|x_M\|_1\le {1\over256},\
S_j\le {1\over512}\right\},
\tag{20}
\]

the exact signed ratio map is

\[
\Phi_M(x_{M'},x_{M''})_\ell
={\left[\mathsf K_M(e_j+\beta_M)\right]_\ell\over
\left[\mathsf K_M(e_j+\beta_M)\right]_j},\qquad
\beta_{M,j}=1,\quad
\beta_{M,\ell}={\lambda_\ell\over\lambda_j}
x_{M',\ell}x_{M'',\ell},
\tag{21}
\]

and is a self-map of \(\mathcal C_j\). In the metric

\[
d_{\rm ch}(x,\widetilde x):=\max_M\|x_M-\widetilde x_M\|_1,
\tag{22}
\]

it satisfies

\[
d_{\rm ch}(\Phi(x),\Phi(\widetilde x))\le {1\over4}
d_{\rm ch}(x,\widetilde x).
\tag{23}
\]

If \(p_M(x)\) is the unit vector in \({\rm span}(G_M)\) with correlation
ratios \(x_M\) and positive target correlation, then

\[
\|p_M(x)-p_M(\widetilde x)\|_2\le4\|x_M-\widetilde x_M\|_1.
\tag{24}
\]

\begin{proof}
Let \(v=e_j+x\), and define

\[
p_M(x):={G_M\mathsf K_M^{-1}v\over
\|G_M\mathsf K_M^{-1}v\|_2}.
\tag{25}
\]

The Neumann bound applies because \(q<1\), and substitution gives
\(G_M^\top p_M(x)\) proportional to \(v\); hence (25) is the unique oriented
span direction with the stated ratios. The Gram eigenvalues lie in
\([1-q,1+q]\). Consequently
\[
\|G_M\mathsf K_M^{-1}(x-\widetilde x)\|_2
\le(1-q)^{-1/2}\|x-\widetilde x\|_2,
\quad
\|G_M\mathsf K_M^{-1}v\|_2\ge(1+q)^{-1/2}\|v\|_2.
\]
Since \(\|e_j+x\|_2\ge1\) on the chart, the normalization inequality gives
(24) with the conservative factor \(4\).

For (21), divide the held-mode target coefficients exactly as in (2); the
common target scale cancels. On (20),
\(\|\beta_{M,-j}\|_\infty\le\gamma/10^4\), and the target denominator in
(21) is at least \(999/1000\). The off-target numerator has coordinatewise
bound \(\gamma/10^4+q<1/100\). Its \(\ell_1\) bound is
\(\gamma(1+q)S_j+q<1/256\), by the row/column summation used in (7).
For every output pair, the product mass is at most
\((1/100)(1/256)<1/512\). Thus \(\Phi\) is a self-map.

For two chart points, with \(d=d_{\rm ch}(x,\widetilde x)\), the product rule
gives

\[
\|\beta_M-\widetilde\beta_M\|_1
\le2\gamma(1/100)d.
\tag{26}
\]

The off-target block of \(\mathsf K_M\) has induced column mass at most \(1+q\), the
target row has mass at most \(q\), and both quotient denominators are at least
\(999/1000\). Using the self-map \(\ell_1\) bound in the denominator variation
term yields

\[
\|\Phi_M(x)-\Phi_M(\widetilde x)\|_1
\le\left({1+q\over999/1000}+
{q\over(999/1000)^2\,256}\right)
\|\beta_M-\widetilde\beta_M\|_1
< {1\over4}d.
\tag{27}
\]

Taking the maximum over modes proves (23). 


\end{proof}
\end{lemma}
\begin{proposition}[Finite projective certificate and stored-state interface]\label{prop:step-004-certificate}
Under Assumption~\ref{assump:subquadratic-rank}, on the coverage interface
of Proposition~\ref{prop:step-003-witness-interface}, and under
Proposition~\ref{prop:step-004-envelope} and
Lemma~\ref{lem:step-004-chart},
on every covered witness trajectory the chart \(\mathcal C_j\) is entered no
later than \(L_{\rm burn}\), remains invariant, and for every
\(t\ge L_{\rm burn}\),

\[
\zeta(h^t)\le {1\over32}\,4^{-(t-L_{\rm burn})}.
\tag{28}
\]

With the universal choice \(C_{\rm cert}=64\) and
\(L_{\rm cert}=\lceil C_{\rm cert}\log r\rceil\),
\(\zeta(h^{L_{\rm burn}+L_{\rm cert}})\le\tau_r=q_*^2/(10^4r)\). Thus the
algorithm stores a state at or before \(L_{\rm prop}=L_{\rm burn}+L_{\rm cert}\);
the stored state is \(h^t\), not its look-ahead image. If the procedure stores
an earlier state because its observable \(\zeta\) is already below the
threshold, the continuation used in this proof is only a proof device and
the earlier stored state already satisfies the exported certificate.

\begin{proof}
By (11), its \(1/100\) and \(1/256\) bounds imply
that the state at \(L_{\rm burn}\) has
\(S_j\le(1/100)(1/256)<1/512\), so it is in \(\mathcal C_j\); the self-map
keeps all subsequent states there. Let
\(\Delta_t=d_{\rm ch}(x^t,x^{t-1})\). The first in-chart displacement is at
most \(2/256=1/128\). For \(t\ge L_{\rm burn}+1\), (23) gives the recursive
inequality below; the displayed closed form also holds at
\(t=L_{\rm burn}\) by the diameter bound:

\[
\Delta_{t+1}\le {1\over4}\Delta_t,\qquad
\Delta_{t+1}\le {1\over128}\,4^{-(t-L_{\rm burn})}.
\tag{29}
\]

For each mode, the oriented state at time \(t\) is \(p_M(x_M^t)\), while the
oriented look-ahead Jacobi state is \(p_M(x_M^{t+1})\). Equation (24) and the
sign-minimum in \(\zeta_M\) give

\[
\zeta_M(h^t)\le4\|x_M^t-x_M^{t+1}\|_1,
\]

which with (29) proves (28). Finally,
\(4^{-L_{\rm cert}}/32\le\tau_r\) holds for all \(r\ge3\) under \(C_{\rm cert}=64\):
after substituting \(\tau_r=q_*^2/(10^4r)\), the required logarithm is
\(\log(32\cdot10^4r/q_*^2)/\log4\), and \(64\log r\) dominates it at \(r=3\)
and hence for all larger \(r\). If the bound first holds at \(t=L_{\rm prop}\),
the algorithm still evaluates \(\zeta(h^t)\) and stores that old state; no
look-ahead state is exported. 

\end{proof}
\end{proposition}
\begin{proof}
On \(E_{\rm sm}\cap E_{\rm cov}\), choose the witness from Proposition~\ref{prop:step-003-witness-interface} \(i(j)\) for
each target \(j\). The argument in fact applies to every pair \((i,j)\) for
which the window event occurs; \(E_{\rm cov}\) is used only to
ensure that at least one such pair exists for every target. The initial interface gives nonzero target
coordinates and \(R_0\le19/20\), \(S_0\le rR_0\). Lemma~\ref{lem:step-004-recurrence}
proves the exact displayed \(R,S\) recurrences, including the first
simultaneous update, and proves every target denominator before division.
Proposition~\ref{prop:step-004-envelope} gives the noncircular invariant,
finite burn entry, and modewise \(\ell_1\) chart bounds. The chart lemma
supplies an independent contraction mechanism for the same old-state map;
Proposition~\ref{prop:step-004-certificate} converts successive displacement
to the exact algorithmic Euclidean \(\zeta\) and guarantees a stored
certificate within the prescribed finite horizon. These conclusions define
\(E_{\rm RS},E_{\rm chart},E_{\rm chart}^{(1)},E_{\rm cert}\) as generated
outputs. No target label, comparator, future ALS output, or landing output is
used by the procedure. A zero target denominator or contraction on an
uncovered tape remains an observable unsuccessful branch; on the covered
branch, (5) rules it out before the first commit.

Concretely, \(E_{\rm RS}\) is the finite intersection, over all the window pairs and all \(t\le L_{\rm prop}\), of (3)--(4) and (9);
\(E_{\rm chart}\) and \(E_{\rm chart}^{(1)}\) are the corresponding post-burn
bounds (20), and \(E_{\rm cert}\) records that each such slot stores its
first old state with \(\zeta\le\tau_r\). These are proof events, not extra
algorithmic tests.
\end{proof}

\subsection{Signed certificate equation}\label{app:step-005}

Proposition~\ref{prop:step-001-realized-gram} gives, on the event supplied by
Proposition~\ref{prop:step-001-geometry}, unit diagonal matrices
$H_M=G_M^\top G_M$ with off-diagonal row and column masses at most
$q_{\rm real}\le q_*$. Proposition~\ref{prop:step-004-certificate} says that
a covered slot stores its old state $h^t$, not its look-ahead image, with
$\zeta_M(h^t)\le\tau_r$, after all modes have entered the corresponding
realized target span. We also use
$\|(H-I)v\|_1\le q\|v\|_1$,
$\|x\circ y\|_1\le\|x\|_2\|y\|_2$, and
$\|G_M^\top z\|_2\le\sqrt{1+q}\|z\|_2$, each instantiated below.

Throughout, let $m_c:=|\mathcal I_{\rm cert}|$, index certified slots by
$a\in\mathcal I_{\rm cert}$, and put
$g_{U,j}=u_j,\ g_{V,j}=v_j,\ g_{W,j}=w_j$.  For a slot $a$, write
$p_{M,a}$ for its three unit vectors and define the raw correlation vector

\[
c_{M,a}:=G_M^\top p_{M,a},\qquad
k_{M,a}:=D_\lambda\bigl(c_{M',a}\circ c_{M'',a}\bigr),
\tag{1}
\]

where $M',M''$ are the two held modes.  The symbols in (1) are slot-local;
the setting's $K_M,D_M$ for selected representatives are recovered by
restricting the slot index set to the selected $r$ columns.


\begin{lemma}[Signed certificate equation and score-preserving orientation]\label{lem:step-005-signed-equation}
Under the certified-state conclusion of Proposition~\ref{prop:step-004-certificate}
and the static event of Proposition~\ref{prop:step-001-geometry}, for every
certified slot (a) and mode (M), let

\[
d_{M,a}:=\|G_Mk_{M,a}\|_2>0,
\]

and choose $\varepsilon_{M,a}\in\{\pm1\}$ attaining the certified
projective residual.  Then, with

\[
r_{M,a}:=d_{M,a}p_{M,a}-\varepsilon_{M,a}G_Mk_{M,a},
\]

\[
\|r_{M,a}\|_2\le d_{M,a}\tau_r,
\qquad
d_{M,a}p_{M,a}=\varepsilon_{M,a}G_Mk_{M,a}+r_{M,a}.
\tag{2}
\]

Writing $\theta_a:=\langle T,p_{U,a}\otimes p_{V,a}\otimes p_{W,a}\rangle$,
all three signs satisfy
\[
\varepsilon_{U,a}=\varepsilon_{V,a}=\varepsilon_{W,a}
=\operatorname{sgn}(\theta_a),
\qquad \theta_a\ne0,
\tag{3}
\]

and, with $s_a:=\operatorname{sgn}(\theta_a)$, the proof-only orientation

\[
\bar p_{U,a}:=s_ap_{U,a},\qquad
\bar p_{V,a}:=p_{V,a},\qquad
\bar p_{W,a}:=p_{W,a}
\tag{4}
\]

has
\[
|\theta_a|\bar p_{U,a}\otimes\bar p_{V,a}\otimes\bar p_{W,a}
=\theta_a p_{U,a}\otimes p_{V,a}\otimes p_{W,a}.
\]
Defining

\[
\bar c_{M,a}:=G_M^\top\bar p_{M,a},\quad
\bar k_{M,a}:=D_\lambda(\bar c_{M',a}\circ\bar c_{M'',a}),\quad
\bar r_{M,a}:=\eta_{M,a}r_{M,a},
\]

where $(\eta_{U,a},\eta_{V,a},\eta_{W,a})=(s_a,1,1)$, gives the
right-sided barred equation
\[
d_{M,a}\bar p_{M,a}=G_M\bar k_{M,a}+\bar r_{M,a},
\qquad \|\bar r_{M,a}\|_2\le d_{M,a}\tau_r.
\tag{5}
\]

Equivalently, if \(P_M^{\rm raw}\), \(K_M^{\rm raw}\), and \(R_M^{\rm raw}\)
stack the columns \(p_{M,a}\), \(k_{M,a}\), and \(r_{M,a}\), and if
\(D_M^{\rm cert}:=\operatorname{diag}(d_{M,a})\) and
\(\Sigma_M:=\operatorname{diag}(\varepsilon_{M,a})\), then
\[
P_M^{\rm raw}D_M^{\rm cert}
=G_MK_M^{\rm raw}\Sigma_M+R_M^{\rm raw},\qquad
\|R_M^{\rm raw}(:,a)\|_2\le d_{M,a}\tau_r.
\tag{5a}
\]
The barred stacks satisfy
\[
\bar P_MD_M^{\rm cert}=G_M\bar K_M+\bar R_M,\qquad
\|\bar R_M(:,a)\|_2\le d_{M,a}\tau_r.
\tag{5b}
\]

In particular, writing $\sigma_a:=|\theta_a|$,
\[
|d_{M,a}-\sigma_a|\le d_{M,a}\tau_r,
\qquad
\frac{\sigma_a}{d_{M,a}}\in[1-\tau_r,1+\tau_r].
\tag{6}
\]

\begin{proof}
The stored-state interface gives a nonzero
contraction and
\[
\left\|p_{M,a}-\varepsilon_{M,a}
\frac{G_Mk_{M,a}}{d_{M,a}}\right\|_2\le\tau_r.
\]
Multiplying by $d_{M,a}$ proves (2).  By multilinearity,

\[
p_{M,a}^{\top}G_Mk_{M,a}
=\sum_{j=1}^r\lambda_j
\langle g_{U,j},p_{U,a}\rangle
\langle g_{V,j},p_{V,a}\rangle
\langle g_{W,j},p_{W,a}\rangle
=\theta_a.
\]

Taking the inner product of (2) with $p_{M,a}$ therefore gives

\[
d_{M,a}-\varepsilon_{M,a}\theta_a
=p_{M,a}^{\top}r_{M,a},
\qquad
|d_{M,a}-\varepsilon_{M,a}\theta_a|\le d_{M,a}\tau_r.
\tag{7}
\]

Because $\tau_r=q_*^2/(10^4r)<1$, (7) implies
$\varepsilon_{M,a}\theta_a>0$, proving (3) simultaneously for all modes.
For the orientation (4), the held-mode sign products are

\[
\bar k_{U,a}=k_{U,a},\qquad
\bar k_{V,a}=s_ak_{V,a},\qquad
\bar k_{W,a}=s_ak_{W,a}.
\]

Multiplying the $U$ equation in (2) by $s_a$, and leaving the other two
equations unchanged, gives (5); the rank-one identity follows because the
product of the three orientation signs is $s_a$.  Finally,
\[
\bar p_{M,a}^{\top}\bar r_{M,a}
=d_{M,a}-\bar p_{M,a}^{\top}G_M\bar k_{M,a}
=d_{M,a}-\sigma_a,
\]
so (6) follows from the residual bound.  This common scalar identity is the
normalization cancellation: the three modes share one $\sigma_a$, and no
independent first-order scalar error is charged in (5).  


\end{proof}
\end{lemma}
\begin{proposition}[Projected Gram-leak and finite residual ledger]\label{prop:step-005-projected-ledger}
Under the hypotheses of
Lemma~\ref{lem:step-005-signed-equation}, fix a barred certified slot $a$.
Put $H_M:=G_M^\top G_M$, $\bar c_M:=\bar c_{M,a}$,
$\bar k_M:=\bar k_{M,a}$, $d_M:=d_{M,a}$, and
$\bar r_M:=\bar r_{M,a}$.  For every selector
$w=(w_1,\ldots,w_r)\in[0,1]^r$,
\[
\sum_{j=1}^r w_j
\left|d_M\bar c_{M,j}-\bar k_{M,j}\right|
\le q\|\bar k_M\|_1+\sqrt{r(1+q)}\,d_M\tau_r,
\tag{8}
\]

and, in particular,
\[
\left|d_M\bar c_{M,j}-\lambda_j\bar c_{M',j}\bar c_{M'',j}\right|
\le q\lambda_{\max}+d_M\tau_r.
\tag{9}
\]

Every stored state is in $\operatorname{range}(G_M)$, so if
$\bar p_{M,a}=G_M\alpha_{M,a}$, then with
$\Delta_{M,a}:=(H_M-I)\alpha_{M,a}$ there is an exact decomposition
\[
d_M\alpha_{M,a}
=B_{M,a}+L_{M,a}+X_{M,a}+\varrho_{M,a},
\tag{10}
\]
where
\[
\begin{aligned}
B_{M,a}&:=D_\lambda(\alpha_{M',a}\circ\alpha_{M'',a}),\\
L_{M,a}&:=D_\lambda(\alpha_{M',a}\circ\Delta_{M'',a}
 +\Delta_{M',a}\circ\alpha_{M'',a}),\\
X_{M,a}&:=D_\lambda(\Delta_{M',a}\circ\Delta_{M'',a}),
\end{aligned}
\]

and
\[
\|L_{M,a}\|_1\le\frac{2\lambda_{\max}q}{1-q},\qquad
\|X_{M,a}\|_1\le\frac{\lambda_{\max}q^2}{1-q},\qquad
\|\varrho_{M,a}\|_1
\le\frac{\sqrt r\,d_{M,a}\tau_r}{\sqrt{1-q}}.
\tag{11}
\]

The first normalized identity is
\[
\alpha_{M,a}-\frac{B_{M,a}}{d_{M,a}}
=\frac{L_{M,a}+X_{M,a}+\varrho_{M,a}}{d_{M,a}},
\tag{12}
\]
so the scalar normalization term is not double-counted.  If $d_{M,a}$ is
replaced by $\sigma_a$, the exact identity is
\[
\alpha_{M,a}-\frac{B_{M,a}}{\sigma_a}
 =\frac{L_{M,a}+X_{M,a}+\varrho_{M,a}}{d_{M,a}}
  +\left(\frac1{d_{M,a}}-\frac1{\sigma_a}\right)B_{M,a}.
\tag{12a}
\]
The additional term therefore has $\ell_1$ norm at most
$\tau_r(1-\tau_r)^{-1}\|B_{M,a}\|_1/d_{M,a}$.

\begin{proof}
Projecting (5) by $G_M^\top$ gives the exact
coordinate identity
\[
d_M\bar c_M=H_M\bar k_M+G_M^\top\bar r_M,
\qquad
d_M\bar c_M-\bar k_M=(H_M-I)\bar k_M+G_M^\top\bar r_M.
\tag{13}
\]
The row and column bounds from Proposition~\ref{prop:step-001-geometry} imply
$\|(H_M-I)\bar k_M\|_1\le q\|\bar k_M\|_1$.  Also
$\|G_M^\top\bar r_M\|_1\le\sqrt r\|G_M^\top\bar r_M\|_2
\le\sqrt{r(1+q)}\,d_M\tau_r$, proving (8).  Since all correlations are inner
products of unit vectors, $|\bar c_{M,j}|\le1$, which proves (9).

For the coefficient form, write $\bar p_M=G_M\alpha_M$.  The Gram
eigenvalue bounds and $\bar p_M^\top\bar p_M=1$ give
\[
\|\alpha_M\|_2\le(1-q)^{-1/2},\qquad
\|\Delta_M\|_2\le q(1-q)^{-1/2},\qquad
\bar c_M=\alpha_M+\Delta_M.
\tag{14}
\]
Substitution of (14) into
$\bar k_M=D_\lambda(\bar c_{M'}\circ\bar c_{M''})$ gives $B+L+X$.
Because both sides of (5) are in $\operatorname{range}(G_M)$, the residual
is in that range and $\varrho_M:=G_M^\dagger\bar r_M$ satisfies
$\|\varrho_M\|_2\le d_M\tau_r/\sqrt{1-q}$.  The Cauchy--Schwarz
product inequality yields (11), and applying $G_M^\dagger$ to (5) gives
(10)--(12).  The final scalar conversion follows from (6). 


\end{proof}
\end{proposition}
\begin{proposition}[Observable score floor and threshold/tail ledger]\label{prop:step-005-threshold-ledger}
Under the certified-state conclusion of Proposition~\ref{prop:step-004-certificate}
and its coverage-witness interface, define
\[
a_{\rm win}:=\sqrt{\frac{1-q_*}{1+1/65536}},\qquad
u_{\rm win}:=\frac1{100\,256^2},\qquad
s_{\rm win}:=a_{\rm win}^3(1-1.01u_{\rm win}).
\tag{15}
\]
Then $a_{\rm win}>99/100$ and $s_{\rm win}>9/10$.  For every coverage
witness for target $j$, $\sigma_a\ge s_{\rm win}\lambda_j$, hence
\[
\sigma_{\max}\ge s_{\rm win}\lambda_{\min}>\frac9{10}\lambda_{\min}.
\tag{16}
\]
Let
\[
\mathcal H:=\{a\in\mathcal I_{\rm cert}:\sigma_a\ge0.85\sigma_{\max}\}
\tag{17}
\]
be the observable retained pool.  For each $a\in\mathcal H$,
\[
\sigma_a>\frac34\lambda_{\min},\qquad
d_{M,a}>\frac34\lambda_{\min},\qquad
\frac{\lambda_{\max}}{d_{M,a}}<\frac75,\qquad
\max_j|\bar c_{M,a,j}|>\frac34
\tag{18}
\]
for every mode $M$.  Put
\[
\eta_{\rm raw}(r):=\frac75(1+q_*)q_*+
\sqrt{r(1+q_*)}\,\tau_r,
\]
\[
\eta_{\rm score}(r):=\eta_{\rm raw}(r)+
\frac{\frac75(1+q_*)\tau_r}{1-\tau_r}.
\tag{19}
\]
Then $\eta_{\rm score}(r)<1/2000$ for every $r\ge3$.  For every
$a\in\mathcal H$, every mode $M$, and every selector $I\subseteq[r]$,
\[
\sum_{j\in I}\left|
\bar c_{M,a,j}-\frac{\lambda_j}{\sigma_a}
\bar c_{M',a,j}\bar c_{M'',a,j}
\right|\le\eta_{\rm score}(r).
\tag{20}
\]
Consequently, with the proof-only latent-coordinate $1/8$-threshold sets
\[
T_{M,a}:=\{j:|\bar c_{M,a,j}|\ge1/8\},
\tag{21}
\]
each $T_{M,a}$ is nonempty, and if $j_M$ maximizes
$|\bar c_{M,a,j}|$ in mode $M$, then
\[
\frac{\lambda_{j_M}}{\sigma_a}
|\bar c_{M',a,j_M}\bar c_{M'',a,j_M}|
>\frac34-\eta_{\rm score}(r),
\tag{22}
\]
and $j_M\in T_{M',a}\cap T_{M'',a}$.  Finally, for any mode $M$, any
$t\in[0,1]$, and $I_{M,a}(t):=\{j:|\bar c_{M,a,j}|<t\}$,
\[
\sum_{j\in I_{M,a}(t)}\lambda_j
|\bar c_{U,a,j}\bar c_{V,a,j}\bar c_{W,a,j}|
\le\lambda_{\max}(1+q_*)t.
\tag{23}
\]
In particular, the $1/8$ tail is at most
$\lambda_{\max}(1+q_*)/8$.  For a coverage witness $j$, the sharper
all-certified tail is
\[
\sum_{\ell\ne j}\lambda_\ell
|\bar c_{U,a,\ell}\bar c_{V,a,\ell}\bar c_{W,a,\ell}|
\le\lambda_{\max}|\bar c_{U,a,j}\bar c_{V,a,j}\bar c_{W,a,j}|u_{\rm win}.
\tag{24}
\]
Equations (8), (20), (21), (23), and (24), taken for every certified slot
and mode, define the finite event the tail conclusion of Proposition~\ref{prop:step-005-threshold-ledger}; it is formed before any
target label, representative ordering, or clustering operation.

\begin{proof}
On a coverage witness, let
$\rho_{M,\ell}:=|\bar c_{M,a,\ell}/\bar c_{M,a,j}|$.
Proposition~\ref{prop:step-004-certificate} gives
$\sum_{\ell\ne j}\rho_{M,\ell}\le1/256$ and
$\max_\ell\rho_{M,\ell}\le1/100$.  Since the stored state is in the
target span, (14) and the Gram floor give
$\|\bar c_M\|_2^2\ge1-q_*$.  Therefore
\[
|\bar c_{M,a,j}|
\ge\frac{\sqrt{1-q_*}}
{\sqrt{1+\sum_{\ell\ne j}\rho_{M,\ell}^2}}
\ge a_{\rm win}.
\tag{25}
\]
The off-target score sum is bounded by
\[
\sum_{\ell\ne j}\lambda_\ell
|\bar c_{U,a,\ell}\bar c_{V,a,\ell}\bar c_{W,a,\ell}|
\le\lambda_{\max}|\bar c_{U,a,j}\bar c_{V,a,j}\bar c_{W,a,j}|
\sum_{\ell\ne j}\rho_{U,\ell}\rho_{V,\ell}\rho_{W,\ell}
\le\lambda_{\max}|\bar c_{U,a,j}\bar c_{V,a,j}\bar c_{W,a,j}|u_{\rm win}.
\tag{26}
\]
The reverse triangle inequality applied to the target summand and the
off-target sum in (26) gives
$\sigma_a\ge |\bar c_{U,a,j}\bar c_{V,a,j}\bar c_{W,a,j}|
(\lambda_j-\lambda_{\max}u_{\rm win})
\ge a_{\rm win}^3\lambda_j(1-1.01u_{\rm win})$,
which is (16). The displayed numerical bounds on $a_{\rm win},s_{\rm win}$
follow by substituting $q_*=1/4096$.

For an arbitrary stored state, $\|\bar c_M\|_2\le\sqrt{1+q_*}$, and hence
\[
\sigma_a\le\lambda_{\max}(1+q_*)^{3/2},\qquad
\sigma_a\le\lambda_{\max}(1+q_*)\max_j|\bar c_{M,a,j}|.
\tag{27}
\]
The first inequality, (16), and the score filter give the first two bounds in
(18).  The second inequality in (27), the same filter, and
the weight ratio $\lambda_{\max}/\lambda_{\min}\le1.01$ give
$\max_j|\bar c_{M,a,j}|>3/4$ for every mode. The $d$-bound and
$\lambda_{\max}/d<7/5$ follow from (6) and the numerical slack
$\tau_r<10^{-6}$.

Apply (8) with a selector and divide by $d_{M,a}$. The bound
$\|\bar k_M\|_1\le\lambda_{\max}(1+q_*)$ and (18) give
$\eta_{\rm raw}$. By (6),
$|1-d_{M,a}/\sigma_a|\le\tau_r/(1-\tau_r)$; the additional conversion
from $d^{-1}$ to $\sigma^{-1}$ is therefore at most the last term in (19),
proving (20). Taking a singleton selector gives (22). Since
$\lambda_j/\sigma_a\le\lambda_{\max}/\sigma_a<7/5$ on the retained pool,
(22) gives
$|\bar c_{M',a,j_M}\bar c_{M'',a,j_M}|>(3/4-1/2000)/(7/5)>1/3$.
Each held correlation is therefore larger than $1/3>1/8$, proving the
threshold intersection claim and nonemptiness.

For (23), on $I_{M,a}(t)$ use $|\bar c_{M,a,j}|<t$ and Cauchy--Schwarz:
\[
\sum_{j\in I_{M,a}(t)}\lambda_j|\bar c_{U,j}\bar c_{V,j}\bar c_{W,j}|
\le\lambda_{\max}t\|\bar c_{M'}\|_2\|\bar c_{M''}\|_2
\le\lambda_{\max}(1+q_*)t.
\]
Equation (24) is the same calculation with the witness ratios and
$\sum\rho_U\rho_V\rho_W\le u_{\rm win}$. All these inequalities are
slotwise finite deterministic statements on the certified-state conclusion of Proposition~\ref{prop:step-004-certificate} and the event from Proposition~\ref{prop:step-001-geometry}; no label is
read by the procedure. 

\end{proof}
\end{proposition}
\begin{proof}
Fix the event from Proposition~\ref{prop:step-001-geometry} and the certified-state conclusion of Proposition~\ref{prop:step-004-certificate}.  For each certified slot,
Lemma~\ref{lem:step-005-signed-equation} starts from the exact old-state
Jacobi contraction and multiplies the observable projective residual by the
same contraction norm.  It proves the raw right-sided equation, shows that
the three certificate signs are one common score sign, and applies a
score-preserving proof-only orientation to obtain the barred equation.  The
orientation preserves the represented rank-one term, so the tensor remains
the one represented by $(G_U,G_V,G_W,D_\lambda)$.

Proposition~\ref{prop:step-005-projected-ledger} then projects that same
equation with the realized target Gram.  Its selector inequality separates
the Gram leakage, the two first-order coordinate leaks, their quadratic cross
term, and the certificate residual. The residual is charged once as
$\sqrt{r(1+q)}\,\tau_r$, rather than being multiplied by a selection or ALS
horizon. The exact $d$-normalized identity records the scalar cancellation;
conversion to the score $\sigma_a$ adds only the
explicit one-time term in (19).

Finally, Proposition~\ref{prop:step-005-threshold-ledger} uses only the
already generated coverage witnesses to lower-bound $\sigma_{\max}$. It
then instantiates the selector inequality for every observable high-score
slot and every proof-only (1/8)-threshold/tail set, while retaining the unnormalized
ledger for low-score slots.  The resulting finite conjunction is precisely
the score-equation conclusion of Proposition~\ref{prop:step-005-threshold-ledger} and the tail conclusion of Proposition~\ref{prop:step-005-threshold-ledger}.  It is produced before the score-filter output
is used to form clusters, and all target indices in the witness argument are
proof-only choices rather than algorithmic inputs.  Thus the exact displayed
claim is proved with no landing, selected representative, or cyclic
condition.
\end{proof}

\subsection{Observable selection and gauge}\label{app:step-006}

Proposition~\ref{prop:step-001-realized-gram} supplies the $H_M$ row/column
bound and its spectral consequence. Lemma~\ref{lem:step-005-signed-equation}
and Proposition~\ref{prop:step-005-threshold-ledger} supply the barred
correlations $c_{M,a}=G_M^\top\bar p_{M,a}$, the positive score
$\sigma_a=|\theta_a|$, and, for every selector $I\subseteq[r]$,
\[
\sum_{j\in I}\left|c_{M,a,j}
-\frac{\lambda_j}{\sigma_a}c_{M',a,j}c_{M'',a,j}\right|
\le\eta_{\rm score}(r).
\tag{A}
\]
On the retained pool ${\cal H}=\{a:\sigma_a\ge0.85\sigma_{\max}\}$, the
same proposition gives the retained score floor, the dominant-coordinate
$3/4$ bound, and the coverage-witness score bounds. We use only these named
interfaces and the elementary Cauchy--Schwarz and unit-vector identities.

Fix the event from Proposition~\ref{prop:step-001-geometry} and the finite certified pool. For a barred slot
\(a\), write \(c_M=c_{M,a}\), \(\sigma=\sigma_a\), and set
\[
 \eta:=\eta_{\rm score}(r)<{1\over2000},\qquad
 \rho_a:={\lambda_{\max}\over\sigma_a}<{7\over5}. \tag{1}
\]
All indices introduced in this derivation are latent proof indices.


\begin{proposition}[Common dominant support and small-root tail]\label{prop:step-006-support}
Under the static event of Proposition~\ref{prop:step-001-geometry} and the
score/threshold conclusions of Proposition~\ref{prop:step-005-threshold-ledger}, for
every \(a\in\mathcal H\) there is a unique index \(\pi(a)\) such that
\[
 |c_{U,a,\pi(a)}|,\ |c_{V,a,\pi(a)}|,\ |c_{W,a,\pi(a)}|>{3\over4}. \tag{2}
\]
Moreover, with
\[
 x_a:=\max_{M\in\{U,V,W\},\,\ell\ne\pi(a)}|c_{M,a,\ell}|,\qquad
 L_{M,a}:=\sum_{\ell\ne\pi(a)}|c_{M,a,\ell}|, \tag{3}
\]
one has
\[
 x_a<2\eta,\qquad
 L_{M,a}\le {\eta\over1-(14/5)\eta},\qquad
 \sum_{\ell\ne\pi(a)}c_{M,a,\ell}^2
 \le \chi:= {2\eta^2\over1-(14/5)\eta}. \tag{4}
\]

\begin{proof}
Proposition~\ref{prop:step-005-threshold-ledger} gives, for each mode, a
coordinate of magnitude \(>3/4\). There cannot be two such coordinates,
because \(\|c_M\|_2^2\le1+q_*<18/16\). Denote the unique indices by
\(j_U,j_V,j_W\).

We first rule out \(j_U\ne j_V\). Put \(S:=1+q_*\) and \(A:=9/16\),
\(u_i:=|c_{U,i}|^2\), \(v_i:=|c_{V,i}|^2\). If \(a=j_U\) and
\(b=j_V\ne a\), then \(u_a>A\), \(v_b>A\), and
\(\sum_i u_i,\sum_i v_i\le S\). With \(t=u_a\) and \(z=v_a\),
\[
\begin{aligned}
 \sum_i u_i v_i
 &\le tz+(S-t)\max_{i\ne a}v_i\\
 &\le tz+(S-t)(S-z).
\end{aligned} \tag{5}
\]
Here \(t\in[A,S]\) and \(z\in[0,S-A]\), since \(v_b>A\). The last
expression is bilinear in \(t,z\); checking its four rectangle corners gives
\[
 \sum_i u_i v_i\le2A(S-A)
 ={9\over8}\left({7\over16}+q_*\right)<{1\over2}. \tag{6}
\]
Consequently, by the weighted triangle inequality and Cauchy--Schwarz,
\[
 \sigma
 \le\lambda_{\max}\sum_i|c_{U,i}c_{V,i}c_{W,i}|
 \le\lambda_{\max}\|c_U\circ c_V\|_2\|c_W\|_2
 <{71\over100}\lambda_{\max}. \tag{7}
\]
On the other hand, the coverage score floor and the observable filter imply
for every retained slot
\[
 \sigma\ge0.85\sigma_{\max}>0.765\lambda_{\min}
 \ge {0.765\over1.01}\lambda_{\max}>{3\over4}\lambda_{\max}, \tag{8}
\]
contradicting (7). The same argument for the pairs \(U,W\) and \(V,W\)
proves \(j_U=j_V=j_W=:\pi(a)\). This is the promised
singleton-versus-multi-support separation: a mismatched multi-support
pattern has score \(<0.71\lambda_{\max}\), whereas the retained singleton
floor is \(>0.75\lambda_{\max}\).

For \(\ell\ne\pi(a)\), taking \(I=\{\ell\}\) in (A) gives
\[
 |c_{M,\ell}|\le {\lambda_\ell\over\sigma}
 |c_{M',\ell}c_{M'',\ell}|+\eta
 \le {7\over5}x_a^2+\eta. \tag{9}
\]
The dominant coordinate and norm upper bound also give
\[
 x_a^2\le S-{9\over16}<\left({2\over3}\right)^2. \tag{10}
\]
If \(\eta=0\), (9) immediately gives \(x_a=0\). Otherwise, on
\([2\eta,2/3]\) the concave function \(f(x)=x-(7/5)x^2\) is bounded below
by its endpoint values:
\[
 f(2\eta)>\eta,\qquad f(2/3)={2\over45}>\eta.
\]
Thus (9)--(10) force \(x_a<2\eta\). Summing (A) over
\(I=[r]\setminus\{\pi(a)\}\), and using one factor bounded by \(x_a\), gives
\[
 L_{M,a}\le {7\over5}x_a L_{M'',a}+\eta. \tag{11}
\]
Taking the maximum over modes and using \(x_a<2\eta\) proves the second
bound in (4). Finally,
\(\sum_{\ell\ne\pi(a)}c_{M,a,\ell}^2\le x_aL_{M,a}\) gives the last
bound. 


\end{proof}
\end{proposition}
\begin{proposition}[Observable graph components]\label{prop:step-006-graph}
Under the hypotheses and conclusion of
Proposition~\ref{prop:step-006-support}, retain exactly the certified slots
\[
{\cal H}:=\{a:\sigma_a\ge0.85\sigma_{\max}\},
\]
and join two retained slots $a,b$ exactly when
\[
|\langle p_{M,a},p_{M,b}\rangle|\ge1-64q_*
\qquad\text{for every }M\in\{U,V,W\}.
\]
Then every target $j$ has at least one retained slot, and this observable
graph has exactly the $r$ nonempty components
\[
 \mathcal H_j:=\{a\in\mathcal H:\pi(a)=j\},\qquad j\in[r]. \tag{12}
\]
Consequently the exactly-$r$ component test passes. Choosing in each component
the minimum-$\zeta$ member, with score as the tie-break, returns one
representative per target without using target labels.

\begin{proof}
Let \(a(j)\) be the coverage witness supplied by the tail conclusion of Proposition~\ref{prop:step-005-threshold-ledger} for
target \(j\). Its score obeys
\(\sigma_{a(j)}\ge s_{\rm win}\lambda_j>(9/10)\lambda_{\min}\).
For every certified slot, including the maximizer,
\(\sigma_i\le\lambda_{\max}(1+q_*)^{3/2}\). Hence
\[
 {\sigma_{a(j)}\over\sigma_{\max}}
 \ge {s_{\rm win}\lambda_{\min}\over
 \lambda_{\max}(1+q_*)^{3/2}}
 >{0.9\over1.01(1+1/4096)^{3/2}}>0.85, \tag{13}
\]
so \(a(j)\in\mathcal H\) and \(\mathcal H_j\ne\varnothing\).

Set
\[
 \chi_*:={2\eta_{\rm score}(r)^2\over
 1-(14/5)\eta_{\rm score}(r)},\qquad
 d_*:=\sqrt{2(q+\chi_*)}. \tag{14}
\]
For \(a\in\mathcal H_j\), Proposition~\ref{prop:step-006-support} and
the lower norm bound in the tail conclusion of Proposition~\ref{prop:step-005-threshold-ledger} imply
\[
 |c_{M,a,j}|^2\ge1-q-\chi_*,\qquad
 \|\widehat p_{M,a}-g_{M,j}\|_2^2
 =2(1-|c_{M,a,j}|)\le2(q+\chi_*)=d_*^2, \tag{15}
\]
where \(\widehat p_{M,a}:=\operatorname{sgn}(c_{M,a,j})\bar p_{M,a}\) and
\(g_{U,j}=u_j,g_{V,j}=v_j,g_{W,j}=w_j\). The sign is well-defined by
(15).

Since \(\eta<1/2000\), direct substitution gives
\[
 q+\chi_*<{1\over4000},\qquad d_*^2<{1\over2000},\qquad d_*<{1\over40}. \tag{16}
\]
If \(a,b\in\mathcal H_j\), then in every mode
\[
 \begin{aligned}
 |\langle p_{M,a},p_{M,b}\rangle|
 &=|\langle\widehat p_{M,a},\widehat p_{M,b}\rangle|\\
 &\ge1-\tfrac12(\|\widehat p_{M,a}-g_{M,j}\|_2
             +\|\widehat p_{M,b}-g_{M,j}\|_2)^2\\
 &\ge1-2d_*^2>1-{1\over1000}>1-64q_*.
 \end{aligned} \tag{17}
\]
Thus all vertices in \(\mathcal H_j\) are joined. If \(a\in\mathcal H_j\),
\(b\in\mathcal H_\ell\), and \(j\ne\ell\), then
\(|\langle g_{M,j},g_{M,\ell}\rangle|\le q\) and (15)--(16) give
\[
 |\langle p_{M,a},p_{M,b}\rangle|
 \le q+2d_*+d_*^2<{1\over8}<1-64q_*.
 \tag{18}
\]
Therefore no cross-target edge exists. The graph is the disjoint union of
the \(r\) nonempty cliques in (12), so its connected components are exactly
those cliques. The component-count test therefore passes, and the
minimum-$\zeta$ rule with score tie-breaking selects exactly one observable
member from each component.


\end{proof}
\end{proposition}
\begin{lemma}[Product-preserving target permutation and sign gauge]\label{lem:step-006-gauge}
Under Proposition~\ref{prop:step-006-graph},
the observable representative rule yields a bijection between the \(r\)
representatives and the realized target columns. There is a proof-only
permutation and a sign gauge that preserves every represented rank-one term.

\begin{proof}
The graph components are the nonempty sets \(\mathcal H_j\). The
minimum-\(\zeta\) (then score) rule selects one member \(a(j)\) from each,
and the arbitrary algorithmic ordering induces a bijection
\(\pi:[r]\to[r]\) with representative \(a\) assigned to target
\(\pi(a)\). This is a proof-only permutation; the algorithm never reads
\(\pi\).

For a representative \(a\), use the score-preserving barred orientation from
Proposition~\ref{prop:step-005-threshold-ledger}, so that
\[
 \sigma_a\,\bar p_{U,a}\otimes\bar p_{V,a}\otimes\bar p_{W,a}
 =\theta_a\,p_{U,a}\otimes p_{V,a}\otimes p_{W,a}. \tag{19}
\]
Let \(\xi_{M,a}=\operatorname{sgn}\langle g_{M,\pi(a)},\bar p_{M,a}\rangle\),
which is defined by (15), and let \(\chi_a=\xi_{U,a}\xi_{V,a}\xi_{W,a}\).
The target-aligned directions and signed scalar
\[
 \widetilde p_{M,a}=\xi_{M,a}\bar p_{M,a},\qquad
 \widetilde\theta_a=\chi_a\sigma_a \tag{20}
\]
satisfy
\[
 \widetilde\theta_a\bigotimes_M\widetilde p_{M,a}
 =\sigma_a\bigotimes_M\bar p_{M,a}, \tag{21}
\]
because \(\chi_a^2=1\). Thus (20) is a proof-only sign chart that aligns
all target coordinates and preserves the represented tensor. For the literal
product-one sign convention used by the quotient chart, set
\[
 (\varepsilon_{U,a},\varepsilon_{V,a},\varepsilon_{W,a})
 :=(\xi_{U,a},\xi_{V,a},\chi_a\xi_{W,a}).
\]
Then \(\varepsilon_{U,a}\varepsilon_{V,a}\varepsilon_{W,a}=1\), so
\(\bigotimes_M\varepsilon_{M,a}\bar p_{M,a}
=\bigotimes_M\bar p_{M,a}\); the coefficient \(\sigma_a\) is unchanged.
The all-positive chart (20) and this product-one chart differ only by
placing the sign \(\chi_a\) in the W scalar, exactly as in the observable
seed definition. Positive componentwise rescalings with product one are
unchanged by this sign bookkeeping. Hence the resulting permutation/gauge is
exactly the interface exported as the gauge conclusion of Lemma~\ref{lem:step-006-gauge}. 

\end{proof}
\end{lemma}
\begin{proof}
Fix the event from Proposition~\ref{prop:step-001-geometry} and the
interfaces supplied by the preceding propositions. The weighted ledger
supplies the retained score floor, the \(3/4\) dominant coordinate in each
mode, and (A). Proposition~\ref{prop:step-006-support} first compares two
different dominant indices through (5)--(8), proving that every retained
slot has one common target support. The same proposition then applies the
small-root branch of (A), giving a dimension-free off-target \(\ell_1\) tail
and direction error.

Coverage witnesses are retained by (13). Proposition~\ref{prop:step-006-graph}
uses the direction error and the realized Gram off-diagonal bound to prove
within-target edges and cross-target nonedges at the exact observable
threshold \(1-64q_*\). It follows that the algorithm's graph has exactly
\(r\) components and that the selected representative in each component is
associated with a distinct realized target. Finally,
Lemma~\ref{lem:step-006-gauge} supplies the proof-only permutation and
product-preserving sign chart. Together these results produce
\[
 E_{\rm support},\qquad E_{\rm cluster},\qquad E_{\rm gauge},
\]
before any selected-state coefficient, landing output, or cyclic invariant is
used. The construction is label-free at the procedure level and preserves
the represented rank-one terms exactly.
\end{proof}

\subsection{Two-orientation coefficient closure}\label{app:step-007}

The geometry bounds used below are supplied by
Proposition~\ref{prop:step-001-geometry}. The signed equation and the support,
graph, and gauge conclusions are supplied by
Lemma~\ref{lem:step-005-signed-equation}, Proposition~\ref{prop:step-006-graph},
and Lemma~\ref{lem:step-006-gauge}; the latter contributes only the
proof-only permutation and term-preserving product-one option. No landing
state or later ALS invariant is imported. The remaining tools are
Cauchy--Schwarz, the Hadamard-product norm inequality, and induced row/column
$\ell_1$ estimates, all derived below.

Reindex by the proof-only target permutation and put
\[
q:=q_{\rm real},\qquad t:=r\tau_r,\qquad
\omega:=q_*^2+t,\qquad \eta:=\eta_{\rm score}(r).
\]
The definitions of $\eta_{\rm score}$ and
$\tau_r=q_*^2/(10000r)$, together with $q\le q_*$, imply
\[
\eta\le\frac32q_*,\qquad \ell_*\le\frac{31}{20}q_*,\qquad
x_*:=2\eta\le3q_*,\qquad x_*\ell_*\le\frac{93}{20}q_*^2.
\tag{A}
\]
Indeed,
\[
\sqrt{r(1+q_*)}\,\tau_r\le\frac{q_*}{10^6},\qquad
\frac{1.4(1+q_*)\tau_r}{1-\tau_r}\le\frac{q_*}{10^6},
\]
and hence
$\eta\le1.4(1+q_*)q_*+2q_*/10^6<3q_*/2$. Since
$\eta<1/2000$, division by $1-(14/5)\eta$ gives the stated $\ell_*$ bound.
All constants in (A) are universal and remain valid when $q_{\rm real}=0$.


\begin{lemma}[sign-covariant dual equation, positive diagonal, and projected residual]\label{lem:step-007-dual-equation}

Under the static event of Proposition~\ref{prop:step-001-geometry}, the
certificate/score ledger of Proposition~\ref{prop:step-005-threshold-ledger},
and the support and graph conclusions of Propositions~\ref{prop:step-006-support}
and \ref{prop:step-006-graph}, first keep the score-oriented objects barred.
After the proof-only target permutation, write
\[
\bar c_{M,ja}=(G_M^\top\bar P_M)(j,a),\qquad
\bar\sigma_a=|\theta_a|
 =\sum_j\lambda_j\bar c_{U,ja}\bar c_{V,ja}\bar c_{W,ja}
 >\frac34\lambda_{\max}.
\tag{A.1}
\]
For each $a$ let $\xi_{M,a}=\operatorname{sign}(\bar c_{M,aa})$; then
$\xi_{U,a}\xi_{V,a}\xi_{W,a}=1$.  Define
\[
 p_{M,a}=\xi_{M,a}\bar p_{M,a},\quad
 c_{M,ja}=\xi_{M,a}\bar c_{M,ja},\quad
 k_{M,a}=\xi_{M',a}\xi_{M'',a}\bar k_{M,a}=\xi_{M,a}\bar k_{M,a},\quad
 r_{M,a}=\xi_{M,a}\bar r_{M,a}.
\tag{A.2}
\]
Stacking these columns gives $P_M,K_M,R_M$ and the exact identities
\[
 P_MD_M=G_MK_M+R_M,\qquad c_M=G_M^\top P_M,\qquad
 A_M=G_M^\dagger P_M=H_M^{-1}G_M^\top P_M,\qquad c_M=H_MA_M.
\tag{B0}
\]
Set $\sigma_a=\bar\sigma_a$ and relabel without bars.  With
$a_{M,a}=A_M(a,a)$,
\[
 0.74<a_{M,a}<1.001,\qquad
 \Delta_M=\operatorname{diag}(a_{M,a}),\qquad
 Z_M=A_M\Delta_M^{-1},\qquad
 h_{M,ja}=(H_MZ_M)(j,a)=\frac{c_{M,ja}}{a_{M,a}}.
\tag{A.3}
\]
Then $h_{M,aa}>2/3$.  If
\[
 \varepsilon_{M,ja}:=\frac{(G_M^\dagger R_M)(j,a)}{d_{M,a}},
\]
then $|\varepsilon_{M,ja}|\le2\tau_r$, and for $j\ne a$,
\[
 |Z_M(j,a)|\le4|h_{M',ja}h_{M'',ja}|+3\tau_r.
\tag{B}
\]

\begin{proof}

The barred score identity and retained-pool bounds are the pre-gauge facts.
If $\xi_{U,a}\xi_{V,a}\xi_{W,a}<0$, the dominant term would be negative and
the off-target bounds would give
\[
\bar\sigma_a\le\lambda_{\max}x_*
 \|\bar c_{M',-a}\|_2\|\bar c_{M'',-a}\|_2
 \le\lambda_{\max}x_*(1+q_*)<\frac34\lambda_{\max}.
\tag{A.4}
\]
This contradicts the retained score floor, so the signs have product one.
Multiplying the barred same-state equation in mode $M$ by $\xi_{M,a}$ gives
(B0), because $\xi_{M',a}\xi_{M'',a}=\xi_{M,a}$. The residual column obeys
\[
\|R_M(:,a)\|_2\le d_{M,a}\tau_r.
\tag{A.5}
\]
The product of the transformed correlations equals the barred product, so the
score identity and every absolute tail in (A) are unchanged. This is a
product-one sign gauge; by construction $c_{M,aa}=|\bar c_{M,aa}|>3/4$.

Since $P_M\in{\rm range}(G_M)$, write $P_M=G_MA_M$. The unit-column norm
and eigenvalue floor give $\|A_M(:,a)\|_2\le(1-q)^{-1/2}$, and hence
\[
 |c_{M,aa}-A_M(a,a)|
 \le\|H_M(a,-a)\|_2\|A_M(-a,a)\|_2
 \le {q\over\sqrt{1-q}}.
\tag{C}
\]
Thus $0.74<a_{M,a}<1.001$ and $h_{M,aa}>2/3$. Applying $G_M^\dagger$ to
(B0) gives
\[
d_{M,a}A_M(:,a)=K_M(:,a)+(G_M^\dagger R_M)(:,a).
\tag{D}
\]
The pseudoinverse bound gives
\[
\|\varepsilon_{M,:,a}\|_2\le{\tau_r\over\sqrt{1-q}}<2\tau_r,
\tag{E}
\]
and the directly projected residual separately satisfies
\[
{ |(G_M^\top R_M)(j,a)|\over d_{M,a}}
 \le\sqrt{1+q}\,\tau_r<2\tau_r.
\tag{E'}
\]
Also
\[
d_{M,a}\le\|G_M\|_2\|K_M(:,a)\|_2
\le\lambda_{\max}(1+q).
\tag{A.6}
\]
For $j\ne a$, division of the $j$-coordinate of (D) by the $a$-coordinate
has denominator at least
\[
{9\over16}-{1.02\tau_r\over\sqrt{1-q}}>0.56.
\tag{F}
\]
The product part of the numerator divided by $\lambda_a$ is at most
$1.02|h_{M',ja}h_{M'',ja}|$, and the residual part is at most
$1.02\tau_r$. Division by (F) gives (B), with the residual bounded entrywise
before any sum. QED.



\end{proof}
\end{lemma}
\begin{proposition}[normalized dual-coordinate column bootstrap]\label{prop:step-007-column}

Under Lemma~\ref{lem:step-007-dual-equation},

\[
\max_M\|Z_M-I\|_{\rm col,1}\le42\omega.
\tag{G}
\]

\begin{proof}

For each selected column $a$, (A), (C), and $a_{M,a}>0.74$ give
\[
\|h_{M,-a}\|_2\le{\|c_{M,-a}\|_2\over a_{M,a}}
\le{3\over2}\sqrt{x_*\ell_*}.
\tag{H}
\]

Summing (B) over $j\ne a$ and using Cauchy--Schwarz gives
\[
\sum_{j\ne a}|Z_M(j,a)|
\le4\|h_{M',-a}\|_2\|h_{M'',-a}\|_2+3(r-1)\tau_r
\le9x_*\ell_*+3t\le42(q_*^2+t)=42\omega.
\tag{I}
\]

Taking the maximum over M and a proves (G).  This is a genuine column
bootstrap: no row mass, landing output, or cyclic condition occurs.
\end{proof}


\end{proposition}
\begin{proposition}[two-orientation coefficient closure]\label{prop:step-007-row}

Under Lemma~\ref{lem:step-007-dual-equation} and Proposition
\ref{prop:step-007-column}, let

\[
R:=\max_M\|Z_M-I\|_{\rm row,1},\qquad
C:=\max_M\|Z_M-I\|_{\rm col,1}.
\]

Then

\[
R\le5\omega,\qquad
\max_M\|A_M-I\|_{\rm row,1}\le7\omega,\qquad
\max_M\|A_M-I\|_{\rm col,1}\le44\omega.
\tag{J}
\]

The normalized residual entries obey
$|\varepsilon_{M,ja}|\le2\tau_r$, and each row or column sum is charged once
by at most $2r\tau_r$.

\begin{proof}

Fix a mode $M$ and row $j$. For $a\ne j$, set
$v_{M,ja}:=(H_MZ_M)(j,a)$. From (G),

\[
\max_{a\ne j}|v_{M,ja}|\le q+(1+q)C=:u.
\tag{K}
\]

Here the direct $H_M(j,a)$ term has row mass $q$, the off-diagonal $Z_M$
entry is at most $C$, and the remaining Gram row contributes at most $qC$.

\[
\sum_{a\ne j}|v_{M,ja}|
\le q+(1+q)\|Z_M-I\|_{\rm row,1}.
\tag{L}
\]

To verify (L), the $l=j$ term in $H_MZ_M$ contributes the row error. Each
$l\ne j$ term contributes its off-diagonal row mass and its diagonal entry
$Z_M(l,l)=1$, weighted by the $H_M$ row mass $q$.

Summing (B) over $a\ne j$, using (K) for one factor and (L) for the other,
and taking the maximum over modes gives

\[
R\le4u\{q+(1+q)R\}+3t.
\tag{M}
\]

By (G), $t=q_*^2/10000$ and $q\le q_*$, so

\[
u\le q+42(1+q_*)(1+10^{-4})q_*^2<{51\over50}q_*,\qquad
4u(1+q)<{1\over1000}.
\tag{N}
\]

Thus (M) has a positive denominator and

\[
R\le{4uq+3t\over1-4u(1+q)}<5\omega.
\tag{O}
\]

This is the small-root selection.  Before inserting (G), (M) is the usual
quadratic/small-gain envelope; (G) puts the solution in the small branch,
and (N) excludes the large branch without assuming a row tube.

It remains to transfer from normalized columns to $A_M$. Let $z$ be a column
of $Z_M$, with $z_a=1$ and $\|z-e_a\|_1\le C$, and write $a_0=a_{M,a}$.
The unit-norm identity gives
\[
|z^\top H_Mz-1|\le2qC+(1+q)C^2=:\delta_a.
\tag{P}
\]

At $q\le q_*$ and $C\le42\omega$, $\delta_a<1/100$ and
$|a_0-1|\le2\delta_a\le\omega$; the last inequality follows directly from
$\omega=(1+10^{-4})q_*^2$ and $q_*=1/4096$.
Since $a_0<1.001$, (G), (O), and this diagonal bound imply

\[
\|A_M-I\|_{\rm col,1}\le1.001C+\omega<44\omega,\qquad
\|A_M-I\|_{\rm row,1}\le1.001R+\omega<7\omega.
\]

This proves (J). Equation (E) gives the entrywise $2\tau_r$ residual, and
summing over at most $r$ entries gives $2r\tau_r$ in either orientation,
never $r^2\tau_r$.\quad\Box

\end{proof}
\end{proposition}
\begin{proof}
The support and gauge interfaces first provide one selected
stored state per target, a proof-only permutation, and a term-preserving
product-one sign option.  Lemma~\ref{lem:step-007-dual-equation} starts from the
barred score identity, proves the dominant sign product, explicitly
transforms the equation and residual, and only then applies the true-factor
dual.  It proves the positive diagonal and bounds every projected residual
entry before summation.  Proposition~\ref{prop:step-007-column} uses
only the per-column correlation tails to produce the $42\omega$
column interface.  Proposition~\ref{prop:step-007-row} then uses that interface
to bound held correlation rows, solve the positive-denominator small-gain
inequality, and transfer from normalized dual columns to $A_M-I$. Therefore
the column-closure conclusion of Proposition~\ref{prop:step-007-column}, the projected-residual conclusion of Proposition~\ref{prop:step-007-row}, and the row-closure conclusion of Proposition~\ref{prop:step-007-row} hold with both induced orientations at
$O(\omega)$, with one residual charge $r\tau_r$. All objects are from the
selected stored state; no landing or cyclic object is consumed.
Because the stored-state interface places every selected column in
$\operatorname{range}(G_M)$, the orthogonal complement field is exactly zero
at this stage and therefore satisfies the later budget $\tau_r$.
\end{proof}

\subsection{Balanced observable seed}\label{app:step-008}

Lemma~\ref{lem:step-006-gauge} supplies one selected representative for each
target, a proof-only permutation, and a product-one sign gauge. In the
resulting orientation, Proposition~\ref{prop:step-006-support} gives
\[
\max_{M,\ell\ne j}|c_{M,\ell j}|\le
x_*:=2\eta_{\rm score}(r),\qquad
\sum_{\ell\ne j}|c_{M,\ell j}|\le
\ell_*:=\frac{\eta_{\rm score}(r)}
{1-(14/5)\eta_{\rm score}(r)},
\tag{A}
\]
with $x_*\le3q_*$ and $\ell_*\le(31/20)q_*$. The score identity is
\[
\sigma_j:=|\theta_j|
=\sum_{\ell=1}^r\lambda_\ell
c_{U,\ell j}c_{V,\ell j}c_{W,\ell j}>0.
\tag{B}
\]
Proposition~\ref{prop:step-007-row} gives the exact target-span coefficient
interface
\[
P_M=G_MA_M,\qquad
\|A_M-I\|_{\rm row,1}\le7\omega,\qquad
\|A_M-I\|_{\rm col,1}\le44\omega,
\tag{C}
\]
and the columns of $P_M$ are unit vectors. The remaining tools are orthogonal
projection, induced $\ell_1$ inequalities, Cauchy--Schwarz, and
$|\log(1+u)|\le |u|/(1-|u|)$ for $|u|<1$.

Put
\[
 q:=q_{\rm real},\qquad t:=r\tau_r,\qquad
 \omega:=q_*^2+t,\qquad \beta:=44\omega. \tag{1}
\]
The setting gives $t=q_*^2/10^4$, hence
$\omega=(1+10^{-4})q_*^2<10^{-6}$ and $\beta<10^{-4}$. All numerical
inequalities below are uniform in $r,n,\kappa_0,\rho$ and the confidence
parameters.


\begin{lemma}[observable seed decomposition]\label{lem:step-008-seed-decomp}
Under the static event of Proposition~\ref{prop:step-001-geometry}, the
permutation/gauge conclusion of Lemma~\ref{lem:step-006-gauge}, and the
two-orientation closure of Proposition~\ref{prop:step-007-row}, after the
proof-only permutation and product-one sign chart, let $P_M^0$ be the
normalized active seed factors. Then, for every mode $M$,
\[
 P_M^0=G_M(I+C_M^0)+N_M^0,\qquad G_M^\top N_M^0=0, \tag{2}
\]
with
\[
 N_M^0=0,\qquad
 \|C_M^0\|_{\rm row,1}\vee\|C_M^0\|_{\rm col,1}\le48\omega. \tag{3}
\]

\begin{proof}

For selected representative $j$, the setting forms
\[
 x_j^0=|\theta_j|^{1/3}p_{U,j},\quad
 y_j^0=|\theta_j|^{1/3}p_{V,j},\quad
 z_j^0=\operatorname{sgn}(\theta_j)|\theta_j|^{1/3}p_{W,j}. \tag{4}
\]
The score-preserving orientation from the gauge conclusion of Lemma~\ref{lem:step-006-gauge} has the same rank-one product
as (4) with scalar $\sigma_j=|\theta_j|$. Hence the two triples of unit
directions differ only by a product-one sign gauge. Applying that proof-only
gauge and the target permutation gives
\[
 P_M^0=P_M, \tag{5}
\]
where $P_M$ is the oriented selected matrix in Proposition~\ref{prop:step-007-row}.
This is an identity of the actual observable seed, not a new
algorithmic state; represented rank-one terms are unchanged.

By Proposition~\ref{prop:step-001-geometry}, $\lambda_{\min}(H_M)\ge1-q>0$,
so $G_M$ has full column rank. Proposition~\ref{prop:step-007-row} gives
$P_M=G_MA_M$ with (C). The orthogonal
decomposition relative to the exact target span is unique:
\[
 A_M:=G_M^\dagger P_M,\qquad
 N_M^0:=(I-G_MG_M^\dagger)P_M. \tag{6}
\]
Because $P_M\in{\rm range}(G_M)$, (6) gives $N_M^0=0$ and
$P_M^0=G_M(I+C_M^0)$ with $C_M^0=A_M-I$. Thus (C) implies
\[
 \|C_M^0\|_{\rm row,1}\le7\omega\le48\omega,\qquad
 \|C_M^0\|_{\rm col,1}\le44\omega\le48\omega.
\]
The requested perpendicular-column bound follows from the stronger identity
$\max_j\|N_M^0(:,j)\|_2=0\le\tau_r$. QED.


\end{proof}
\end{lemma}
\begin{lemma}[best-scalar transfer]\label{lem:step-008-best-scalar}
Under the static event of Proposition~\ref{prop:step-001-geometry}, the
support conclusion of Proposition~\ref{prop:step-006-support}, and the
closure conclusion of Proposition~\ref{prop:step-007-row}, with the oriented matrices in
Lemma~\ref{lem:step-008-seed-decomp}, for every target column $j$,
\[
 |c_{M,jj}-1|\le2\omega\quad(M\in\{U,V,W\}),\qquad
 \left|\frac{\sigma_j}{\lambda_j}-1\right|\le9\omega. \tag{7}
\]
Here $c_M=G_M^\top P_M^0$ and $\sigma_j=|\theta_j|>0$ is the observable
best-scalar in (B).

\begin{proof}

Fix $M,j$, write $H=H_M=I+E$, $A=A_M$, and let
\[
 a:=A_{jj},\qquad \alpha:=A(:,j)-e_j.
\]
From (C), $\|\alpha\|_1\le\beta$, so $a>0$ and
$|a-1|\le\beta<1/100$. Put $z:=A(:,j)/a=e_j+u$. Then $u_j=0$ and
\[
 B:=\|u\|_1\le\frac{\beta}{1-\beta}\le45\omega. \tag{8}
\]
The column $P_M^0(:,j)=G_M(az)$ has norm one. Since the diagonal of $H$
is one and the diagonal of $E$ is zero,
\[
 1=a^2 z^\top H z,\qquad
 |z^\top H z-1|\le2qB+(1+q)B^2=:d. \tag{9}
\]
Indeed, $\|u\|_2^2\le B^2$, $|e_j^\top Eu|\le q\|u\|_\infty\le qB$,
and $|u^\top Eu|\le q\|u\|_1\|u\|_\infty\le qB^2$. Substituting
$B\le45\omega$, $q\le q_*$, and $\omega<10^{-6}$ gives
\[
 d\le90q_*\omega+(1+q_*)2025\omega^2<\frac{\omega}{20}<\frac1{20}. \tag{10}
\]
The elementary inverse-square-root bound
$|(1+v)^{-1/2}-1|\le2|v|$ for $|v|\le1/2$, applied to (9), yields
$|a-1|\le2d<\omega/10$, and in particular $|a-1|\le\omega$. Now
$c_M=HA$ and $E_{jj}=0$, so
\[
 c_{M,jj}-1=(a-1)+e_j^\top E\alpha.
\]
Consequently, using $\|\alpha\|_\infty\le\beta$,
\[
 |c_{M,jj}-1|\le\omega+q\beta
 \le(1+44q_*)\omega<2\omega. \tag{11}
\]

It remains to transfer this diagonal control to the scalar. By the support ledger (A), every off-target column satisfies the same $x_*,\ell_*$
bounds after the sign gauge. Write $c_{M,jj}=1+\delta_M$, so
$|\delta_M|\le2\omega$. Then
\[
 \left|\prod_M c_{M,jj}-1\right|
 \le3(2\omega)+3(2\omega)^2+(2\omega)^3
 <7\omega. \tag{12}
\]
Using (B), $\lambda_\ell/\lambda_j\le\Gamma$, and two off-target maxima
together with one off-target $\ell_1$ sum, we obtain
\[
\begin{aligned}
 \left|\frac{\sigma_j}{\lambda_j}-1\right|
 &\le 7\omega+
 \Gamma\sum_{\ell\ne j}
 |c_{U,\ell j}c_{V,\ell j}c_{W,\ell j}|\\
 &\le 7\omega+\Gamma x_*^2\ell_*.
\end{aligned} \tag{13}
\]
The endpoint bounds $x_*\le3q_*$, $\ell_*\le(31/20)q_*$,
$\Gamma\le1.01$, and $\omega\ge q_*^2$ give
\[
 \Gamma x_*^2\ell_*
 \le1.01\cdot9\cdot\frac{31}{20}q_*^3
 <\omega, \tag{14}
\]
so the scalar bound in (7) follows. In particular
$\sigma_j/\lambda_j>0$, consistently with the observable nonzero-score
branch. QED.


\end{proof}
\end{lemma}
\begin{proposition}[balanced scale and quotient entry]\label{prop:step-008-seed-interface}
Under the conclusions of Lemma~\ref{lem:step-008-seed-decomp} and
Lemma~\ref{lem:step-008-best-scalar}, the
setting's equal-norm seed satisfies
\[
 \gamma_j^0:=d_{U,j}^0d_{V,j}^0d_{W,j}^0=\sigma_j,\qquad
 s_j^0:=\log(\gamma_j^0/\lambda_j),\qquad
 \|s^0\|_\infty\le128\omega, \tag{15}
\]
and, with $E^0={\rm diag}(e^{s_j^0/3})$,
\[
 D_U^0=D_V^0=D_W^0=D_\lambda^{1/3}E^0,\qquad
 e^{-128\omega/3}\le\frac{d_{M,j}^0}{\lambda_j^{1/3}}
 \le e^{128\omega/3}. \tag{16}
\]

\begin{proof}

The three factors in (4) all have norm $\sigma_j^{1/3}$; the sign in the
third factor does not change its norm. Therefore their product of norms is
exactly $\gamma_j^0=\sigma_j$. By Lemma~\ref{lem:step-008-best-scalar}, write
$\sigma_j/\lambda_j=1+u_j$ with $|u_j|\le9\omega<1/2$. Then
\[
 |s_j^0|=|\log(1+u_j)|
 \le\frac{|u_j|}{1-|u_j|}
 \le18\omega<128\omega. \tag{17}
\]
The definition of equal-norm balancing now gives
\[
 d_{M,j}^0=(\sigma_j)^{1/3}
 =\lambda_j^{1/3}e^{s_j^0/3},
\]
which is (16) and the displayed diagonal identity. QED.

\end{proof}
\end{proposition}
\begin{proof}
Lemma~\ref{lem:step-006-gauge} first identifies, purely for proof,
one selected representative per realized target and a product-preserving
orientation. Lemma~\ref{lem:step-008-seed-decomp} checks that the normalized
observable best-scalar seed is exactly this oriented selected state. Because
the selected state is already in the exact realized target span, its
orthogonal complement is identically zero; the row/column closure
then gives the public $48\omega$ directional interface in both induced
orientations and the stronger $N^0=0\le\tau_r$ bound.

Lemma~\ref{lem:step-008-best-scalar} uses only the same selected-state correlations,
the target Gram margin, and the off-target ledger. It first derives
the diagonal correlation error, then bounds the weighted off-target score
contribution before taking a logarithm. This proves the observable ratio
estimate without a latent scale or any post-landing object.
Proposition~\ref{prop:step-008-seed-interface} converts the ratio to the exact relative
product-log coordinate and records the positive scale congruence. Together
these results produce
\[
 E_{\rm seed}:\quad
 P_M^0=G_M(I+C_M^0)+N_M^0,\quad G_M^\top N_M^0=0,\quad
 \|C_M^0\|_{\rm row,1}\vee\|C_M^0\|_{\rm col,1}\le48\omega,\quad
 \max_j\|N_M^0(:,j)\|_2\le\tau_r,\quad
 \|s^0\|_\infty\le128\omega.
\]
This event holds before any frozen design or landing solve and supplies the
input to the synchronized landing argument and the separate exact-baseline
specialization; no chronological predecessor, post-solve normalizer, or
cyclic invariant enters the proof.
\end{proof}

\subsection{Frozen landing reserves}\label{app:step-009}

The static geometry interface is Proposition~\ref{prop:step-001-geometry}; it
gives \(q\le q_*\), \(\Gamma\le1.01\), and the displayed induced row/column
Gram bounds for every realized \(G_M\). The balanced seed interface is
Proposition~\ref{prop:step-008-seed-interface}; it gives the decomposition,
orthogonality, \(48\omega\) row/column fields, the \(\tau_r\)
perpendicular-column bound, unit normalized columns, and
\(D_M^0=D_\lambda^{1/3}E^0\). No landing object is imported.

For $\nu\in\{{\rm row},{\rm col}\}$ we use
$\|AB\|_\nu\le\|A\|_\nu\|B\|_\nu$ and
$\|A\circ B\|_\nu\le\|A\|_\nu\|B\|_{\max}$. When $\|E\|_\nu<1$, the
Neumann identity $(I+E)^{-1}=\sum_{m\ge0}(-E)^m$ holds in the same induced
norm. All hypotheses are checked below.

For each \(M\in\{U,V,W\}\), let \(M',M''\) denote the other two modes in
the order used by the setting's frozen Khatri--Rao design, and define
\[
 H_M:=G_M^\top G_M,\qquad
 L_M:=G_M^\top P_M^0,\qquad
 F_M:=(P_M^0)^\top P_M^0.
\]
The setting-derived budgets are
\[
\delta_L=(1+q_*)c_0,\qquad
\delta_F=2(1+q_*)c_0+(1+q_*)c_0^2+r\tau_r^2,
\]
\[
\delta_{FL}=(1+q_*)c_0+(1+q_*)c_0^2+r\tau_r^2,\qquad
\alpha_0=q_*+\delta_F,
\]
\[
\eta_J=\alpha_0^2,\qquad
\eta_{QJ}=2\delta_{FL}+\delta_L^2+\delta_F^2,\qquad
\eta_A=\frac{\eta_{QJ}}{1-\eta_J}.
\]


\begin{proposition}[fixed numerical slack]\label{prop:step-009-numerical-slack}
Under Assumption~\ref{assump:subquadratic-rank} and the geometry conclusion
of Proposition~\ref{prop:step-001-geometry}, in particular
\(\Gamma\le1.01\), with \(q_*=1/4096\) and
\(\tau_r=q_*^2/(10^4r)\), the setting-defined budgets satisfy
\[
c_0<49q_*^2,\quad \delta_L<50q_*^2,\quad
\delta_F<101q_*^2,\quad \delta_{FL}<51q_*^2, \tag{N1}
\]
\[
\eta_J<1.051q_*^2<10^{-6},\qquad
\eta_{QJ}<103q_*^2,\qquad
\eta_A<104q_*^2, \tag{N2}
\]
and, using the \(\Gamma\le1.01\),
\[
\eta_J\le\frac1{64},\qquad
\Gamma\eta_A\le128\omega<\frac1{64}. \tag{N3}
\]

\begin{proof}
Since
\(\omega=(1+10^{-4})q_*^2\),
\[
c_0=48(1+10^{-4})q_*^2<49q_*^2.
\]
Also \(r\tau_r^2=q_*^4/(10^8r)\le q_*^4\) because \(r\ge3\). Therefore
\[
\delta_L<50q_*^2,\qquad
\delta_F<100q_*^2+4803q_*^4<101q_*^2,
\]
\[
\delta_{FL}<50q_*^2+4803q_*^4<51q_*^2,
\]
where \(4803q_*^2<1\). Consequently
\(\alpha_0<q_*(1+101/4096)<1.025q_*\), so
\(\eta_J<1.051q_*^2<10^{-6}\). Further,
\[
\eta_{QJ}<2(51q_*^2)+50^2q_*^4+101^2q_*^4
<103q_*^2.
\]
Thus \(1/(1-\eta_J)<1.000002\), giving
\(\eta_A<104q_*^2\). Finally,
\[
\Gamma\eta_A<1.01\cdot104q_*^2
<128(1+10^{-4})q_*^2=128\omega<1/64.
\]
This proves (N1)--(N3). 


\end{proof}
\end{proposition}
\begin{lemma}[two-orientation frozen Gram expansion]\label{lem:step-009-raw-expansion}
Under the geometry conclusion of Proposition~\ref{prop:step-001-geometry}
and the balanced-seed conclusion of Proposition~\ref{prop:step-008-seed-interface},
for every mode \(M\),
\[
\|L_M-H_M\|_{\rm row,1}\vee\|L_M-H_M\|_{\rm col,1}\le\delta_L, \tag{1}
\]
\[
\|F_M-H_M\|_{\rm row,1}\vee\|F_M-H_M\|_{\rm col,1}\le\delta_F, \tag{2}
\]
\[
\|F_M-L_M\|_{\rm row,1}\vee\|F_M-L_M\|_{\rm col,1}\le\delta_{FL}. \tag{3}
\]

\begin{proof}

The geometry event gives
\(\|H_M\|_{\rm row,1}\vee\|H_M\|_{\rm col,1}\le1+q_*\). From the seed
decomposition and \(G_M^\top N_M^0=0\),
\[
L_M=H_M(I+C_M^0),
\]
so \(L_M-H_M=H_MC_M^0\). Submultiplicativity in each orientation gives
(1).

Expanding the full seed Gram, with the two cross terms involving \(N_M^0\)
vanishing by \(G_M^\top N_M^0=0\), gives
\[
F_M-H_M=(C_M^0)^\top H_M+H_MC_M^0
 +(C_M^0)^\top H_MC_M^0+(N_M^0)^\top N_M^0. \tag{4}
\]
For every pair \(i,j\),
\(\left|((N_M^0)^\top N_M^0)_{ij}\right|\le\tau_r^2\), hence each induced row and
column sum is at most \(r\tau_r^2\). Applying the induced-norm product
inequality to (4), and using that transpose exchanges row and column norms,
proves (2).

Subtracting \(L_M=H_M(I+C_M^0)\) from \(F_M\) instead yields
\[
F_M-L_M=(C_M^0)^\top H_M(I+C_M^0)+(N_M^0)^\top N_M^0. \tag{5}
\]
The same two-orientation calculation proves (3). All bounds use the fixed
upper value \(q_*\), so the result is simultaneous over \(M=U,V,W\).



\end{proof}
\end{lemma}
\begin{proposition}[normalized pair-Gram and direct Q-J reserves]\label{prop:step-009-pair-reserves}
Under Proposition~\ref{prop:step-009-numerical-slack}
and the conclusion of Lemma~\ref{lem:step-009-raw-expansion}, define
\[
J_M^0:=F_{M'}\circ F_{M''},\qquad
Q_M^0:=L_{M'}\circ L_{M''}.
\]
Then, in both induced orientations,
\[
\|F_M-I\|_{\rm row,1}\vee\|F_M-I\|_{\rm col,1}\le\alpha_0, \tag{6}
\]
\[
\|J_M^0-I\|_{\rm row,1}\vee\|J_M^0-I\|_{\rm col,1}\le\eta_J, \tag{7}
\]
\[
\|Q_M^0-J_M^0\|_{\rm row,1}\vee\|Q_M^0-J_M^0\|_{\rm col,1}
\le\eta_{QJ}. \tag{8}
\]
Moreover, \(\eta_J\le1/64\) and
\[
\lambda_{\min}(J_M^0)\ge1-\eta_J,\qquad
\|(J_M^0)^{-1}\|_{\rm row,1}\vee\|(J_M^0)^{-1}\|_{\rm col,1}
\le\frac1{1-\eta_J}, \tag{9}
\]
\[
\|(J_M^0)^{-1}-I\|_{\rm row,1}\vee\|(J_M^0)^{-1}-I\|_{\rm col,1}
\le\frac{\eta_J}{1-\eta_J}. \tag{10}
\]

\begin{proof}

Because \(P_M^0\) and \(G_M\) have unit columns, every entry of \(F_M\),
\(L_M\), and \(H_M\) has absolute value at most one. In particular the
diagonal of \(F_M\) is one. Set \(E_M:=F_M-I\). Equation (2) and
\(\|H_M-I\|_\nu\le q_*\) give (6), and \(E_M\) has zero diagonal. Therefore
\[
J_M^0=(I+E_{M'})\circ(I+E_{M''})
  =I+E_{M'}\circ E_{M''}. \tag{11}
\]
For either \(\nu={\rm row}\) or \(\nu={\rm col}\),
\[
\|E_{M'}\circ E_{M''}\|_\nu
\le\|E_{M'}\|_\nu\|E_{M''}\|_\nu
\le\alpha_0^2,
\]
which proves (7).

For the direct paired cancellation, use the exact identity
\[
Q_M^0-J_M^0
 =(L_{M'}-F_{M'})\circ L_{M''}
  +F_{M'}\circ(L_{M''}-F_{M''}). \tag{12}
\]
The unit-column observation gives
\(\|L_{M''}\|_{\max}\vee\|F_{M'}\|_{\max}\le1\). Thus (3) and the
Hadamard bound yield the stronger estimate
\[
\|Q_M^0-J_M^0\|_\nu\le2\delta_{FL}
\le2\delta_{FL}+\delta_L^2+\delta_F^2=\eta_{QJ},
\]
proving (8) without inverting either pair. The two nonnegative square terms
are retained because they are part of the binding reserve definition.

Proposition~\ref{prop:step-009-numerical-slack} gives \(\eta_J<1\), so the
Neumann series for \(J_M^0=I+(J_M^0-I)\) converges in each induced norm and
gives (9)--(10).
The Schur-product/Khatri--Rao identity says that \(J_M^0\) is the Gram
matrix of the columnwise Kronecker products
\(P_{M'}^0\odot P_{M''}^0\), so it is symmetric positive semidefinite.
Together with the row bound in (7) (or
\(\|J_M^0-I\|_2\le\eta_J\)), Gershgorin gives
\(\lambda_{\min}(J_M^0)\ge1-\eta_J\). 


\end{proof}
\end{proposition}
\begin{lemma}[positive diagonal congruence and ordinary frozen solve]\label{lem:step-009-congruence}
Under Proposition~\ref{prop:step-009-pair-reserves} and the scale identity
in Proposition~\ref{prop:step-008-seed-interface}, let
\[
P_{-M}^0:=P_{M'}^0\odot P_{M''}^0,\qquad
K_{-M}:=G_{M'}\odot G_{M''},\qquad
D_{-M}^0:=D_{M'}^0D_{M''}^0.
\]
Then \(D_{-M}^0=D_\lambda^{2/3}(E^0)^2\succ0\) and
\[
(H_M^0)^\top H_M^0=D_{-M}^0J_M^0D_{-M}^0. \tag{13}
\]
Consequently the Moore--Penrose frozen-input update is ordinary on this
event and, for every mode,
\[
\widetilde X_M=G_M\widetilde B_M,\qquad
\widetilde B_M=D_\lambda Q_M^0(J_M^0)^{-1}(D_{-M}^0)^{-1}. \tag{14}
\]

\begin{proof}

The scale identity from the balanced-seed conclusion of Proposition~\ref{prop:step-008-seed-interface} gives
\[
D_{-M}^0=(D_\lambda^{1/3}E^0)^2
=D_\lambda^{2/3}(E^0)^2,
\]
which is positive diagonal. Since
\(X_{M'}^0=P_{M'}^0D_{M'}^0\) and
\(X_{M''}^0=P_{M''}^0D_{M''}^0\), the frozen Khatri--Rao design in the
setting is
\[
H_M^0=P_{-M}^0D_{-M}^0.
\]
The Khatri--Rao Gram identity gives
\[
(P_{-M}^0)^\top P_{-M}^0
=F_{M'}\circ F_{M''}=J_M^0,
\]
proving (13). Proposition~\ref{prop:step-009-pair-reserves} makes both
factors on the right invertible, so the Moore--Penrose inverse equals the
ordinary inverse.

For completeness, the current unfolding identity is
\[
T_{(M)}=G_MD_\lambda K_{-M}^\top,\qquad
K_{-M}^\top P_{-M}^0=L_{M'}\circ L_{M''}=Q_M^0. \tag{15}
\]
Using (13) and the diagonal inverse gives
\[
\begin{aligned}
T_{(M)}H_M^0((H_M^0)^\top H_M^0)^{-1}
 &=G_MD_\lambda Q_M^0D_{-M}^0
   (D_{-M}^0)^{-1}(J_M^0)^{-1}(D_{-M}^0)^{-1}\\
 &=G_MD_\lambda Q_M^0(J_M^0)^{-1}(D_{-M}^0)^{-1},
\end{aligned}
\]
which is (14). This calculation is done for all three frozen inputs before
any output is committed. On a singular or otherwise off-event tape the
setting's Moore--Penrose rule remains the observable fallback; the ordinary
formula and quantitative reserves are asserted only on this generated
pre-solve event. 


\end{proof}
\end{lemma}
\begin{proposition}[scale-conjugated coefficient and normalizer reserve]\label{prop:step-009-normalizer}
Under Proposition~\ref{prop:step-009-numerical-slack}
and Lemma~\ref{lem:step-009-congruence}, for
every mode define
\[
L_M^0:=D_\lambda(D_{-M}^0)^{-1}
   =D_\lambda^{1/3}(E^0)^{-2},\qquad
\ell_{M,j}^0:=(L_M^0)_{jj}
   =\lambda_j^{1/3}e^{-2s_j^0/3}>0,
\]
and
\[
\Xi_M:=D_\lambda Q_M^0(J_M^0)^{-1}D_\lambda^{-1}.
\]
Then, in both induced orientations,
\[
\|\Xi_M-I\|_{\rm row,1}\vee\|\Xi_M-I\|_{\rm col,1}
\le\Gamma\eta_A\le128\omega<\frac1{64}. \tag{16}
\]
If \(\widetilde d_{M,j}:=\|G_M\widetilde B_M(:,j)\|_2\), then
\[
(1-\Gamma\eta_A)\ell_{M,j}^0
\le\widetilde d_{M,j}
\le(1+\Gamma\eta_A)\ell_{M,j}^0, \tag{17}
\]
so every output normalizer is strictly positive before the joint commit.

\begin{proof}

By Proposition~\ref{prop:step-009-numerical-slack}, \(\eta_J\le1/64\) and
\(\Gamma\eta_A\le128\omega<1/64\); the explicit arithmetic is (N1)--(N3).

By (14), \(\widetilde B_M=\Xi_M L_M^0\). From (8)--(10),
\[
Q_M^0(J_M^0)^{-1}-I
=(Q_M^0-J_M^0)(J_M^0)^{-1},
\]
whose row and column norms are at most \(\eta_A\). Diagonal similarity
multiplies either induced norm by at most
\(\max_i\lambda_i/\min_i\lambda_i=\Gamma\), proving the first inequality
in (16).

For a fixed column \(j\), write \(\Xi_M(:,j)=e_j+v_j\). The column part of
(16) gives \(\|v_j\|_1\le\Gamma\eta_A\). Since every column of \(G_M\) has
norm one,
\[
\|G_M\Xi_M(:,j)-G_M(:,j)\|_2
\le\|v_j\|_1\le\Gamma\eta_A.
\]
The reverse and forward triangle inequalities give
\(1-\Gamma\eta_A\le\|G_M\Xi_M(:,j)\|_2
\le1+\Gamma\eta_A\). Multiplication by the positive
\(\ell_{M,j}^0\) proves (17). All three modes satisfy the same bounds in a
single pre-solve pass. 

\end{proof}
\end{proposition}
\begin{proof}
Lemma~\ref{lem:step-009-raw-expansion} expands the exact frozen seed in the
target Gram/cross-Gram convention and supplies \(\delta_L,\delta_F,\delta_{FL}\)
in both orientations, including the shared perpendicular charge
\(r\tau_r^2\). Proposition~\ref{prop:step-009-pair-reserves} uses the zero
diagonal of \(F_M-I\) for the Hadamard cancellation in \(J_M^0\), and uses
the direct \(Q_M^0-J_M^0\) identity before either pair is inverted. Its
Neumann bounds and the numerical check establish the normalized pair floor
and both inverse reserves simultaneously for \(M=U,V,W\).

Lemma~\ref{lem:step-009-congruence} transfers those normalized reserves to the
actual frozen designs through the exact positive diagonal congruence
\((H_M^0)^\top H_M^0=D_{-M}^0J_M^0D_{-M}^0\), and supplies the
ordinary-inverse coefficient formula for all three independent Moore--Penrose
calls. Finally, Proposition~\ref{prop:step-009-normalizer} controls the only
scale similarity by the generated \(\Gamma\le1.01\), proves
\(\Gamma\eta_A\le128\omega<1/64\), and gives strictly positive output
normalizers before any solve is committed. These four results give the
pre-solve reserve conclusion of Proposition~\ref{prop:step-009-normalizer},
with no post-solve or predecessor input.
\end{proof}

\begin{proposition}[Complete frozen pre-solve reserve interface]
\label{prop:complete-frozen-reserve-interface}
Under the event supplied by Proposition~\ref{prop:step-001-geometry} and the
balanced-seed conclusion of Proposition~\ref{prop:step-008-seed-interface},
define
\[
H_M:=G_M^\top G_M,\qquad L_M:=G_M^\top P_M^0,\qquad
F_M:=(P_M^0)^\top P_M^0,
\]
\[
D_{-M}^0:=D_{M'}^0D_{M''}^0,\qquad
L_M^0:=D_\lambda(D_{-M}^0)^{-1},\qquad
\ell_{M,j}^0:=(L_M^0)_{jj}>0.
\]
Put $c_0:=48\omega$ and define
\[
\delta_L=(1+q_*)c_0,\qquad
\delta_F=2(1+q_*)c_0+(1+q_*)c_0^2+r\tau_r^2,
\]
\[
\delta_{FL}=(1+q_*)c_0+(1+q_*)c_0^2+r\tau_r^2,\qquad
\alpha_0=q_*+\delta_F,\qquad
\eta_J=\alpha_0^2,\qquad
\eta_{QJ}=2\delta_{FL}+\delta_L^2+\delta_F^2,\qquad
\eta_A=\frac{\eta_{QJ}}{1-\eta_J}.
\]
For every mode $M$ and $\nu\in\{{\rm row},{\rm col}\}$, the frozen seed
satisfies
\[
\|L_M-H_M\|_\nu\le\delta_L,\qquad
\|F_M-H_M\|_\nu\le\delta_F,\qquad
\|F_M-L_M\|_\nu\le\delta_{FL},
\]
and, with $J_M^0=F_{M'}\circ F_{M''}$ and
$Q_M^0=L_{M'}\circ L_{M''}$,
\[
\|J_M^0-I\|_\nu\le\eta_J\le\frac1{64},\qquad
\|(J_M^0)^{-1}\|_\nu\le\frac1{1-\eta_J},\qquad
\|Q_M^0-J_M^0\|_\nu\le\eta_{QJ}.
\]
Moreover,
\[
\Xi_M:=D_\lambda Q_M^0(J_M^0)^{-1}D_\lambda^{-1},\qquad
\|\Xi_M-I\|_\nu\le\Gamma\eta_A\le128\omega<\frac1{64},
\]
\[
(H_M^0)^\top H_M^0=D_{-M}^0J_M^0D_{-M}^0,\qquad
\widetilde X_M=G_M\Xi_ML_M^0,\qquad
(1-\Gamma\eta_A)\ell_{M,j}^0
\le\|\widetilde x_{M,j}\|_2
\le(1+\Gamma\eta_A)\ell_{M,j}^0.
\]
In particular, every frozen design has full column rank and every output
normalizer is positive before any mode is committed.

\begin{proof}
Lemma~\ref{lem:step-009-raw-expansion} gives the three first inequalities in
both orientations. Proposition~\ref{prop:step-009-pair-reserves} applies the
zero-diagonal cancellation and Neumann bounds, yielding the $J_M^0$ floor,
inverse, and direct $Q_M^0-J_M^0$ estimates. Lemma~\ref{lem:step-009-congruence}
gives the positive diagonal congruence and ordinary-inverse coefficient
identity. Proposition~\ref{prop:step-009-normalizer} applies the only scale
similarity and proves the $\Xi_M$ and normalizer bounds. These are all
pre-commit conclusions from the same seed.
\end{proof}
\end{proposition}

\subsection{Synchronized frozen landing}\label{app:step-010}

Proposition~\ref{prop:complete-frozen-reserve-interface} supplies, before any
landing commit,
\[
\widetilde X_M
=G_MD_\lambda Q_M^0(J_M^0)^{-1}(D_{-M}^0)^{-1}
=G_M\Xi_ML_M^0. \tag{C1}
\]
Right multiplication preserves the left column range, positive diagonal
congruences are invertible, and for $0<\beta<1$ and
$a\in[1-\beta,1+\beta]$,
\[
|\log a|\le-\log(1-\beta)\le\frac{\beta}{1-\beta}. \tag{C2}
\]
The interval for each normalizer is supplied by the same pre-solve result.

For \(M\in\{U,V,W\}\), let \(M',M''\) be the other modes and put
\[
K_{-M}:=G_{M'}\odot G_{M''},\qquad
P_{-M}^0:=P_{M'}^0\odot P_{M''}^0.
\]
These are proof-local translations of the frozen Khatri--Rao factors.


\begin{lemma}[unconditional unfolding-range inclusion]\label{lem:step-010-range}
Under the exact tensor decomposition, for every frozen
design, including a singular one,
\[
\widetilde X_M:=T_{(M)}H_M^0\big((H_M^0)^\top H_M^0\big)^\dagger
\]
has every column in \({\rm range}(G_M)\).

\begin{proof}
The exact unfolding identity is
\[
T_{(M)}=G_MD_\lambda K_{-M}^{\top}. \tag{1}
\]
Hence \({\rm range}(T_{(M)})\subseteq{\rm range}(G_M)\), and right
multiplication by \(H_M^0((H_M^0)^\top H_M^0)^\dagger\) preserves that
left column range.  This uses no inverse or good event and applies
independently to all three modes.


\end{proof}
\end{lemma}
\begin{proposition}[three independent frozen solves and positive normalizers]\label{prop:step-010-frozen-solves}
Under the pre-solve reserve supplied by Proposition~\ref{prop:step-009-normalizer}
and Lemma~\ref{lem:step-009-congruence}, for each mode
\[
\widetilde X_M=G_M\widetilde B_M,\qquad
\widetilde B_M=D_\lambda Q_M^0(J_M^0)^{-1}(D_{-M}^0)^{-1}
=\Xi_ML_M^0. \tag{2}
\]
Writing \(\xi_{M,j}:=\Xi_M(:,j)\),
\[
\widetilde x_{M,j}=\ell_j^0G_M\xi_{M,j},\quad
\widetilde d_{M,j}:=\|\widetilde x_{M,j}\|_2
=\ell_j^0a_{M,j},\quad a_{M,j}:=\|G_M\xi_{M,j}\|_2, \tag{3}
\]
and \(1-\beta\le a_{M,j}\le1+\beta\), so
\(\widetilde d_{M,j}>0\).  All three formulas are evaluated from the one
frozen seed before any mode is committed.

\begin{proof}
The frozen scaling and the congruence give
\[
H_M^0=P_{-M}^0D_{-M}^0,\qquad
(H_M^0)^\top H_M^0=D_{-M}^0J_M^0D_{-M}^0. \tag{4}
\]
Lemma~\ref{lem:step-009-congruence} gives $D_{-M}^0\succ0$, while the
Step~009 pair-Gram reserve recorded in
Proposition~\ref{prop:complete-frozen-reserve-interface} gives
$\lambda_{\min}(J_M^0)\ge1-\eta_J>0$. Hence $H_M^0$ has full column rank and
its Moore--Penrose inverse is the ordinary inverse. Since
\[
K_{-M}^{\top}P_{-M}^0
=(G_{M'}^\top P_{M'}^0)\circ(G_{M''}^\top P_{M''}^0)
=Q_M^0,
\]
the unfolding identity yields
\[
\begin{aligned}
T_{(M)}H_M^0\big((H_M^0)^\top H_M^0\big)^{-1}
&=G_MD_\lambda Q_M^0D_{-M}^0
 (D_{-M}^0)^{-1}(J_M^0)^{-1}(D_{-M}^0)^{-1}\\
&=G_MD_\lambda Q_M^0(J_M^0)^{-1}(D_{-M}^0)^{-1},
\end{aligned}
\]
which is (2).  Factoring \(L_M^0\) gives (3).  For
\(v_{M,j}:=\xi_{M,j}-e_j\), the coefficient reserve from
Proposition~\ref{prop:step-009-normalizer} and the unit columns of $G_M$ imply
\[
|a_{M,j}-1|\le\|G_Mv_{M,j}\|_2\le\|v_{M,j}\|_1\le\beta. \tag{5}
\]
Thus the interval and positive normalizer follow. The singular
Moore--Penrose branch was not altered; the ordinary formula is used only
after the Step~009 congruence and pair-Gram floor have been established.


\end{proof}
\end{proposition}
\begin{proposition}[normalized direction fields]\label{prop:step-010-directions}
On the pre-solve reserve event supplied by Proposition~\ref{prop:step-009-normalizer},
assume the conclusions of Lemma~\ref{lem:step-009-congruence} and
Proposition~\ref{prop:step-009-normalizer}; in particular, for every
\(M\in\{U,V,W\}\),
\[
\beta:=\Gamma\eta_A\le128\omega<1/64<1,\qquad
\|\Xi_M-I\|_{\rm row,1}\vee\|\Xi_M-I\|_{\rm col,1}\le\beta,
\]
where \(\Xi_M:=D_\lambda Q_M^0(J_M^0)^{-1}D_\lambda^{-1}\), and the positive
diagonal congruence and ordinary frozen-solve identity hold.
Let the three outputs and their positive normalizers be those supplied by
Proposition~\ref{prop:step-010-frozen-solves}, so that
\[
\widetilde d_{M,j}>0,\qquad
a_{M,j}:=\|G_M\Xi_M(:,j)\|_2\in[1-\beta,1+\beta].
\]
Define \(D_{a,M}:={\rm diag}(a_{M,1},\ldots,a_{M,r})\) and
\[
P_M^{\rm land}:=\big[\widetilde x_{M,j}/\widetilde d_{M,j}\big]_{j=1}^r,
\qquad C_M^{\rm land}:=\Xi_MD_{a,M}^{-1}-I. \tag{6}
\]
Then \(P_M^{\rm land}=G_M(I+C_M^{\rm land})\), its perpendicular field is
zero, and, for \(\nu\in\{{\rm row},{\rm col}\}\),
\[
\|C_M^{\rm land}\|_\nu\le\frac{2\beta}{1-\beta}<3\beta,\qquad
\max_j\|G_MC_M^{\rm land}(:,j)\|_2
\le\frac{2\beta}{1-\beta}<3\beta. \tag{7}
\]
Every output direction has positive inner product with its corresponding
target column \(G_M(:,j)\).

\begin{proof}
From (5), \(D_{a,M}\) is positive and
\(\|D_{a,M}^{-1}\|_\nu\le(1-\beta)^{-1}\).  The identity
\[
\Xi_MD_{a,M}^{-1}-I
=(\Xi_M-I)D_{a,M}^{-1}+(D_{a,M}^{-1}-I) \tag{8}
\]
and Proposition~\ref{prop:step-009-normalizer} give two terms each at most
\(\beta/(1-\beta)\) in either
orientation, proving the first bound in (7).  Since
\(2/(1-\beta)=128/63<3\), the displayed strict bound follows.
For a column,
\[
\left\|\frac{G_M\xi_{M,j}}{a_{M,j}}-G_M(:,j)\right\|_2
\le\frac{\|G_Mv_{M,j}\|_2+|a_{M,j}-1|}{a_{M,j}}
\le\frac{2\beta}{1-\beta}.
\]
The range lemma gives zero perpendicular field.  Also
\(\langle G_M(:,j),G_M\xi_{M,j}\rangle\ge1-\beta>0\), so the proof-only
target sign chart is unchanged.


\end{proof}
\end{proposition}
\begin{proposition}[joint commit and product-preserving rebalance]\label{prop:step-010-rebalance}
On the pre-solve reserve event supplied by Proposition~\ref{prop:step-009-normalizer}, under
the setting's rule that all three frozen-input designs are formed before any
commit and their outputs are committed together, let
\[
(\widetilde X_U,\widetilde X_V,\widetilde X_W)
\]
be the three outputs from the common frozen seed
\((X^0,Y^0,Z^0)\) supplied by
Proposition~\ref{prop:step-010-frozen-solves}.  Assume, as that proposition
concludes, that every active normalizer satisfies
\(\widetilde d_{M,j}>0\) before commitment.  With
\[
\widetilde\gamma_j:=\prod_{M\in\{U,V,W\}}\widetilde d_{M,j}>0,
\]
set
\[
X_M^{\rm land}(:,j):=\widetilde\gamma_j^{1/3}
\frac{\widetilde X_M(:,j)}{\widetilde d_{M,j}}. \tag{9}
\]
This preserves each active rank-one product exactly and leaves normalized
directions unchanged.

\begin{proof}
Proposition~\ref{prop:step-010-frozen-solves}
states that every mode formula uses only the same frozen
\((X^0,Y^0,Z^0)\), so the three outputs are computed and then committed as
one transaction.  Its positive-normalizer conclusion gives
\(\widetilde\gamma_j>0\), and
\[
X_U^{\rm land}(:,j)\otimes X_V^{\rm land}(:,j)\otimes X_W^{\rm land}(:,j)
=\frac{\widetilde\gamma_j}
{\widetilde d_{U,j}\widetilde d_{V,j}\widetilde d_{W,j}}
\widetilde x_j\otimes\widetilde y_j\otimes\widetilde z_j
=\widetilde x_j\otimes\widetilde y_j\otimes\widetilde z_j. \tag{10}
\]
Only positive scales change, so \(P_M^{\rm land}\) and \(C_M^{\rm land}\)
are unchanged; inactive columns remain zero.


\end{proof}
\end{proposition}
\begin{proposition}[quotient entry after one rebalance]\label{prop:step-010-quotient-entry}
On the conjunction of the pre-solve reserve from Proposition~\ref{prop:step-009-normalizer}
and the seed conclusion of Proposition~\ref{prop:step-008-seed-interface}, assume
the seed-scale conclusion \(\|s^0\|_\infty\le128\omega\), and the reserve
\(\beta=\Gamma\eta_A\le128\omega<1/64\), and the following named
conclusions: Lemma~\ref{lem:step-010-range} gives exact target-span
membership for the Moore--Penrose outputs; Proposition~\ref{prop:step-010-frozen-solves}
gives \(\widetilde d_{M,j}=\ell_j^0a_{M,j}\) with
\(a_{M,j}\in[1-\beta,1+\beta]\); Proposition~\ref{prop:step-010-directions}
gives \(P_M^{\rm land}=G_M(I+C_M^{\rm land})\), zero perpendicular field, and
positive target orientation; and Proposition~\ref{prop:step-010-rebalance} gives the common
frozen-input joint commit and product-preserving rebalance.  Then
\[
\|s^{\rm land}\|_\infty\le1024\omega,\qquad
d_Q(e^{\rm land},0)\le1024\omega.
\tag{11}
\]
Consequently \(d_Q(e^{\rm land},0)\le4096\omega<\rho_{\rm ALS}/3\).

\begin{proof}
Since
\(\ell_j^0=\lambda_j^{1/3}e^{-2s_j^0/3}\), (3) gives
\[
\frac{\widetilde\gamma_j}{\lambda_j}
=e^{-2s_j^0}\prod_{M\in\{U,V,W\}}a_{M,j}.
\tag{12}
\]
The quotient log is therefore
\[
s_j^{\rm land}=-2s_j^0+\sum_{M\in\{U,V,W\}}\log a_{M,j}. \tag{13}
\]
By (C2) and (5), \(|\log a_{M,j}|\le\beta/(1-\beta)<2\beta\le256\omega\).
Thus
\[
|s_j^{\rm land}|\le2(128\omega)+3(256\omega)=1024\omega. \tag{14}
\]
The definition of \(d_Q\) takes the maximum of the three direction
quantities and the scale log.  Proposition~\ref{prop:step-010-directions}
bounds every direction quantity by \(3\beta\le384\omega\), so (11) follows.
Finally,
\[
\omega=q_*^2+r\tau_r=q_*^2(1+10^{-4}),\quad
q_*={1\over4096},\quad \rho_{\rm ALS}={1\over1024},
\]
and hence
\[
4096\omega={1+10^{-4}\over4096}
<{1\over3072}={\rho_{\rm ALS}\over3},
\tag{15}
\]
because \(1+10^{-4}<4/3\).

\end{proof}
\end{proposition}
\begin{proof}
Lemma~\ref{lem:step-010-range} gives exact target-span membership for each
Moore--Penrose output without assuming a nonsingular design.  On the derived
event \(E_{\rm land}^{\rm pre}\), Proposition~\ref{prop:step-010-frozen-solves}
invokes the positive diagonal congruence and pair-Gram floor supplied by
Lemma~\ref{lem:step-009-congruence} and
Proposition~\ref{prop:step-009-normalizer} to replace each pseudoinverse by
its own ordinary inverse, derives the exact coefficient identity, and proves
all active normalizers positive from the explicit \(\beta<1\) reserve.  Its
statement also records that all three outputs use the one frozen seed.

Under the same event, Proposition~\ref{prop:step-010-directions} applies the
normalizers to obtain the target-chart coefficient fields, both induced
direction bounds, zero perpendicular fields, and the fixed target
orientation.  The setting's synchronized joint-commit rule together with the
positive normalizers is discharged by
Proposition~\ref{prop:step-010-rebalance}, which proves exact preservation of
every represented term.  Finally, under the \(E_{\rm seed}\) scale
bound, Proposition~\ref{prop:step-010-quotient-entry} gives
\(d_Q(e^{\rm land},0)\le1024\omega\le4096\omega\) and the
\(\rho_{\rm ALS}/3\) entry.  No landing output is used as an input to another
landing solve; singular/off-event tapes retain the observable
Moore--Penrose fallback without an unsupported basin claim.
\end{proof}

\subsection{Cyclic target-span contraction}\label{app:step-011}

The landing and static-geometry inputs for this section are the conclusions of
Propositions~\ref{prop:step-010-quotient-entry} and
\ref{prop:step-001-geometry}.  We use the quotient interpretation in
Uschmajew~\cite{Uschmajew2012} only at the scope of Assumption~1, Lemma~3.2,
and Theorems~3.3 and~3.5 of that paper: the active factors are nonzero, the
block normal equations are nonsingular, and cyclic ALS is read modulo the
componentwise product-one scaling kernel with block-Gauss--Seidel chronology.
Here the object map is $G_M\in\{U,V,W\}$, the held Khatri--Rao design is
$H_M=X_{M'}\odot X_{M''}$, and the quotient representative is
$(C_U,C_V,C_W,s)$ with the setting's $\operatorname{Refresh}_s$ operation.
The target-span inclusion, the radius $\rho_{\rm ALS}$, every inverse and
contraction constant, and the refresh identity are proved in the current
notation below; none is imported from the cited source.

For induced row/column $\ell_1$ norms, submultiplicativity, the Neumann inverse bound, and the Khatri--Rao Gram identity are used after their hypotheses are displayed.

For matrices with unit diagonal, the following direct inequality will be
   used repeatedly:
   \[
   \|A\circ B\|_\nu
   \le \|A-\operatorname{diag}(A)\|_\nu
      \max_{i\ne j}|B_{ij}|
   +\max_i|A_{ii}B_{ii}|,
   \qquad \nu\in\{{\rm row},{\rm col}\}.
   \tag{C1}
   \]
   It follows by separating the diagonal term in each row or column sum.

Throughout this derivation, \(M'\) and \(M''\) denote the two modes other
than \(M\). For convenience, write
\[
\varrho:=\rho_{\rm ALS}=1/1024
\]
to distinguish the chart radius from the smoothing scale $\rho$.


\begin{lemma}[target-span quotient representation and balanced scales]\label{lem:step-011-quotient-chart}
Under the static event of Proposition~\ref{prop:step-001-geometry}, the
landing conclusion of Proposition~\ref{prop:step-010-quotient-entry}, and
Assumption~\ref{assump:base-scale}, let an active state have nonzero columns in the exact
target spans, and orient its columns by the landing permutation/sign chart.
If its represented component products are positive, then there are unique
matrices \(C_M\) and a unique vector \(s\) such that
\[
 P_M=G_M(I+C_M),\qquad
 s_j=\log(\gamma_j/\lambda_j),
 \qquad \gamma_j=\prod_M\|X_M(:,j)\|_2.
 \tag{1}
\]
After the product-preserving equal-norm representative is chosen,
\[
 D_M=\operatorname{diag}(\|X_M(:,j)\|_2)
  =D_\lambda^{1/3}E,
 \qquad E:=\operatorname{diag}(e^{s_j/3}),
 \tag{2}
\]
for all three modes. The exact target state is \(0=(0,0,0,0)\).

\begin{proof}
From Proposition~\ref{prop:step-001-geometry}, every \(G_M^\top G_M\) is positive
definite by Gershgorin, so the coefficient matrix in
\(P_M=G_MB_M\) is unique. Set \(C_M=B_M-I\). The landing signs and
permutation fix the column chart; later positive normalizations preserve it
when the coefficient error is below one. The positive product definition
gives the unique logarithm \(s\). Equal-norm balancing replaces the three
norms by \(\gamma_j^{1/3}\), which is positive and preserves each rank-one
product, yielding (2). No factor outside the target span is introduced by
this scaling.


\end{proof}
\end{lemma}
\begin{proposition}[scale equivariance and literal refresh identity]\label{prop:step-011-refresh-equivariance}
Under the hypotheses of
Lemma~\ref{lem:step-011-quotient-chart}, suppose the held design for a mode
\(M\) has full column rank and the solved active output satisfies
\(\min_j\|\widetilde X_M(:,j)\|_2>0\). Let positive diagonal matrices
\(R_U,R_V,R_W\) satisfy \(R_UR_VR_W=I\). Rescaling a literal state by
\(X_M\mapsto X_MR_M\) changes the mode-\(M\) ALS output to
\(\widetilde X_M\mapsto\widetilde X_MR_M\), and hence leaves its quotient
coordinates unchanged. Consequently, if \(e^t\) is the balanced quotient
representative of a literal state, then the three records
\[
 e\longmapsto e_U\longmapsto e_V\longmapsto e_W
 \tag{3}
\]
defined by the setting's Refresh$_s$ operation are quotient-equivalent to
the literal U/V/W post-block states, in that chronological order.

\begin{proof}
Write the literal mode design as
\(H_M=X_{M'}\odot X_{M''}\). Under the product-one rescaling,
\(H_M\mapsto H_MR_{M'}R_{M''}=H_MR_M^{-1}\). Full column rank lets the
Moore--Penrose solve be written as the ordinary normal-equation solve, so
\[
\begin{aligned}
T_{(M)}H_MR_M^{-1}
 \big((R_M^{-1})H_M^\top H_MR_M^{-1}\big)^{-1}
 &=T_{(M)}H_M(H_M^\top H_M)^{-1}R_M\\
 &=\widetilde X_MR_M.
\end{aligned}
\tag{4}
\]
Thus positive product-one scaling is a gauge symmetry of the literal update.
The displayed positive-output condition makes the canonical equal-norm map
well-defined; it changes only positive scales and preserves every rank-one
product, so it produces the same quotient point. The U record uses
the old V,W and writes the new U product log globally; the V record then uses
that new U and old W, and the W record uses the two current fields. Applying
(4) successively proves that (3) is exactly the quotient image of the
literal chronological orbit. This is the block-Gauss--Seidel chronology
described qualitatively by the cited quotient interpretation, but the identity is
the displayed current-notation calculation.


\end{proof}
\end{proposition}
\begin{lemma}[two-orientation dynamic Gram reserve]\label{lem:step-011-gram-reserve}
On the event supplied by Proposition~\ref{prop:step-001-geometry}, let
\(e=(C_U,C_V,C_W,s)\) be an exact target-span balanced state with
\(x:=d_Q(e,0)\le\varrho\). For each mode set
\[
 A_M:=I+C_M,\quad P_M:=G_MA_M,\quad
 L_M:=G_M^\top P_M=H_MA_M,\quad
 F_M:=P_M^\top P_M=A_M^\top H_MA_M.
\tag{5}
\]
Then, for \(\nu\in\{{\rm row},{\rm col}\}\), simultaneously for all
three modes,
\[
\begin{aligned}
 \|F_M-H_M\|_\nu&\le2.01x,\\
 \|(L_M-F_M)-\operatorname{diag}(L_M-F_M)\|_\nu&\le1.01x,\\
 \max_j|L_{M,jj}-F_{M,jj}|&\le x^2/2,\\
 \max_{i\ne j}|(L_M)_{ij}|&\le q+1.01x,\\
 \max_{i\ne j}|(F_M)_{ij}|&\le q+2.01x.
\end{aligned}
\tag{6}
\]
For a block M define
\[
 Q_M:=L_{M'}\circ L_{M''},\quad
 J_M:=F_{M'}\circ F_{M''},\quad
 K_M:=H_{M'}\circ H_{M''}.
\tag{7}
\]
Then
\[
 \|Q_M-J_M\|_\nu\le2.1qx+4.1x^2,
 \qquad
 \|J_M-I\|_\nu\le q^2+4.1qx+4.1x^2<1/64,
\tag{8}
\]
and therefore \(J_M\) is positive definite and
\[
 \|J_M^{-1}\|_{\rm row,1}\vee\|J_M^{-1}\|_{\rm col,1}\le64/63.
\tag{9}
\]

\begin{proof}
Write \(H=H_M\), \(A=I+C\), and note from
the definition of \(d_Q\) that both induced norms of C and every
\(\|GC(:,j)\|_2\) are at most x. Since \(\|H-I\|_\nu\le q\),
\(\|H\|_\nu\le1+q\), and \(x\le1/1024\),
\[
 \|F-H\|_\nu
 \le (1+q)x+(1+q)x+(1+q)x^2\le2.01x.
\tag{10}
\]
Here we used \(F-H=HC+C^\top H+C^\top HC\). Also
\[
 L-F=HA-A^\top HA=-C^\top HA,
\tag{11}
\]
so its induced norm is at most \((1+q)(1+x)x\le1.01x\).
Unit columns of P give the exact diagonal identity
\[
 L_{jj}-F_{jj}=L_{jj}-1=-\tfrac12\|GC(:,j)\|_2^2,
\tag{12}
\]
which proves the second and third lines of (6). The expansion
\(L-H=HC\) gives the fourth line of (6); the expansion in (10) gives the
fifth line.

For (8), apply the direct diagonal/off-diagonal inequality (C1) to
\[
 Q-J=(L_{M'}-F_{M'})\circ L_{M''}
   +F_{M'}\circ(L_{M''}-F_{M''}).
\tag{13}
\]
The first term is at most
\(1.01x(q+1.01x)+x^2/2\), and the second is at most
\((q+2.01x)(1.01x)+x^2/2\). Their sum is at most
\(2.02qx+4.051x^2\), hence the first bound in (8). Similarly, with
\(R_M=F_M-H_M\), whose diagonal is zero,
\[
 J-K=R_{M'}\circ F_{M''}+H_{M'}\circ R_{M''}
\tag{14}
\]
has norm at most \(4.1qx+4.1x^2\). The target Gram satisfies
\(\|K-I\|_\nu\le q^2\), because its off-diagonal entries are products of
the two target off-diagonal entries. This proves the second inequality in
(8). Finally \(q\le q_*\le\varrho/4\) and \(x\le\varrho\), so
\[
q^2+4.1qx+4.1x^2
\le(1/16+4.1/4+4.1)\varrho^2<1/64.
\tag{15}
\]
The Neumann series in either induced norm gives (9); because J is the Gram
of \(P_{M'}\odot P_{M''}\), it is symmetric positive semidefinite, and
the same bound gives its positive spectral floor. Thus the Moore--Penrose
update is the ordinary update throughout the chart.


\end{proof}
\end{lemma}
\begin{proposition}[one chronological block contraction]\label{prop:step-011-block-contraction}
On the event supplied by Proposition~\ref{prop:step-001-geometry},
and under the hypotheses and conclusions of
Lemmas~\ref{lem:step-011-quotient-chart} and
\ref{lem:step-011-gram-reserve}, fix a mode M and let
\(x=d_Q(e,0)\le\varrho\) for an exact-span balanced positive-product state.
Form the literal exact ALS update from the two held modes. The full-rank
reserve below first proves positive solved columns; only then apply the
product-preserving equal-norm representative. The landing conclusion of
Proposition~\ref{prop:step-010-quotient-entry}
supplies one such in-chart input, but the estimate below is uniform over
the entire displayed chart.
The new mode direction and its refreshed product log satisfy, in both
induced orientations and in the Euclidean column part of \(d_Q\),
\[
 \|C_M^+\|_{\rm row,1}\vee\|C_M^+\|_{\rm col,1}
 \vee\max_j\|G_MC_M^+(:,j)\|_2
 \vee\|s_M^+\|_\infty
 \le \ell x,
 \qquad \ell:=8q_*+32\varrho.
\tag{16}
\]
The ordinary full-rank output is in \({\rm range}(G_M)\), has positive
target orientation, and has positive active norm, so the canonical
representative is defined.

\begin{proof}
In the balanced representative, write
\(D_M=D_\lambda^{1/3}E\) as in (2) and
\(D_{-M}=D_{M'}D_{M''}=D_\lambda^{2/3}E^2\). The unfolding identity and
the full-rank conclusion of Lemma~\ref{lem:step-011-gram-reserve} give the
current normal-equation formula
\[
 \widetilde X_M
 =G_MD_\lambda Q_MJ_M^{-1}D_{-M}^{-1}
 =G_M\Xi_M D_\lambda^{1/3}E^{-2},
 \qquad
 \Xi_M:=D_\lambda Q_MJ_M^{-1}D_\lambda^{-1}.
\tag{17}
\]
Indeed, \(T_{(M)}=G_MD_\lambda(G_{M'}\odot G_{M''})^\top\) and
\((G_{M'}\odot G_{M''})^\top(P_{M'}\odot P_{M''})=Q_M\).
Using \(Q_MJ_M^{-1}-I=(Q_M-J_M)J_M^{-1}\), (8)--(9), and
\(\Gamma\le1.01\), diagonal similarity gives, for either \(\nu\),
\[
 \|\Xi_M-I\|_\nu
 \le1.01\frac{64}{63}(2.1qx+4.1x^2)
 \le2.2qx+4.3x^2.
\tag{18}
\]
Set \(\beta:=2.2qx+4.3x^2\). Since
\(2.2q_*+4.3\varrho<1/128\), we have \(\beta<1/128\). For each column
\(\xi_j=\Xi_M(:,j)\), let \(a_j=\|G_M\xi_j\|_2\). The column bound in
(18) gives \(|a_j-1|\le\beta\), hence \(a_j>0\). Moreover
\(\langle G_M(:,j),G_M\xi_j\rangle\ge1-\beta>0\), so the target
orientation is preserved. Put
\[
 b_{M,j}:=\lambda_j^{1/3}e^{-2s_j/3}>0,
 \qquad \widetilde x_{M,j}=b_{M,j}G_M\xi_j,
 \qquad \|\widetilde x_{M,j}\|_2=b_{M,j}a_j>0.
\tag{19a}
\]
Thus the positive solved-column condition in
Proposition~\ref{prop:step-011-refresh-equivariance} is discharged before
canonical normalization. The normalized coefficient field is
\[
 C_M^+=\Xi_M\operatorname{diag}(a_j^{-1})-I.
\tag{19}
\]
Consequently, in either induced orientation,
\[
 \|C_M^+\|_\nu
 \le {2\beta\over1-\beta}
 \le4.5qx+8.8x^2,
\tag{20}
\]
and the same bound holds for each Euclidean column error after applying
$G_M$. To identify the refreshed product log, the two held balanced norms are
\(\lambda_j^{1/3}e^{s_j/3}\) each. Therefore, before the canonical rebalance,
\[
 \gamma_j^+
 = (b_{M,j}a_j)
   (\lambda_j^{1/3}e^{s_j/3})^2
 = \lambda_j a_j.
\tag{20a}
\]
The rebalance preserves this rank-one product, so its global product register
satisfies the exact gauge cancellation
\[
 s_{M,j}^+
 =\log\!\left({\gamma_j^+\over\lambda_j}\right)=\log a_j.
\tag{20b}
\]
The elementary logarithm bound on \([1-\beta,1+\beta]\) now gives
\[
 |s_{M,j}^+|=|\log a_j|
 \le {\beta\over1-\beta}
 \le2.3qx+4.4x^2.
\tag{21}
\]
Since \(q\le q_*\) and \(x\le\varrho\), each right side in (20)--(21)
is at most
\[
 (8q_*+32\varrho)x=\ell x.
\tag{22}
\]
On the maintained chart, Lemma~\ref{lem:step-011-gram-reserve} makes the
current design Gram invertible, so the displayed ordinary formula (17) has a
left factor \(G_M\) and proves the output is in \({\rm range}(G_M)\). No
singular/off-chart range assertion is used in this cyclic step. The positive
normalizer just proved makes the normalized direction and canonical rebalance
well-defined. This proves (16).


\end{proof}
\end{proposition}
\begin{proposition}[chronological contraction and target-span invariant]\label{prop:step-011-chronological-contraction}
Under the event supplied by Proposition~\ref{prop:step-001-geometry}, let
Proposition~\ref{prop:step-010-quotient-entry} supply the entry state
\(e^{\rm land}\). Using
Proposition~\ref{prop:step-011-block-contraction} and then
Proposition~\ref{prop:step-011-refresh-equivariance} (the latter only after
the former has supplied positive solved columns), take any exact-span
balanced positive-product state \(e^0=e\) in the fixed orientation chart
with \(x:=d_Q(e,0)\le\varrho\), and define the setting's chronological states
\[
 e_U=(C_U^+,C_V^0,C_W^0,s_U^+),\quad
 e_V=(C_U^+,C_V^+,C_W^0,s_V^+),\quad
 e_W=(C_U^+,C_V^+,C_W^+,s_W^+),
 \qquad \Psi^q(e):=e_W.
\tag{23}
\]
Then all three blocks are defined, their outputs remain in the exact
target-span chart with positive active products, and their refreshed records
are quotient-equivalent to the literal ALS post-block states. Uniformly for
every such in-chart input,
\[
 d_Q(\Psi^q(e),0)\le\ell d_Q(e,0),
 \qquad \ell=8q_*+32\rho_{\rm ALS}=17/512<1/16.
\tag{24}
\]
In particular, starting from \(e^{\rm land}\), the chronological orbit
remains in the chart and defines the basin event; writing
\(e^{t+1}:=\Psi^q(e^t)\) and \(e^0=e^{\rm land}\), it has no additive forcing:
\[
 d_Q(e^t,0)\le\ell^t d_Q(e^{\rm land},0),
 \quad
 d_Q(e^t,0)-d_Q(e^{t+1},0)\ge(1-\ell)d_Q(e^t,0),
 \quad
 \sum_{t\ge0}d_Q(e^t,0)\le{d_Q(e^{\rm land},0)\over1-\ell}.
\tag{25}
\]

\begin{proof}
Fix an arbitrary input state in the domain of the
proposition and write \(x=d_Q(e,0)\le\varrho\). Lemma~\ref{lem:step-011-gram-reserve}
gives the full-rank held-design reserve for the U block. Proposition~\ref{prop:step-011-block-contraction}
therefore gives
\[
  \|C_U^+\|_{\rm row,1}\vee\|C_U^+\|_{\rm col,1}
  \vee\max_j\|G_UC_U^+(:,j)\|_2\vee\|s_U^+\|_\infty
  \le\ell x,
\tag{26}
\]
The positivity and target-span conclusions in the same proposition discharge
the local condition of Proposition~\ref{prop:step-011-refresh-equivariance};
the full-rank reserve and equation (4) then identify the U record with the
literal U post-block state. The unchanged V,W fields and their scale register
are still bounded by x, while the newly written U fields/register are bounded
by \(\ell x\), so \(d_Q(e_U,0)\le x\). Apply
the same two propositions to V with the new U field and old W field. Its input
is still in the chart, so the new V fields and \(s_V^+\) are at most
\(\ell x\), and \(d_Q(e_V,0)\le x\); the V record is again the literal
post-block state. Applying them to W gives the same bound for the new W
fields and \(s_W^+\). At the end of the sweep, U,V,W and the common product
register are all at most \(\ell x\), proving the uniform map inequality
(24), while equation (17) proves exact target-span membership before every
subsequent block update.

For the landing state, Proposition~\ref{prop:step-010-quotient-entry} gives
\(d_Q(e^{\rm land},0)<\varrho/3\). Since \(\ell<1\), the uniform map
inequality keeps every \(e^t\) in the chart and gives the first inequality
in (25) by induction. Subtracting successive nonnegative terms and summing
the geometric series gives the potential drop and finite budget in (25).

Finally,
\[
 8q_*+32\rho_{\rm ALS}
 ={8\over4096}+{32\over1024}
 ={17\over512}< {1\over16},
\tag{27}
\]
so the contraction is strict with the fixed sketch constant. A singular or
zero-normalizer tape outside the maintained chart receives no quantitative
conclusion and follows the setting's observable unsuccessful/cap branch.

\end{proof}
\end{proposition}
\begin{proof}
The landing conclusion of Proposition~\ref{prop:step-010-quotient-entry},
together with the event from Proposition~\ref{prop:step-001-geometry},
supplies the legal positive exact-span entry.
Lemma~\ref{lem:step-011-quotient-chart} supplies the unique exact-span
quotient coordinates and the balanced scale convention consumed by every
block. Proposition~\ref{prop:step-011-block-contraction} first proves
positive solved columns, target-span output, and the one-block bound on every
in-chart input; only after that discharge does
Proposition~\ref{prop:step-011-refresh-equivariance} prove that
positive product-one scaling is a genuine gauge symmetry of each full-rank
literal normal-equation solve and that the setting's chronological
Refresh$_s$ records are the quotient representatives of the literal U/V/W
orbit. Lemma~\ref{lem:step-011-gram-reserve} derives, rather than assumes,
the dynamic pair-Gram reserve from the realized interference bound
and the current radius; it also proves that every design used on the good
trajectory is full rank. Proposition~\ref{prop:step-011-block-contraction}
then retains the exact coefficient identity, displays the product/gauge
scale cancellation
\(s_M^+=\log a_j\), positive normalizers, target orientation, and all three
parts of the one-block \(d_Q\) bound. Finally,
Proposition~\ref{prop:step-011-chronological-contraction} composes the held
input bounds in the actual U/V/W order for every exact-span positive-product
input with \(d_Q\le\rho_{\rm ALS}\), proves target-span invariance before
each block update, and exports the uniform map estimate (24) and the
no-forcing recurrence (25). Therefore the generated basin event and
\(\Psi^q\) satisfy the exact displayed claim, with no predecessor
comparator, post-solve landing input, or hidden basin assumption.

At the exact orthogonal equal-weight certified state,
\(H_M=I\), \(C_M=0\), and \(s=0\). Hence \(Q_M=J_M=I\),
\(\Xi_M=I\), \(a_j=1\), every Refresh$_s$ register remains zero, and all
three cyclic updates are stationary. The quotient error and the represented
tensor residual are exactly zero, preserving the required baseline rather
than replacing it by the finite-radius bound.
\end{proof}

\subsection{Stopping, runtime, and restarts}\label{app:step-012}

Proposition~\ref{prop:step-003-coupon} and
Assumption~\ref{assump:random-initialization} give, conditional on a fixed
instance in the event of Proposition~\ref{prop:step-001-geometry},
\[
\Pr(E_{\rm cov}\mid{\cal F}_{\rm sm})\ge p_0,
\qquad p_0:=\frac{26}{27}.
\tag{A1}
\]
Proposition~\ref{prop:step-003-witness-interface} ensures that its witnesses
are proof-only. Lemma~\ref{lem:step-011-quotient-chart} and
Proposition~\ref{prop:step-011-chronological-contraction} give, on the
generated covered path,
\[
x_0:=d_Q(e^{\rm land},0)\le4096\omega<\rho_{\rm ALS}/3,
\qquad x_{t+1}\le\ell x_t,\qquad
\ell:=\frac{17}{512}<\frac1{16}.
\tag{A2}
\]
The same proposition gives exact-span invariance, positive products, and
literal ALS/refresh equivalence.

For the residual comparison we use
$(A\odot B)^\top(A\odot B)=(A^\top A)\circ(B^\top B)$ and Schur order:
if $H\succeq0$ has unit diagonal and $aI\preceq K\preceq bI$, then
$aI\preceq H\circ K\preceq bI$. We also use
$\|A\otimes B\|_2=\|A\|_2\|B\|_2$,
$\|C\|_2\le\sqrt{\|C\|_{\rm row,1}\|C\|_{\rm col,1}}$, and the isometry
$\Delta(z)=\sum_jz_je_j\otimes e_j\otimes e_j$. Every finite matrix has a
Moore--Penrose pseudoinverse; this gives existence and dense cost on singular
branches, not an inverse reserve. Conditional on the tensor, complete-run
tapes are independent by Assumption~\ref{assump:random-initialization}, so the
joint failure probability of completed runs is a product.

\begin{lemma}[balanced-chart residual transfer]\label{lem:step-012-residual-transfer}
Under Assumption~\ref{assump:base-scale}, the static-geometry interface of
Proposition~\ref{prop:step-001-geometry}, and the coordinate
conclusion of Lemma~\ref{lem:step-011-quotient-chart}, let
\(e=(C_U,C_V,C_W,s)\) be a state supplied by Proposition~\ref{prop:step-011-chronological-contraction} represented in
the exact-span, positively oriented, equal-norm chart, and suppose
\(x=d_Q(e,0)\le\rho_{\rm ALS}=1/1024\).  If \(\widehat T(e)\) denotes the
rank-\(r\) tensor represented by its active factors, with the inactive
\(k-r\) columns zero, then
\[
 { \|T-\widehat T(e)\|_F\over\|T\|_F}\le5x\le8\kappa_0^2x.
\tag{1}
\]

\begin{proof}
Write \(G_M=U,V,W\) in the corresponding mode,
\(A_M=I+C_M\), and \(\lambda=(\lambda_1,\ldots,\lambda_r)^\top\). The fixed orientation chart
and positive product register give
\[
 p_{M,j}=G_MA_Me_j,\qquad
 \widehat T(e)={\mathscr G}{\cal A}\Delta(\lambda\odot e^s),
 \qquad T={\mathscr G}\Delta(\lambda),
\tag{2}
\]
where \({\mathscr G}=G_U\otimes G_V\otimes G_W\),
\({\cal A}=A_U\otimes A_V\otimes A_W\), and \(e^s\) has entries \(e^{s_j}\).
This is the componentwise identity
\(\lambda_je^{s_j}p_{U,j}\otimes p_{V,j}\otimes p_{W,j}\), so it does not
add a factor or change the represented tensor.

By the definition of \(d_Q\), both induced norms of each \(C_M\) are at most
\(x\), so \(\|C_M\|_2\le x\) and
\[
 \|{\cal A}\|_2\le(1+x)^3,\qquad
 \|{\cal A}-I\|_2\le(1+x)^3-1.
\tag{3}
\]
For \(0\le x<1\), the power-series comparison \(e^x\le1/(1-x)\) gives
\[
 \|\lambda\odot e^s-\lambda\|_2
 \le(e^x-1)\|\lambda\|_2
 \le {x\over1-x}\|\lambda\|_2.
\tag{4}
\]
At \(x\le1/1024\), \((1+x)^3<1.003\),
\(x/(1-x)<1.001x\), and \((1+x)^3-1<3.004x\). The isometry of \(\Delta\)
therefore yields
\[
\begin{aligned}
 \|{\cal A}\Delta(\lambda\odot e^s)-\Delta(\lambda)\|_F
 &\le\|{\cal A}\|_2\|\Delta(\lambda\odot e^s-\lambda)\|_F
  +\|{\cal A}-I\|_2\|\Delta(\lambda)\|_F\\
 &<(1.003)(1.001)x\|\lambda\|_2+3.004x\|\lambda\|_2
 <4.1x\|\lambda\|_2.
\end{aligned}
\tag{5}
\]
The static geometry gives
\(\|G_M\|_2^2=\lambda_{\max}(G_M^\top G_M)\le1+q_*\), and hence
\[
 \|{\mathscr G}\|_2\le(1+q_*)^{3/2}.
\tag{6}
\]
It remains to compare the denominator with the same coefficient norm used in
(5).  Put \(H_U=U^\top U\), \(H_V=V^\top V\), and \(H_W=W^\top W\), and
define the two-mode and full three-mode Khatri--Rao Grams
\[
 K_{VW}:=(W\odot V)^\top(W\odot V)=H_W\circ H_V,
 \qquad
 K_{UVW}:=(W\odot V\odot U)^\top(W\odot V\odot U)
      =H_U\circ H_V\circ H_W,
\]
where column \(j\) of \(W\odot V\odot U\) is
\(w_j\otimes v_j\otimes u_j\).  For each row \(j\),
\[
 \sum_{\ell\ne j}|(K_{VW})_{j\ell}|
 \le \left(\max_{\ell\ne j}|(H_W)_{j\ell}|\right)
   \sum_{\ell\ne j}|(H_V)_{j\ell}|
 \le q_*^2.
\]
Thus symmetry and Gershgorin give
\((1-q_*^2)I\preceq K_{VW}\preceq(1+q_*^2)I\).  Since
\(H_U\succeq0\) and \(\operatorname{diag}(H_U)=I\), the checked Schur-order
fact in the displayed auxiliary fact gives the full sandwich
\[
 (1-q_*^2)\|\lambda\|_2^2
 \le \lambda^\top K_{UVW}\lambda
 =\left\|\sum_{j=1}^r\lambda_j
    u_j\otimes v_j\otimes w_j\right\|_F^2
 =\|T\|_F^2
 \le(1+q_*^2)\|\lambda\|_2^2.
\tag{7}
\]
This is the required full three-mode Gram identity; the two-mode Gram is used
only as the Schur-order input.
Combining (5)--(7), direct substitution of \(q_*=1/4096\) gives
\[
 { \|T-\widehat T(e)\|_F\over\|T\|_F}
 \le { (1+q_*)^{3/2}\over\sqrt{1-q_*^2}}\,4.1x
 <5x.
\tag{8}
\]
Finally Assumption~\ref{assump:base-scale} gives \(\kappa_0\ge1\), hence
\(5x\le8\kappa_0^2x\).  At \(x=0\), (2) is exactly the target tensor and
the residual is zero.  This is a conditional exact-entry statement; it does
not assert that the frozen landing produces that entry in the baseline
specialization.



\end{proof}
\end{lemma}
\begin{proposition}[covered-path stopping before the cap]\label{prop:step-012-stop-cap}
Under Assumptions~\ref{assump:base-scale} and
\ref{assump:accuracy-confidence}, Proposition~\ref{prop:step-011-chronological-contraction}, and
Lemma~\ref{lem:step-012-residual-transfer}, suppose a covered run has the
generated \(E_{\rm basin}\) entry
\(x_0=d_Q(e^{\rm land},0)\le4096\omega<1\), and let \(x_t\) be its quotient
error after \(t\) completed cyclic sweeps. With \(\ell=17/512\),
\[
 { \|T-\widehat T(e^t)\|_F\over\|T\|_F}
 \le8\kappa_0^2\ell^t.
\tag{9}
\]
Choose any universal \(C_{\rm stop}\) with
\(C_{\rm stop}\log(512/17)\ge1\) (for example \(C_{\rm stop}=1\)), and set
\[
 m_{\rm cap}=\left\lceil C_{\rm stop}\log{8\kappa_0^2\over\epsilon}\right\rceil.
\tag{10}
\]
Then the covered run has an original relative-Frobenius residual test at most
\(\epsilon\) no later than this counter.

\begin{proof}
Proposition~\ref{prop:step-011-chronological-contraction} gives
\(x_t\le\ell^t x_0\). Since \(x_0<1\),
Lemma~\ref{lem:step-012-residual-transfer} gives (9). Put
\(A=8\kappa_0^2/\epsilon\). Assumptions~\ref{assump:base-scale} and
\ref{assump:accuracy-confidence} give \(A>1\). Since \(\ell=17/512\),
(10) and the displayed choice of \(C_{\rm stop}\) imply
\[
 \ell^{m_{\rm cap}}
 \le\exp\{-C_{\rm stop}\log A\log(512/17)\}\le A^{-1}.
\tag{11}
\]
The residual at the \(m_{\rm cap}\)-th completed sweep is therefore at most
\(8\kappa_0^2A^{-1}=\epsilon\). The setting tests the original residual after
each completed sweep, so the first hit is no later than the cap. If an exact
quotient entry \(x_0=0\) is supplied, the contraction forces
\(x_t=0\) for every completed sweep, and
Lemma~\ref{lem:step-012-residual-transfer} makes the first tested residual
zero. This conditional statement does not prove the baseline landing trace.



\end{proof}
\end{proposition}
\begin{proposition}[observable cap and Moore--Penrose branch separation]\label{prop:step-012-finite-tapes}
Under Assumption~\ref{assump:accuracy-confidence}, fix an instance in the
event supplied by Proposition~\ref{prop:step-001-geometry}; equation (7)
gives \(\|T\|_F>0\). Define the
algorithmically admissible raw domain of one complete run and of all \(J\)
restarts by
\[
 \Omega_{\rm raw}^{(1),\circ}:=
 \left\{(\xi_i^{(M)})_{i\in[k],\,M\in\{U,V,W\}}:
   \|\xi_i^{(M)}\|_2>0\ \text{for every }i,M\right\},
 \qquad
 \Omega_{\rm raw}^{(J),\circ}:=
   (\Omega_{\rm raw}^{(1),\circ})^J.
\tag{12}
\]
For every tape in \(\Omega_{\rm raw}^{(1),\circ}\), the setting's procedure
is initialized and has a finite observable outcome. A zero proposal
contraction, empty certified pool, wrong cluster count, zero \(\theta_a\), or
zero active norm at the landing commit terminates that run as unsuccessful.
Otherwise every prescribed Moore--Penrose solve is defined, including for a
singular design. The cyclic phase executes no more than \(m_{\rm cap}\)
complete sweeps and returns either the first original-residual hit or an
unsuccessful timeout. No singular or off-event tape receives the generated
\(E_{\rm basin}\) recurrence. Under the conditional Gaussian law,
\(\Pr(\Omega_{\rm raw}^{(J),\circ}\mid{\cal F}_{\rm sm})=1\).

\begin{proof}
The domain (12) is exactly where every initial normalization
\({\cal N}(\xi_i^{(M)})\) in the written procedure is defined. A nondegenerate
Gaussian vector in finite dimension is zero with probability zero; a finite
union over \(3kJ\) raw vectors proves the final probability-one assertion.
No statement is made about a raw-zero array outside (12).

There are \(k\) proposal slots and each has at most
\(L_{\rm prop}+1\) old-state evaluations. A zero contraction is checked before
normalization, so that gate has a finite failure branch. Certification,
score filtering, graph construction, and connected components are finite
operations on the k-slot list. The setting declares the other listed
failures observable before a cyclic update.

For a surviving landing call, each frozen design is a finite matrix. The
Moore--Penrose rule returns a matrix even when its Gram is singular; no
ordinary inverse is silently substituted. The positive-norm landing gate
either accepts all active columns or marks the run unsuccessful. After an
accepted landing, the literal cyclic update also uses a Moore--Penrose
solution at every mode. A singular cyclic design is thus a legal finite
algebraic branch, but it is outside the hypotheses of the quotient
$\operatorname{Refresh}_s$ argument; no $d_Q$, target-span, or contraction
claim is made there.
If a proof-only canonical refresh would require division by a zero output
norm, that refresh is not invoked and the literal tape remains on the
observable capped branch. The original residual is still evaluable for the
literal factors because \(\|T\|_F>0\).

The counter is initialized before the first cyclic solve and ranges over the
finite set \(\{1,\ldots,m_{\rm cap}\}\). After each counted completed sweep,
the original residual is tested; a hit returns success, and the final nonhit
returns unsuccessful. Thus any admissible noncovered, singular, off-chart,
or nonhitting tape cannot cause unbounded work or justify the good-path
recurrence. This deterministic completion property is the event
\(E_{\rm run}^{\rm cap}\).



\end{proof}
\end{proposition}
\begin{proposition}[dense per-run runtime]\label{prop:step-012-dense-cost}
For the setting-defined finite dimensions and horizons, under the completion
conclusion of Proposition~\ref{prop:step-012-finite-tapes}, put in the dense
exact-arithmetic model
\[
 B_s:=n^3s+n^2s^2+s^3,\qquad
 L:=L_{\rm prop},\qquad m:=m_{\rm cap}.
\tag{13}
\]
There is a universal $c_{\rm cost}$ such that one complete run, including all
unsuccessful branches, costs at most
\[
 W_{\rm run}\le c_{\rm cost}\left[
 n^3+kL(n^3+n)+kn^3+k^2n+(m+1)B_k+kn\right].
\tag{14}
\]
This bound is independent of every generated condition number.

\begin{proof}
Materializing the three dense unfoldings costs
$O(n^3)$ once. A contraction of an $n\times n^2$ unfolding with one
Khatri--Rao vector costs $O(n^3)$; three contractions per Jacobi commit,
including certification look-ahead, give $O(kLn^3)$. Generating and
normalizing raw vectors costs $O(kn)$. Scores use at most $O(kn^3)$ dense
tensor contractions.
The three modewise inner products for every retained pair and the connected
graph cost $O(k^2n)$, covering filtering, clustering, and tie handling.

The frozen landing pass has three products of an n by n^2 unfolding with an
n^2 by k design, dense Gram formation, and three k-column Moore--Penrose
factorizations. Their
combined cost is $O(B_k)$; the one rebalance is lower order. A cyclic sweep
has the same dense upper bound $O(B_k)$, including construction of the
reconstructed tensor for the original residual test, which costs $O(n^3k)$.
Inactive $k-r$ columns may be removed for a sharper $B_r$ bound, but retaining
the prescribed zero-padded arrays is covered by $B_k$. Singular pseudoinverses use the same
dimension-controlled factorization, so no inverse condition number enters.
There are at most \(m\) sweeps by
Proposition~\ref{prop:step-012-finite-tapes}, proving (14).



\end{proof}
\end{proposition}
\begin{proposition}[conditional restart amplification]\label{prop:step-012-restart}
Under Assumptions~\ref{assump:random-initialization} and
\ref{assump:accuracy-confidence}, condition on a fixed instance in the event
supplied by Proposition~\ref{prop:step-001-geometry}. Suppose complete runs
use the cap in (10) and the conclusions of
Propositions~\ref{prop:step-003-coupon},
\ref{prop:step-011-chronological-contraction},
\ref{prop:step-012-stop-cap}, and
\ref{prop:step-012-finite-tapes}. Let
\[
 J=\max\{1,\lceil C_{\rm rep}\log(1/\delta_{\rm init})\rceil\},
 \qquad C_{\rm rep}\ge1/\log27.
\tag{15}
\]
Then
\[
 \Pr(\text{all \(J\) runs fail}\mid{\cal F}_{\rm sm})
 \le(1-p_0)^J=27^{-J}\le\delta_{\rm init}.
\tag{16}
\]
On the complementary event, the smallest-residual successful output has at
most \(k\) nonzero terms and original relative residual at most \(\epsilon\).

\begin{proof}
Proposition~\ref{prop:step-003-coupon} gives a one-run coverage probability
at least \(p_0=26/27\). On a fixed \(E_{\rm sm}\) tensor, the deterministic chain through
Proposition~\ref{prop:step-011-chronological-contraction} supplies the
generated \(E_{\rm basin}\) on that covered branch, and
Proposition~\ref{prop:step-012-stop-cap} supplies a residual hit before the
cap. Proposition~\ref{prop:step-012-finite-tapes} makes every admissible tape
a completed success/failure trial, and its domain has conditional probability
one. Fresh complete-run tapes make these outcomes independent conditional on
\({\cal F}_{\rm sm}\), even though the tensor is reused. Therefore the
failure product in (16) holds. Since
\[
 J\log27\ge C_{\rm rep}\log(1/\delta_{\rm init})\log27
 \ge\log(1/\delta_{\rm init}),
\]
the final inequality in (16) follows. The clamp in (15) remains exact for
\(\delta_{\rm init}\) near one.

Every successful run passed the original residual test, so minimum-residual
selection cannot increase its residual. Only the active rank-\(r\) columns
can be nonzero and \(r\le k\); all inactive columns remain exactly zero.


\end{proof}
\end{proposition}
\begin{proof}
Fix a generated \(E_{\rm sm}\) instance.  For completed run
\(a\in[J]\), let \(E_{\rm cov}^{(a)}\) and \(E_{\rm basin}^{(a)}\) denote
the generated interfaces on its fresh tape, and define
\[
 E_{\rm stop}^{(a)}:=E_{\rm cov}^{(a)}\cap E_{\rm basin}^{(a)}\cap
 \left\{\exists t\le m_{\rm cap}:
 {\|T-\widehat T(e^{t,a})\|_F\over\|T\|_F}\le\epsilon\right\},
 \qquad
 E_{\rm restart}:=\bigcup_{a=1}^J E_{\rm stop}^{(a)}.
\tag{17}
\]
Let \(E_{\rm run}^{\rm cap}\) be the event that every one of the \(J\) runs
returns either success or an observable gate/timeout failure within its
declared work.  Proposition~\ref{prop:step-012-finite-tapes} proves this for
every tape in \(\Omega_{\rm raw}^{(J),\circ}\), a conditional
probability-one domain.  Its singular and off-event branches are never fed
into the quotient recurrence.

Proposition~\ref{prop:step-003-coupon} gives
\(\Pr(E_{\rm cov}^{(a)}\mid{\cal F}_{\rm sm})\ge26/27\).  On each covered
branch, Proposition~\ref{prop:step-011-chronological-contraction} produces
the exact-span chronology and contraction (A2).
Lemma~\ref{lem:step-012-residual-transfer}
uses the full three-mode Gram sandwich (7) to transfer that quotient error to
the original relative Frobenius residual, and
Proposition~\ref{prop:step-012-stop-cap} proves
\(E_{\rm cov}^{(a)}\subseteq E_{\rm stop}^{(a)}\).  Therefore
Proposition~\ref{prop:step-012-restart} gives the nested probabilities
\[
 \Pr(E_{\rm restart}\mid{\cal F}_{\rm sm})\ge1-\delta_{\rm init}
 \quad\text{on }E_{\rm sm},
 \qquad
 \Pr(E_{\rm sm}\cap E_{\rm restart})
 \ge(1-\delta_{\rm sm})(1-\delta_{\rm init}).
\tag{18}
\]
The second inequality is the tower calculation
\(\mathbb E[1_{E_{\rm sm}}
\Pr(E_{\rm restart}\mid{\cal F}_{\rm sm})]\); it does not move either
confidence parameter into \(k\).  Every successful output passed the original
residual test and has only its \(r\le k\) active columns nonzero.

It remains to charge all restarts rather than only one run.  Define the
proof-local ceiling bounds
\[
 \bar k:=1+C_{\rm rank}r^{5/3}(\log r)^{5/2},\qquad
 R_\epsilon:=2+C_{\rm stop}
 \left(\log8+2\log\kappa_0+\log(1/\epsilon)\right).
\]
From the setting definitions and \(0<\delta_{\rm init}<1\),
\[
 k\le\bar k,\quad
 L_{\rm prop}\le2+(C_{\rm burn}+C_{\rm cert})\log r,\quad
 m_{\rm cap}+1\le R_\epsilon,\quad
 J\le1+C_{\rm rep}\log(1/\delta_{\rm init}).
\tag{19}
\]
Applying Proposition~\ref{prop:step-012-dense-cost} to each completed run
and summing, including early failed runs, gives the displayed total-work
bound
\[
\begin{aligned}
 W_{\rm total}\ &\le J W_{\rm run}\\
 &\le c_{\rm cost}\bigl(1+C_{\rm rep}\log(1/\delta_{\rm init})\bigr)
 \Big[ n^3
 +\bar k\{2+(C_{\rm burn}+C_{\rm cert})\log r\}(n^3+n)
 +\bar k n^3+\bar k^2n\\
 &\hspace{46mm}
 +R_\epsilon(n^3\bar k+n^2\bar k^2+\bar k^3)+\bar k n\Big].
\end{aligned}
\tag{20}
\]
In particular, after exposing the prescribed rank powers, (20) is
\[
\begin{aligned}
 W_{\rm total}=O\!\Big((1+\log(1/\delta_{\rm init}))\big[&n^3
 +n^3r^{5/3}(\log r)^{7/2}+nr^{10/3}(\log r)^5\\
 &+(1+\log\kappa_0+\log(1/\epsilon))
 \{n^3r^{5/3}(\log r)^{5/2}
 +n^2r^{10/3}(\log r)^5+r^5(\log r)^{15/2}\}\big]\Big).
\end{aligned}
\tag{21}
\]
The hidden constant in (21) depends only on the fixed universal algorithmic
constants.  The bound is independent of \(\rho^{-1}\), hence is in
particular polynomial in the full list of variables allowed by the setting.

For the exact/noiseless baseline, the scope of this step is conditional: if
the exact quotient entry \(e^0=0\) is supplied,
Proposition~\ref{prop:step-011-chronological-contraction} keeps every completed
cyclic state at zero and Lemma~\ref{lem:step-012-residual-transfer} gives an
exactly zero first tested residual for every \(\epsilon>0\).  Neither the cap
nor the restart rule weakens that conclusion.  This step does not prove that
the frozen landing produces \(e^0=0\); the complete landing specialization is
proved separately in Proposition~\ref{prop:step-013-frozen-landing}.
\end{proof}

\begin{proposition}[Rate Specialization Bridge]\label{prop:rate-specialization-bridge}
Under Assumptions~\ref{assump:base-scale}--\ref{assump:accuracy-confidence},
let $E_{\rm sm}$ be the event supplied by
Proposition~\ref{prop:step-001-geometry}, and condition on a fixed instance in
that event. Choose
\[
 k=\left\lceil C_{\rm rank}r^{5/3}(\log r)^{5/2}\right\rceil,
 \quad m_{\rm cap}=\left\lceil C_{\rm stop}\log\frac{8\kappa _0^2}{\epsilon}\right\rceil,
 \quad J=\max\left\{1,\left\lceil C_{\rm rep}\log\frac1{\delta_{\rm init}}\right\rceil\right\}.
\tag{B.1}
\]
Let
\[
 K_r:=r^{5/3}(\log r)^{5/2},\quad
 \bar k:=1+C_{\rm rank}K_r,\quad
 L_r:=2+(C_{\rm burn}+C_{\rm cert})\log r,\quad
 R_\epsilon:=2+C_{\rm stop}\bigl(\log 8+2\log\kappa _0+\log(1/\epsilon)\bigr).
\tag{B.2}
\]
Then
\[
 k\le\bar k,\quad
 L_{\rm prop}\le L_r,\quad
 m_{\rm cap}+1\le R_\epsilon,\quad
 J\le1+C_{\rm rep}\log(1/\delta_{\rm init}).
\tag{B.3}
\]
If $p_0\ge26/27$ is the one-run coverage probability from
Proposition~\ref{prop:step-003-coupon}, then
\[
 (1-p_0)^J\le27^{-J}\le\delta_{\rm init}
 \quad\text{whenever }C_{\rm rep}\ge1/\log 27.
\tag{B.4}
\]
Thus
\[
\Pr(E_{\rm restart}\mid\mathcal F_{\rm sm})\ge1-\delta_{\rm init}
\quad\text{on }E_{\rm sm},\qquad
\Pr(E_{\rm sm}\cap E_{\rm restart})
\ge(1-\delta_{\rm sm})(1-\delta_{\rm init}).
\tag{B.5}
\]
Moreover, the dense per-run bound of
Proposition~\ref{prop:step-012-dense-cost} gives
\[
\begin{aligned}
W_{\rm total}
\le c_{\rm cost}\bigl(1+C_{\rm rep}\log(1/\delta_{\rm init})\bigr)
\Big[&n^3+\bar kL_r(n^3+n)\\
&+\bar k n^3+\bar k^2n
+R_\epsilon(n^3\bar k+n^2\bar k^2+\bar k^3)+\bar k n\Big].
\end{aligned}
\tag{B.6}
\]
For $r\ge3$ and $k\le n$, set
$C_K:=1+C_{\rm rank}$,
$C_L:=2+C_{\rm burn}+C_{\rm cert}$, and
$C_R:=2+C_{\rm stop}(\log8+2)$. Then
\[
K_r\ge1,\qquad \bar k\le C_KK_r,\qquad
L_r\le C_L\log r,\qquad
R_\epsilon\le C_R\bigl(1+\log\kappa _0+\log(1/\epsilon)\bigr),
\tag{B.7}
\]
and every term suppressed from (B.6) is dominated through
\[
\begin{aligned}
&\bar kL_r(n^3+n)+\bar k n^3+\bar k n
 \le4C_K(C_L+1)n^3K_r\log r,\\
&\bar k^2n\le C_K^2nK_r^2,\\
&R_\epsilon(n^3\bar k+n^2\bar k^2+\bar k^3)\\
&\hspace{15mm}\le C_RC_K^3
 (1+\log\kappa _0+\log(1/\epsilon))
 (n^3K_r+n^2K_r^2+K_r^3).
\end{aligned}
\tag{B.8}
\]
Consequently,
\[
\begin{aligned}
W_{\rm total}=O\!\Big((1+\log(1/\delta_{\rm init}))\big[&n^3
+n^3r^{5/3}(\log r)^{7/2}+nr^{10/3}(\log r)^5\\
&+(1+\log\kappa _0+\log(1/\epsilon))
\{n^3r^{5/3}(\log r)^{5/2}
+n^2r^{10/3}(\log r)^5+r^5(\log r)^{15/2}\}\big]\Big).
\end{aligned}
\tag{B.9}
\]
The hidden constant depends only on the fixed universal algorithmic constants.
Assumptions~\ref{assump:base-scale} and
\ref{assump:gaussian-smoothing} give
$\kappa _0\le r^{d_\kappa}$ and $\rho^{-1}\le r^{d_\rho}$; the cost bound is
independent of $\rho^{-1}$ and hence is polynomial in the full declared list.
\end{proposition}

\begin{proof}
The ceiling in the definition of $k$ gives the first inequality in (B.3),
and the definitions of $L_{\rm burn}$ and $L_{\rm cert}$ give the second. Since
$0<\epsilon<1$ and $\kappa _0\ge1$,
\[
\log\frac{8\kappa _0^2}{\epsilon}
 =\log 8+2\log\kappa _0+\log(1/\epsilon),
\]
which proves the third inequality in (B.3); the last follows from the ceiling
and $0<\delta_{\rm init}<1$.  The conditional one-run bound is
Proposition~\ref{prop:step-003-coupon}; independence of complete runs from
Assumption~\ref{assump:random-initialization} gives the first inequality in
(B.4), and the choice of $C_{\rm rep}$ gives the second. The deterministic
covered-run chain through Propositions~\ref{prop:step-011-chronological-contraction}
and~\ref{prop:step-012-stop-cap} converts coverage into success, so (B.4)
gives the conditional statement in (B.5). Proposition~\ref{prop:step-001-geometry}
gives $\Pr(E_{\rm sm})\ge1-\delta_{\rm sm}$; the tower property gives the
joint bound in (B.5).

Substituting (B.3) into the per-run expression of
Proposition~\ref{prop:step-012-dense-cost} and multiplying by the bound on
$J$ gives (B.6). Since $\log r\ge\log3>1$, (B.7) follows directly. Also
$n^3+n\le2n^3$ and $n\le n^3$, which proves the first line of (B.8); the
other two lines follow by substituting $\bar k\le C_KK_r$. Finally,
$K_r^2=r^{10/3}(\log r)^5$ and
$K_r^3=r^5(\log r)^{15/2}$, so (B.8) gives (B.9) with no omitted term.
The two envelope inequalities are primitive parts of the cited assumptions.
No confidence parameter is absorbed into $k$, and the probability conversion
remains conditional until the displayed tower calculation.
\end{proof}

\subsection{Exact orthogonal baseline}\label{app:step-013}

Assume the deterministic scope condition $\mathsf{ExactCertifiedSeed}$ from
the preliminaries.  The proof below checks the frozen-input solves, the
product-preserving rebalance, all chronological mode orderings, and the
original residual directly.  We use only the elementary identities
$(A\odot B)^\top(A\odot B)=(A^\top A)\circ(B^\top B)$ and the
Moore--Penrose formula for a positive scalar multiple of an orthonormal
matrix; their three mode-wise instances are displayed explicitly.

Set
\[
G_U:=U,\quad G_V:=V,\quad G_W:=W,\qquad
K_U:=W\odot V,\quad K_V:=W\odot U,\quad K_W:=V\odot U.
\tag{1}
\]
The Khatri--Rao column \(K_M(:,j)\) is the tensor product of the two held
target directions in the order used by the corresponding unfolding. Under
\(\mathsf{ExactCertifiedSeed}\),
\[
K_U^\top K_U=K_V^\top K_V=K_W^\top K_W=I_r,
\tag{2}
\]
because each factor Gram is \(I_r\). The target tensor and its unfoldings are
\[
T=\lambda\sum_{j=1}^r u_j\otimes v_j\otimes w_j,\quad
T_{(1)}=\lambda U K_U^\top,\quad
T_{(2)}=\lambda V K_V^\top,\quad
T_{(3)}=\lambda W K_W^\top.
\tag{3}
\]


\begin{lemma}[exact orthogonal equal-weight seed and frozen Grams]\label{lem:step-013-baseline-geometry}
Under the local
specialization \(\mathsf{ExactCertifiedSeed}\) and the seed convention from Proposition~\ref{prop:step-008-seed-interface}, the
observable scores satisfy \(\theta_j=\lambda\), the balanced frozen seed is
\[
X^0=\lambda^{1/3}U,\qquad Y^0=\lambda^{1/3}V,\qquad
Z^0=\lambda^{1/3}W,
\tag{4}
\]
and, for every \(M\in\{U,V,W\}\),
\[
P_M^0=G_M,\quad D_M^0=\lambda^{1/3}I_r,\quad
C_M^0=0,\quad N_M^0=0,\quad s^0=0,
\tag{5}
\]
\[
F_M^0=(P_M^0)^\top P_M^0=I_r,\quad
G_M^\top P_M^0=I_r.
\tag{6}
\]
Consequently, in all three held-mode orderings,
\[
\begin{array}{c|c|c}
M & J_M^0=F_{M'}^0\circ F_{M''}^0 &
Q_M^0=(G_{M'}^\top P_{M'}^0)\circ(G_{M''}^\top P_{M''}^0)\\ \hline
U & I\circ I=I & I\circ I=I\\
V & I\circ I=I & I\circ I=I\\
W & I\circ I=I & I\circ I=I
\end{array}
\tag{7}
\]
and every actual frozen pair Gram is \(\lambda^{4/3}I_r\).

\begin{proof}
For an exact representative,
\[
\theta_j=\left\langle
\lambda\sum_{\ell=1}^r u_\ell\otimes v_\ell\otimes w_\ell,\,
u_j\otimes v_j\otimes w_j\right\rangle
=\lambda
\tag{8}
\]
by the orthonormality identities. If a product-one sign chart is used, the three signs multiply to
one and the same inner product and represented rank-one term result. The
setting's observable seed therefore gives (4), and normalizing its columns
gives (5). Equations (6)--(7) follow from the orthonormality identities and the Khatri--Rao identity.
Finally, for each mode \(M\),
\[
H_M^0=(P_{M'}^0D_{M'}^0)\odot(P_{M''}^0D_{M''}^0)
=\lambda^{2/3}K_M,
\tag{9}
\]
so (2) gives \((H_M^0)^\top H_M^0=\lambda^{4/3}I_r\). This also proves
that the exact seed has no perpendicular or product-log defect. The line
\(\lambda^{4/3}I_r\) is valid for each of the U, V, and W orderings in (7).
Thus the actual left-hand defects in the bridge (the differences bounded by
\(\delta_L,\delta_F,\delta_{FL}\), and \(Q_M^0-J_M^0\)) are zero, while the
conservative budget constants \(\eta_J,\eta_{QJ},\eta_A\) may remain positive
at the main-theorem threshold values. All these bounds are therefore sharp
upper bounds at the baseline, and the budget constants become zero in the
optional formal \(q_*,\tau_r\downarrow0\) no-floor limit; no budget is used
to weaken the equalities above.



\end{proof}
\end{lemma}
\begin{proposition}[three exact frozen landing solves]\label{prop:step-013-frozen-landing}
Under \(\mathsf{ExactCertifiedSeed}\), with
the synchronized frozen designs of the setting, the Moore--Penrose outputs
are well-defined and, mode by mode,
\[
\begin{array}{c|c|c|c}
M & H_M^0 & T_{(M)}H_M^0 &
T_{(M)}H_M^0((H_M^0)^\top H_M^0)^\dagger\\ \hline
U & \lambda^{2/3}K_U & \lambda^{5/3}U & \lambda^{1/3}U\\
V & \lambda^{2/3}K_V & \lambda^{5/3}V & \lambda^{1/3}V\\
W & \lambda^{2/3}K_W & \lambda^{5/3}W & \lambda^{1/3}W
\end{array}
\tag{10}
\]
Thus \(\widetilde X_U=\lambda^{1/3}U\),
\(\widetilde X_V=\lambda^{1/3}V\), and
\(\widetilde X_W=\lambda^{1/3}W\), with every active output norm equal to
\(\lambda^{1/3}>0\).

\begin{proof}
The U row follows from (2)--(3) and (9):
\[
T_{(1)}H_U^0
=\lambda U K_U^\top(\lambda^{2/3}K_U)
=\lambda^{5/3}U,\qquad
((H_U^0)^\top H_U^0)^\dagger=\lambda^{-4/3}I_r.
\tag{11}
\]
Their product is \(\lambda^{1/3}U\). Replacing \((U,K_U)\) by
\((V,K_V)\) and \((W,K_W)\) gives, respectively,
\[
T_{(2)}H_V^0=\lambda^{5/3}V,\qquad
T_{(3)}H_W^0=\lambda^{5/3}W,
\tag{12}
\]
with the same Gram inverse and the same output scale. These are three
independent evaluations from the one frozen seed; no output in one row is an
input to another row. Since \(\lambda^{4/3}I_r\) is nonsingular, the
Moore--Penrose rule agrees with the displayed ordinary inverse, while the
calculation remains valid at the exact \(\rho_{\rm sm}=0\) specialization.



\end{proof}
\end{proposition}
\begin{proposition}[joint rebalance and product-one gauge]\label{prop:step-013-rebalance-gauge}
Under \(\mathsf{ExactCertifiedSeed}\) and the
outputs of Proposition~\ref{prop:step-013-frozen-landing}, the setting's
single joint rebalance returns
\[
X_U^{\rm land}=\lambda^{1/3}U,\qquad
X_V^{\rm land}=\lambda^{1/3}V,\qquad
X_W^{\rm land}=\lambda^{1/3}W,
\tag{13}
\]
preserves every represented rank-one product, and has
\(\gamma_j^{\rm land}=\lambda\) and \(s_j^{\rm land}=0\) for every \(j\).
The same represented tensor and quotient state result after any positive
product-one scale gauge or proof-only product-one sign gauge.

\begin{proof}
Proposition~\ref{prop:step-013-frozen-landing}
gives \(\|\widetilde x_j\|_2=\|\widetilde y_j\|_2
=\|\widetilde z_j\|_2=\lambda^{1/3}\). Hence
\[
\widetilde\gamma_j
=\|\widetilde x_j\|_2\|\widetilde y_j\|_2\|\widetilde z_j\|_2
=\lambda,
\tag{14}
\]
and the rebalance multiplier is \(\widetilde\gamma_j^{1/3}
=\lambda^{1/3}\). Dividing each output by its norm therefore gives (13).
More generally, for nonzero factors and positive scalars
\(R_{U,j},R_{V,j},R_{W,j}\) with
\(R_{U,j}R_{V,j}R_{W,j}=1\),
\[
(R_{U,j}x_j)\otimes(R_{V,j}y_j)\otimes(R_{W,j}z_j)
=x_j\otimes y_j\otimes z_j.
\tag{15}
\]
The rebalance uses positive norms and is exactly such a product-one
rescaling. A product-one sign gauge has the same identity with signs in
place of positive scalars. Thus neither gauge changes the represented term,
and the common product remains \(\lambda\), giving
\(s_j=\log(\lambda/\lambda)=0\).



\end{proof}
\end{proposition}
\begin{proposition}[all cyclic normal equations are stationary]\label{prop:step-013-cyclic-fixed}
Under \(\mathsf{ExactCertifiedSeed}\), start a
cyclic active ALS sweep from the balanced state (13), with inactive columns
zero. For each mode \(M\in\{U,V,W\}\), regardless of which mode is first,
the held design is \(\lambda^{2/3}K_M\), its Gram is
\(\lambda^{4/3}I_r\), and the exact normal equation returns
\(\lambda^{1/3}G_M\). After the product-preserving canonical refresh, the
state is unchanged and the refreshed common-product register is \(s=0\).
Consequently every chronological Refresh$_s$ state and every full cyclic
sweep equal the same exact quotient state; this holds for all six
permutations of the U/V/W mode order.

\begin{proof}
For the displayed U/V/W order, the U block uses
the held factors \(Y=\lambda^{1/3}V\), \(Z=\lambda^{1/3}W\), hence
\[
H_U=\lambda^{2/3}K_U,\quad H_U^\top H_U=\lambda^{4/3}I_r,\quad
T_{(1)}H_U=\lambda^{5/3}U,\quad
\widetilde X_U=\lambda^{1/3}U.
\tag{16}
\]
Its solved product with the held factors is \(\lambda\) in each component,
so Refresh$_s$ writes \(s_U^+=0\) and the canonical representative is again
the state (13). The V block then has the unchanged held U and W factors and
similarly obeys
\[
H_V=\lambda^{2/3}K_V,\quad
T_{(2)}H_V=\lambda^{5/3}V,\quad
\widetilde X_V=\lambda^{1/3}V,\quad s_V^+=0.
\tag{17}
\]
The W block gives
\[
H_W=\lambda^{2/3}K_W,\quad
T_{(3)}H_W=\lambda^{5/3}W,\quad
\widetilde X_W=\lambda^{1/3}W,\quad s_W^+=0.
\tag{18}
\]
Thus the U/V/W chronological intermediate states are all exactly (13).
For any other first mode, use the same generic identity
\[
T_{(M)}=\lambda G_MK_M^\top,\qquad
H_M=\lambda^{2/3}K_M,\qquad K_M^\top K_M=I_r,
\tag{19}
\]
and then apply the remaining two identities in their chosen order. This
proves all six mode permutations, not only the displayed cyclic order.

For completeness, the scale-equivariance interface can be checked
directly. If a literal state is changed by positive diagonal
\(R_UR_VR_W=I\), then the mode-M design changes to \(H_MR_M^{-1}\), and
\[
T_{(M)}H_MR_M^{-1}
\big(R_M^{-1}H_M^\top H_MR_M^{-1}\big)^{-1}
=T_{(M)}H_M(H_M^\top H_M)^{-1}R_M.
\tag{20}
\]
Therefore the solved factors differ only by the same gauge, and the
canonical Refresh$_s$ representative is (13). The sign gauge is handled by
the proof-only orientation and preserves each product by (15). Induction
over sweeps proves stationarity forever, with no accumulated forcing term.



\end{proof}
\end{proposition}
\begin{proposition}[exact residual and no-floor stopping]\label{prop:step-013-zero-residual}
Under \(\mathsf{ExactCertifiedSeed}\) and
Assumption~\ref{assump:accuracy-confidence}, after the joint landing and
after every cyclic refresh, the active factors represent exactly
\[
\widehat T
=\lambda\sum_{j=1}^r u_j\otimes v_j\otimes w_j=T,
\tag{21}
\]
and the inactive \(k-r\) columns contribute zero. Hence
\[
\frac{\|T-\widehat T\|_F}{\|T\|_F}=0\le\epsilon,\qquad
\|T\|_F=\lambda\sqrt r>0.
\tag{22}
\]
The first original residual test is therefore successful (whether it is
performed immediately after landing or after the first completed cyclic
sweep), for every \(\epsilon>0\); no \(\omega\)-, \(\tau_r\)-, or positive-radius
error floor is present.

\begin{proof}
Equation (13) gives, componentwise,
\[
x_j^{\rm land}\otimes y_j^{\rm land}\otimes z_j^{\rm land}
=\lambda u_j\otimes v_j\otimes w_j.
\tag{23}
\]
Proposition~\ref{prop:step-013-cyclic-fixed} shows that the same equality is
preserved by every block and every refresh. The tensors
\(u_j\otimes v_j\otimes w_j\) are mutually orthonormal in the Frobenius inner
product under the orthonormality identities, so \(\|T\|_F^2=r\lambda^2\). This proves (21)--(22) and
the stopping assertion. In the deterministic \(\rho_{\rm sm}=0\) limit
there is no perturbation term to dominate, and in the optional formal
no-floor limit \(q_*,\tau_r\downarrow0\) the same exact calculation remains
valid. The positive \(\rho_{\rm ALS}\) chart radius is only a neighborhood
for the main theorem; the baseline state is its exact origin, so no radius
term enters (22).


\end{proof}
\end{proposition}
\begin{proof}
Lemma~\ref{lem:step-013-baseline-geometry} derives the exact observable score,
balanced seed, zero coefficient/perpendicular/product-log fields, and all
three identity pair/cross Grams. Proposition~\ref{prop:step-013-frozen-landing}
then evaluates the three synchronized frozen-input pseudoinverse calls from
that single seed. It explicitly gives the U, V, and W outputs and their
positive normalizers, so no mode-specific predecessor or post-solve design is
used. Proposition~\ref{prop:step-013-rebalance-gauge} applies the one joint
positive rebalance, proves product preservation, and handles both positive
scale and product-one sign gauges. Proposition~\ref{prop:step-013-cyclic-fixed}
checks the chronological U/V/W normal equations and their generic permutation,
and proves that every Refresh$_s$ register remains zero. Finally,
Proposition~\ref{prop:step-013-zero-residual} identifies the represented tensor
with \(T\), computes its strictly positive Frobenius norm, and applies the
residual-test conclusion of Proposition~\ref{prop:step-012-stop-cap}. Thus all three frozen landing
outputs, every cyclic refresh, the common product error, and the original
residual have the exact values required by the stated result. The argument is a
deterministic specialization and leaves the positive-\(\rho\) main-theorem
assumptions untouched.
\end{proof}

\subsection{Proof of the Main Theorem}\label{app:main-proof}
\begin{proof}[Proof of the main theorem]
Proposition~\ref{prop:step-001-geometry} constructs $E_{\rm sm}$, proves
$\Pr(E_{\rm sm})\ge1-\delta_{\rm sm}$, and supplies the realized norm, Gram,
weight-ratio, and Khatri--Rao conclusions. Conditional on a fixed instance in
that event, Proposition~\ref{prop:step-002-joint-window} gives the two-sided
window mass and Proposition~\ref{prop:step-003-coupon} gives coverage
probability at least $26/27$.

Lemma~\ref{lem:step-004-recurrence} proves the displayed $R,S$ recurrence;
Proposition~\ref{prop:step-004-envelope} proves denominator persistence and
finite chart entry, and Proposition~\ref{prop:step-004-certificate} gives the
stored certificate. Proposition~\ref{prop:step-005-threshold-ledger} supplies
the score and coordinate ledger. Proposition~\ref{prop:step-006-graph} applies
the public observable graph rule, while Lemma~\ref{lem:step-006-gauge} gives
the target permutation and product-one gauge. Propositions~
\ref{prop:step-007-column} and~\ref{prop:step-007-row} close both coefficient
orientations, and Proposition~\ref{prop:step-008-seed-interface} gives the
balanced-seed bounds.

Proposition~\ref{prop:complete-frozen-reserve-interface} collects the complete
pre-solve Gram, inverse, congruence, coefficient, and normalizer interface.
It is consumed by Propositions~\ref{prop:step-010-frozen-solves},
\ref{prop:step-010-directions}, and~\ref{prop:step-010-rebalance}, whose
composition in Proposition~\ref{prop:step-010-quotient-entry} produces the
exact-span landing state without a cross-mode landing dependency.
Proposition~\ref{prop:step-011-chronological-contraction} then gives the
literal cyclic orbit identification, chart invariance, and contraction.

Lemma~\ref{lem:step-012-residual-transfer} converts quotient error to the
original Frobenius residual. Propositions~\ref{prop:step-012-stop-cap} and
\ref{prop:step-012-finite-tapes} prove covered-path stopping and finite work on
every tape. Proposition~\ref{prop:step-012-dense-cost} supplies the dense
per-run cost, Proposition~\ref{prop:step-012-restart} gives conditional
restart confidence, and Proposition~\ref{prop:rate-specialization-bridge}
proves the displayed total-work rate and the nested probability conversion.
Combining these exact producer statements proves every conclusion of
Theorem~\ref{thm:main} without an additional assumption or rate loss.
\end{proof}
